\section{Geometric steps II: assembly of the steps}\label{sec:pf:steps2}

\providecommand{\stepsiiH}{H^{\bullet}}
\providecommand{\stepsiiPiece}{\mathrm{Piece}}
\providecommand{\stepsiiOw}{\mathcal{O}_{(w)}}
\providecommand{\stepsiiNw}{\mathrm{N}w}
\providecommand{\stepsiiHk}{P^{\mathrm{Hk}}}
\providecommand{\stepsiirt}{P^{\mathrm{rt}}}
\providecommand{\stepsiiop}{\mathrm{op}}
\providecommand{\stepsiitr}{\mathrm{tr}}
\providecommand{\stepsiisp}{\mathrm{sp}}
\providecommand{\stepsiico}{\mathrm{co}}
\providecommand{\stepsiiSec}{\mathrm{Sec}}
\providecommand{\stepsiicyc}{\mathrm{cyc}}
\providecommand{\stepsiiFr}{\mathrm{Fr}}
\providecommand{\stepsiiCR}{\mathrm{CR}}
\providecommand{\stepsiiGood}{\mathrm{Good}}
\providecommand{\stepsiiCong}{\mathrm{Cong}}
\providecommand{\stepsiiES}{\mathrm{ES}}
\providecommand{\stepsiiJ}{\mathrm{J}}
\providecommand{\stepsiiPhi}{\Phi}
\providecommand{\stepsiiHdg}{\mathrm{Hdg}}
\providecommand{\stepsiiW}{\mathrm{Wt}}
\providecommand{\stepsiiRich}{\mathrm{Rich}}
\providecommand{\stepsiiPar}{\mathrm{Par}}
\providecommand{\stepsiitw}{\mathrm{tw}}
\providecommand{\stepsiim}{\mathrm{m}}

This section proves the proposition $\sfname{Steps}$ of the formal development
(\cref{def:pf:steps2:steps}), a conjunction of seventeen statements \textup{(S1)--(S17)} about the
$\ell$-adic Galois representations that are cut out, by eigenvalues of Hecke correspondences, from
the $\ell$-adic cohomology of a smooth projective variety over a number field. Its proof
(\cref{thm:pf:steps2:main}) is used in the proof of \cref{thm:geoone}, where the seventeen
statements are applied to the Hecke eigen-pieces of Kuga--Sato type towers.

Three of the statements are proved in \cref{sec:pf:steps1}: crystallinity at places above $\ell$
where the tower has a smooth projective model (S3), the occurrence of Hodge degrees among the
Hodge--Tate weights (S4), and the crystalline Eichler--Shimura relation (S5). Ten elementary or
classical ones are proved together with \cref{thm:geoone}: the twisting lemmas (S7)--(S10), two
changes of normalisation (S11), (S13), the cardinality of Satake multisets (S12), the $\ell$-adic
characters attached to algebraic Hecke characters (S15), and the semisimplicity and continuity of
summands (S16), (S17). The remaining four are proved here:
\begin{itemize}[nosep]
\item (S1), \cref{thm:pf:steps2:repex}: every Hecke eigen-piece of the cohomology of a nice tower
  carries a continuous semisimple Galois representation which is a graded part of the Galois action
  on it;
\item (S2), \cref{thm:pf:steps2:esstep}: at a place $w\nmid\ell$ where the tower has a model with a
  \emph{pointwise congruence relation}, the representation is unramified and the inverse of every
  Frobenius element is annihilated by the Hecke polynomial
  $\sum_j(-1)^jq^{\binom j2}\vartheta_jX^{n-j}$;
\item (S6), \cref{thm:pf:steps2:cut}: a model can always be cut down to one on which the covering
  map of generic fibres is surjective;
\item (S14), \cref{thm:pf:steps2:richof}: a lacunary weight at a place above $\ell$ produces a
  lacunary set of Hodge--Tate weights.
\end{itemize}
The proof of (S2) carries most of the weight. Its structure is the classical one (Eichler, Shimura,
Ihara, Deligne; compare \cite{Deligne1971,HarrisTaylor2001}): the eigen-piece descends to the
generic fibre of the model (\cref{lem:pf:steps2:fixdesc}); by the specialization isomorphism, the
inertia group acts trivially and the inverse of Frobenius acts as the relative Frobenius of the
special fibre (\cref{prop:pf:steps2:esgen}); and the congruence relation, which is assumed only as
an identity of multisets of points at the generic points of the special fibre, is promoted to an
identity of operators on cohomology. For the last point the formal development proves that the
action of a weighted finite correspondence on cohomology is determined by its cycle on a dense set
of points (\cref{prop:pf:steps2:traces}(g)), and that the multiplicities of these cycles at points
of the special fibre can be computed by counting sections over valuation rings
(\cref{lem:pf:steps2:wtval}).

Finally, \cref{sec:pf:steps2:found} records foundational results of \'etale cohomology and $p$-adic
Hodge theory which are proved in the formal development and used by several of the proofs of
\cref{sec:pf:steps1}: base change theorems in the forms used, cohomological dimension, Kummer
theory, cohomology of fields, base change of towers, and internal steps of the formal proof of the
crystalline comparison theorem. For the classical ones we give precise references (and the deduction
when the formal form differs from the published one); statements specific to the development are
proved there, except for the long internal steps of the comparison theorem and of the base change of
towers, which are listed in \cref{rem:pf:steps2:deferred}.

\paragraph{Organisation.}
\Cref{sec:pf:steps2:defs} fixes the definitions occurring in the seventeen statements and states
them. \Cref{sec:pf:steps2:assembly} proves the main theorem from the four results of this section
and the results proved elsewhere.
\Cref{sec:pf:steps2:rich,sec:pf:steps2:cut,sec:pf:steps2:repex,sec:pf:steps2:es} prove (S14), (S6),
(S1) and (S2). \Cref{sec:pf:steps2:trproof} proves the trace formalism used in the proof of (S2),
\cref{sec:pf:steps2:aux} proves auxiliary results (residue data, Galois closures, finiteness and
continuity of $\ell$-adic cohomology), and \cref{sec:pf:steps2:found} contains the foundational
statements.

% ===========================================================================
\subsection{Definitions and the step statements}\label{sec:pf:steps2:defs}

\paragraph{Conventions.}
Schemes are arbitrary schemes. For a field $k$, an extension $\Omega/k$ and a $k$-scheme
$q\colon X\to\Spec k$ we write $X_\Omega=X\times_{\Spec k}\Spec\Omega$; for $\sigma\in\Aut_k(\Omega)$
we write $\sigma_X=\mathrm{id}_X\times\Spec(\sigma)\colon X_\Omega\to X_\Omega$. For a morphism
$u\colon X\to Y$ of $k$-schemes, $u_\Omega\colon X_\Omega\to Y_\Omega$ is its base change. A
morphism $f\colon Y\to X$ is \emph{finite \'etale of degree $d$} if it is finite and \'etale and
every geometric point $\Spec K\to X$ ($K$ algebraically closed) has exactly $d$ lifts to $Y$. A
morphism $X\to\Spec R$ is \emph{projective} if it factors as a closed immersion
$X\hookrightarrow\mathbb P^N_R$ followed by the projection. Throughout, $E$ is a number field,
$\Omega$ an algebraic closure of $E$ and $G_E=\Aut_E(\Omega)$ with its Krull topology; $\ell$ is a
prime and $\Qlbar$ the algebraic closure of $\Ql$ used by the development, with its norm topology.
Places of $\Omega$ above a finite place $w$ of $E$, decomposition groups $D_A$, inertia groups
$I_A$, Weil groups $W_A$ with $\deg\colon W_A\to\Z$, and arithmetic Frobenius elements are as in
\cref{sec:galois:local}; explicitly, a valuation ring $A\subset\Omega$ lies over $w$ if for
$x\in\OO_E$ one has $v_A(x)<1$ if and only if $x\in w$, the inertia group $I_A$ consists of the
$s\in D_A$ with $v_A(s(x)-x)<1$ for $x\in A$, and $\sigma$ is an arithmetic Frobenius element for
$(A,q)$ if $\sigma\in D_A$ and $v_A(\sigma(x)-x^q)<1$ for $x\in A$. We write
$\stepsiiNw=\#(\OO_E/w)$ and $\stepsiiOw$ for the valuation ring of $E$ defined by $w$ (the
localization of $\OO_E$ at $w$, a discrete valuation ring with fraction field $E$).

\begin{definition}[\'Etale and $\ell$-adic cohomology]\label{def:pf:steps2:coh}
For a scheme $X$, an integer $q\ge1$ and $n,i\in\Z$, $H^n(X,\Z/q(i))$ denotes the degree-$n$
cohomology of $X$, viewed as a constant simplicial presheaf on the big \'etale site of all
schemes, with coefficients in the Tate twist sheaf $\Z/q(i)$ of $\Z/q$-modules (the group of
homotopy classes of maps from the free chain complex of $X$ to an injective resolution of
$\Z/q(i)$). For a morphism $f$ we write $f^\ast$ for pull-back. For a prime $\ell$,
\begin{gather*}
  H^n(X,\Zl(i))=\varprojlim_k H^n(X,\Z/\ell^k(i)),\qquad
  H^n(X,\Ql(i))=\Ql\otimes_{\Zl}H^n(X,\Zl(i)),\\
  \stepsiiH(X)=\bigoplus_{n\in\Z}H^n(X,\Ql),
\end{gather*}
where the limit is the module of compatible sequences under the reduction maps, and
$H^n(X,\Ql)=H^n(X,\Ql(0))$. If $q\colon X\to\Spec k$ and $\Omega/k$ is an extension, the group
$\Aut_k(\Omega)$ acts on $\stepsiiH(X_\Omega)$ by $\sigma\mapsto\sigma^\ast:=(\sigma_X)^\ast$. One has
$1^\ast=\mathrm{id}$ and $(\sigma\tau)^\ast=\sigma^\ast\circ\tau^\ast$, by functoriality of pull-back
and $(\sigma\tau)_X=\tau_X\circ\sigma_X$. The same notation is used for each $H^n(X_\Omega,\Ql(i))$
and $H^n(X_\Omega,\Zl(i))$.
\formal{GlobalTheory.GalV.Htot}\formal{GlobalTheory.GalV.galTot}
\end{definition}

\begin{definition}[Hecke towers]\label{def:pf:steps2:tower}
A \emph{Hecke tower} $H$ over a field $k$ consists of a set $\Gamma$, a $k$-scheme
$q\colon Z\to\Spec k$, and for each $g\in\Gamma$ a scheme $Y_g$ with morphisms
$s_g,t_g\colon Y_g\to Z$ such that $q\circ s_g=q\circ t_g$. It is \emph{nice} if $q$ is smooth and
projective and every $s_g$ is finite \'etale of some degree. For $\Omega/k$ the base change
$H_\Omega$ has $Z_\Omega$ and the correspondences $s_{g,\Omega},t_{g,\Omega}\colon Y_{g,\Omega}\to
Z_\Omega$; for a field embedding $\tau\colon k\to K$ we write $\mathrm{bc}_\tau H$ for the base change
along $\tau$.
\formal{GlobalTheory.GalV.HTower}\formal{UpperTheory.Gal.Int.Geo.Nice}
\end{definition}

\begin{definition}[Operators of correspondences; Hecke operators]\label{def:pf:steps2:hop}
A finite \'etale surjective morphism $c\colon W\to X$ with an action $a\colon G\to\Aut(W/X)$ of a
finite group ($a(1)=\mathrm{id}$, $a(gh)=a(h)\circ a(g)$) is a \emph{Galois cover} if the map
$\coprod_{g\in G}W\to W\times_XW$, $(a(g),\mathrm{id})$ on the $g$-th summand, is an
isomorphism. A finite family $\sigma_j\colon W\to Y$ ($j\in J$) \emph{splits} $s\colon Y\to X$ over
$c$ if $s\circ\sigma_j=c$ and $\coprod_{j\in J}W\to Y\times_XW$, $(\sigma_j,\mathrm{id})$, is an
isomorphism. For $s\colon Y\to Z$ and $t\colon Y\to Z'$, a $\Ql$-linear map
$T\colon\stepsiiH(Z')\to\stepsiiH(Z)$ is an \emph{operator of the correspondence $(s,t)$} if there are
a Galois cover $c\colon W\to Z$ and a family $(\sigma_j)$ splitting $s$ over $c$ such that
$c^\ast T(x)=\sum_j(t\circ\sigma_j)^\ast x$ for all $x$. For a Hecke tower $H$ over $k$, an extension
$\Omega/k$ and $g\in\Gamma$, the \emph{Hecke operator} $T_g$ on $\stepsiiH(Z_\Omega)$ is an operator
of $(s_{g,\Omega},t_{g,\Omega})$ chosen by Hilbert's choice operator if one exists, and $0$
otherwise.
\formal{GlobalTheory.GalV.IsCorrOp}\formal{GlobalTheory.GalV.hop}
\end{definition}

\begin{definition}[Eigen-pieces, matrices, graded parts]\label{def:pf:steps2:piece}
Let $H$ be a Hecke tower over $k$, $\Omega/k$ an extension, $L\supseteq\Ql$ a field and
$\theta\colon\Gamma\to L$. The \emph{$\theta$-piece} is
\[
  \stepsiiPiece_\theta=\stepsiiPiece_\theta(H,\Omega,L)=\bigcap_{g\in\Gamma}\ker\bigl(T_g\otimes L-\theta(g)\bigr)
  \subseteq L\otimes_{\Ql}\stepsiiH(Z_\Omega),
\]
on which $\Aut_k(\Omega)$ acts by $\sigma\mapsto\sigma^\ast\otimes L$ when it is stable. A homomorphism
$\rho_0\colon\Aut_k(\Omega)\to\GL_d(L)$ is a \emph{matrix of the piece} if there is a basis
$b_1,\dots,b_d$ of $\stepsiiPiece_\theta$ with $\sigma^\ast b_j=\sum_i\rho_0(\sigma)_{ij}b_i$ for all
$\sigma$ and $j$. For homomorphisms $\rho_0,\rho\colon\Gamma'\to\GL_d(L)$ of a group $\Gamma'$, $\rho$
is a \emph{graded part} of $\rho_0$ if there is $c\colon\{1,\dots,d\}\to\N$ such that for every $g$ the
matrix $\rho_0(g)$ is block upper triangular for $c$ (its $(i,j)$-entry vanishes when $c(j)<c(i)$)
and $\rho(g)_{ij}=\rho_0(g)_{ij}$ if $c(i)=c(j)$, $\rho(g)_{ij}=0$ otherwise. Finally, for $k=E$ and
$L=\Qlbar$, a homomorphism $\rho\colon G_E\to\GL_d(\Qlbar)$ is \emph{attached to} $(H,\theta)$ if it is
a graded part of some matrix of the $\theta$-piece.
\formal{UpperTheory.Gal.Int.Geo.piece}\formal{UpperTheory.Gal.Int.Geo.IsMat}\formal{UpperTheory.Gal.Int.Geo.IsGr}\formal{UpperTheory.Gal.Int.Geo.IsPiece}
\end{definition}

The constant grading $c\equiv0$ shows that every matrix of the piece is a graded part of itself.
A homomorphism $\rho\colon\Gamma'\to\GL_d(L)$ is \emph{semisimple} if every $\rho$-invariant subspace
of $L^d$ has a $\rho$-invariant complement.

\begin{definition}[Covers]\label{def:pf:steps2:cover}
Let $H$ be a Hecke tower over $k$, $\theta\colon\Gamma\to L$, and $O\to k$ a ring map. A
\emph{cover} $C$ of $(H,\theta)$ over $O$ consists of: a scheme $u\colon Y\to\Spec O$, smooth of some
relative dimension and projective (the \emph{model}); a morphism $\pi\colon Z\to Y$ with
$u\circ\pi=(\Spec k\to\Spec O)\circ q$; and a finite group $\Delta$ acting on $Z$ by
$(a_\delta)_\delta$ ($a_1=\mathrm{id}$, $a_{\delta\delta'}=a_{\delta'}\circ a_\delta$), such that the
induced morphism $\pi_k\colon Z\to Y_k:=Y\times_{\Spec O}\Spec k$ is finite and \'etale,
$\pi_k\circ a_\delta=\pi_k$, the family $(a_\delta)$ splits $\pi_k$ over $\pi_k$, and for every
$\delta$ there is $g\in\Gamma$ with $s_g$ an isomorphism, $t_g=a_\delta\circ s_g$ and $\theta(g)=1$.
The cover is \emph{onto} if $\pi_k$ is surjective. For $s,t\colon T\to Y$, $g\in\Gamma$ and $m\in\N$,
we say that \emph{$(s,t)$ lies under $g$ with multiplicity $m$} if $m>0$ and there is
$\kappa\colon Y_g\to T$ with $s\circ\kappa=\pi\circ s_g$ and $t\circ\kappa=\pi\circ t_g$ such that the
induced morphism $Y_g\to Z\times_{\pi,Y,s}T$ is finite \'etale of degree $m$. The cover is
\emph{good} if for every $g$ there are $T$ and $s,t\colon T\to Y$ with $s$ finite \'etale of some
degree, $u\circ s=u\circ t$ and $(s,t)$ lying under $g$ with some multiplicity.
\formal{UpperTheory.Gal.Int.Geo.Cover}\formal{UpperTheory.Gal.Int.Geo.Cover.Lies}\formal{UpperTheory.Gal.Int.Geo.Cover.Good}
\end{definition}

The last condition on $\Delta$ says that every deck transformation of the cover is one of the Hecke
correspondences, with eigenvalue $1$; so the $\theta$-piece consists of $\Delta$-invariant classes
(\cref{lem:pf:steps2:fixdesc}).

\begin{definition}[$p$-correspondences and the pointwise congruence relation]\label{def:pf:steps2:ptid}
Let $p$ be a prime and $Y$ a scheme, and write $Y[1/p]$ for the open subscheme where $p$ is
invertible. A \emph{$p$-correspondence} $T$ on $Y$ consists of a scheme $T_Z$, morphisms
$\alpha_T,\beta_T\colon T_Z\to Y$ and $d\in\N$ such that $\alpha_T$ is finite, \'etale over $Y[1/p]$,
and every geometric point $\Spec K\to Y$ with $p\ne0$ in $K$ has exactly $d$ lifts along $\alpha_T$.
For a morphism $y\colon X\to Y$ we write $\stepsiiSec_T(y)$ for the set of $z\colon X\to T_Z$ with
$\alpha_T\circ z=y$. For $\varpi\in\Gamma(Y,\OO_Y)$, a point $x\in Y$ is a \emph{generic point of
$\{\varpi=0\}$} if $\varpi(x)=0$ and every point $x'$ with $\varpi(x')=0$ that specializes to $x$
equals $x$. For $q,n\in\N$ and $p$-correspondences $T=(T_0,\dots,T_n)$, the \emph{pointwise
congruence relation} $\stepsiiCR_{q,\varpi}(T)$ means: for every valuation ring $V$ whose fraction
field is algebraically closed and in which $p\ne0$, every $y\colon\Spec V\to Y$ mapping the closed
point to a generic point of $\{\varpi=0\}$, the endomorphism $\phi\colon a\mapsto a^q$ of the
residue field $\kappa_V$, and all finite sets $I_j$ with bijections
$z_j\colon I_j\to\stepsiiSec_{T_j}(y)$, one has the equality of multisets of $\kappa_V$-points of $Y$
\begin{equation}\label{eq:pf:steps2:altid}
  \sum_{j\ \mathrm{even}}q^{\binom j2}\bigl\{F^{\,n-j}(\bar v_{j,i})\bigr\}_{i\in I_j}
  =\sum_{j\ \mathrm{odd}}q^{\binom j2}\bigl\{F^{\,n-j}(\bar v_{j,i})\bigr\}_{i\in I_j},
\end{equation}
where $\bar v_{j,i}\colon\Spec\kappa_V\to Y$ is the reduction of $\beta_{T_j}\circ z_j(i)$, $F$ is
precomposition with $\Spec\phi$, and $q^{\binom j2}\{\cdots\}$ denotes the multiset repeated
$q^{\binom j2}$ times.
\formal{UpperTheory.Gal.crystallineComparison.Corr}\formal{UpperTheory.Gal.crystallineComparison.PtId}
\end{definition}

\eqref{eq:pf:steps2:altid} is a pointwise form of the classical congruence relation
$\sum_j(-1)^jq^{\binom j2}\,T_j\circ\Frob^{\,n-j}=0$.

\begin{definition}[Congruence data; good reduction]\label{def:pf:steps2:cong}
Let $C$ be a cover over $O$ (\cref{def:pf:steps2:cover}), $p$ a prime, $f,n\in\N$ and
$\vartheta\colon\N\to L$. The cover \emph{carries the congruence datum $(p,f,n,\vartheta)$} if there
are an irreducible element $\varpi\in O$, $p$-correspondences $T_0,\dots,T_n$ on $Y$, elements
$\gamma_j\in\Gamma$ and natural numbers $m_j$ such that $\stepsiiCR_{p^f,u^\ast\varpi}(T)$ holds and
for every $j$: $u\circ\beta_{T_j}=u\circ\alpha_{T_j}$, $(\alpha_{T_j},\beta_{T_j})$ lies under
$\gamma_j$ with multiplicity $m_j$, and $\theta(\gamma_j)=m_j\vartheta_j$. For a Hecke tower $H$ over
$E$, $\theta\colon\Gamma\to L$ and a finite place $w$:
\begin{itemize}[nosep]
\item $\stepsiiGood_w(H,\theta)$: there is an onto good cover over $\stepsiiOw$;
  $\stepsiiGood^\circ_w(H,\theta)$: there is a good cover over $\stepsiiOw$;
\item $\stepsiiCong_w(H,\theta;n,\vartheta)$: there are a cover $C$ over $\stepsiiOw$ and a prime $p\in
  w$ with $\stepsiiNw=p^f$ such that $C$ is onto and carries the congruence datum
  $(p,f,n,\vartheta)$; $\stepsiiCong^\circ_w(H,\theta;n,\vartheta)$ is the same without ``$C$ is
  onto''.
\end{itemize}
\formal{UpperTheory.Gal.Int.Geo.Cover.Cong}\formal{UpperTheory.Gal.Int.Geo.GoodAt}\formal{UpperTheory.Gal.Int.Geo.CongAt}
\end{definition}

\begin{definition}[Hecke polynomials; the unramified condition]\label{def:pf:steps2:es}
For a commutative ring $K$, $q,n\in\N$ and $\vartheta\colon\N\to K$, and a finite multiset $R$ in
$K$,
\[
  \stepsiiHk_{q,n,\vartheta}(X)=\sum_{j=0}^n(-1)^jq^{\binom j2}\vartheta_jX^{n-j},\qquad
  \stepsiirt_R(X)=\prod_{r\in R}(X-r).
\]
Let $\iota\colon\Qlbar\to\C$ be a field isomorphism. For $\vartheta\colon\N\to\Qlbar$ and a multiset
$R\subset\C$, $R$ is an \emph{$\iota$-root multiset for $(w,n,\vartheta)$} if $\vartheta_0\ne0$ and
$\iota(\stepsiiHk_{\stepsiiNw,n,\vartheta})=\iota(\vartheta_0)\,\stepsiirt_R$. For
$\rho\colon G_E\to\GL_d(\Qlbar)$, the \emph{Eichler--Shimura condition} $\stepsiiES_w(\rho,R)$ means:
for every place $A$ of $\Omega$ over $w$, $\rho$ is trivial on $I_A$, and
$\bigl(\iota^{-1}\stepsiirt_R\bigr)(\rho(\sigma^{-1}))=0$ for every arithmetic Frobenius element
$\sigma$ for $(A,\stepsiiNw)$.
\formal{UpperTheory.Gal.crystallineComparison.heckePol}\formal{UpperTheory.Gal.Int.Geo.HasRoots}\formal{GlobalTheory.RoadWGeo.ESAt}
\end{definition}

\begin{definition}[Conditions at places above $\ell$]\label{def:pf:steps2:crys}
The $p$-adic Hodge-theoretic conditions are expressed, as in \cref{def:galois:awayb}(d), relative to
auxiliary data. A \emph{period package} with an $A_{\inf}$-structure $(P,AI)$ is \emph{tied} if
there is an isomorphism, compatible with all structure, with the concrete rings $B_{\mathrm{cris}}$
and $B_{\mathrm{dR}}$ of the development. For a valuation ring $A\subset\Omega$ and a ring $O$, a
\emph{bridge} $B_w$ consists of a finite extension $K_w/\Ql$ inside $\Qlbar$, a ring map
$j\colon\Omega\to\Qlbar$ with $j(E)\subseteq K_w$ and $K_w$ contained in the closure of $j(E)$, a
topological group isomorphism $D_A\cong\Gal(\Qlbar/K_w)$ compatible with $j$, a finite extension
$LL\supseteq K_w$ containing all conjugates of $K_w$, and a topological ring isomorphism
$O\cong\OO_{LL}$. The bridge \emph{induces} $A$ if $A=\{x\in\Omega:|j(x)|\le1\}$; an embedding
$\tau\colon E\to\Qlbar$ is \emph{induced} via a $\Ql$-embedding $\tau'\colon K_w\to LL'$ if
$\tau(x)=\tau'(j(x))$ for $x\in E$. For $r\colon D_A\to\GL_a(\Qlbar)$ and a multiset $J$ of integers,
$J$ is a \emph{jump multiset of $r$ at $\tau$} (relative to $P$ and $B_w$) if $r$ has a model over a
finite extension $LL'$ which is de Rham with respect to $P$, and for some $\tau'$ inducing $\tau$ the
multiplicity of each $n\in\Z$ in $J$ is its multiplicity as a jump of the Hodge filtration on the
$\tau'$-part of $D_{\mathrm{dR}}$. Write ``$\forall^{\mathrm t}$'' for ``for every tied $(P,AI)$, every
$A$ over $w$, every $O$ and every bridge $B_w$ for $A$ that induces $A$''. For
$\rho\colon G_E\to\GL_d(\Qlbar)$:
\begin{itemize}[nosep]
\item $\rho$ is \emph{crystalline at $w$} if, $\forall^{\mathrm t}$, $\rho|_{D_A}$ is crystalline;
\item for a predicate $Q$ on $\Hom(E,\C)\times\Z$, $\stepsiiJ_w(\rho,Q)$ means: $\forall^{\mathrm t}$
  and all $\tau,\tau'$ with $\tau$ induced via $\tau'$, there is a jump multiset $J$ of
  $\rho|_{D_A}$ at $\tau$ with $j\in J$ whenever $Q(\iota\circ\tau,j)$;
\item $\stepsiiPhi_w(\rho,R)$ means: $\forall^{\mathrm t}$, $\rho|_{D_A}$ is crystalline and the
  characteristic polynomial of some $g\in W_A$ of degree $-1$ on the Weil--Deligne representation
  attached to $\rho|_{D_A}$ (via $\iota$) divides a power of $\stepsiirt_R$.
\end{itemize}
The \emph{Hodge predicate} $\stepsiiHdg_{H,\theta}(\tau,j)$ means: $j\ge0$ and the eigensystem
$\iota\circ\theta$ occurs in some Hodge cohomology group $H^b(Z\otimes_{E,\tau}\C,\Omega^j)$.
\formal{Formalism.BasicNotions.LocBridge}\formal{GlobalTheory.Pst.IsJumpsAt}\formal{UpperTheory.Gal.Int.Div.CrysAt}\formal{UpperTheory.Gal.Int.Geo.JOcc}\formal{UpperTheory.Gal.Int.Geo.PhiK}\formal{UpperTheory.Gal.Int.Geo.HodgeAt}
\end{definition}

\begin{definition}[Lacunary weights]\label{def:pf:steps2:lac}
A finite set $J_0\subset\Z$ is \emph{lacunary} if for all disjoint nonempty $A,B\subseteq J_0$ one
has $|B|\sum_{a\in A}a\ne|A|\sum_{b\in B}b$. For a dominant weight
$\lambda=(\lambda_0\ge\dots\ge\lambda_{a-1})\in\Z^a$ put $J(\lambda)=\{k-\lambda_k:0\le k<a\}$, a set
of $a$ integers since $k\mapsto k-\lambda_k$ is strictly increasing. Let $E$ be a CM field, $U$ a
$\C$-linear space of functions on $\Uni_a(\A_{E^+})=\{g\in\GL_a(\A_E):{}^t\bar g\,g=1\}$
(\cref{sec:galois:arch}), and call $\lambda$ a \emph{weight of $U$ at $\tau\colon E\to\C$} if the
algebraic representation of highest weight $\lambda$ is a constituent of $U$ at the archimedean
place given by $\tau$. Then:
\begin{itemize}[nosep]
\item $U$ has a \emph{lacunary weight at $v$} if there are $\tau\colon E\to\C$ with
  $|\iota^{-1}(\tau(x))|<1$ for all $x\in v$, and a weight $\lambda$ of $U$ at $\tau$ with $J(\lambda)$
  lacunary;
\item $\stepsiiW_w(U,\rho)$ means: $\forall^{\mathrm t}$ and all $\tau,\tau'$ with $\tau$ induced via
  $\tau'$, all weights $\lambda$ of $U$ at $\iota\circ\tau$ and all multisets $J\subset\Z$ whose image
  in $\C$ is $\{\frac{a-1}2-z:z\in\mathrm{inf}(\lambda)\}$, where
  $\mathrm{inf}(\lambda)=\{\lambda_k+\frac{a-1}2-k\}_k$, there is a jump multiset $J'$ of
  $\rho|_{D_A}$ at $\tau$ with $J\subseteq J'$ elementwise;
\item $\rho$ is \emph{Hodge--Tate rich of rank $a$}, $\stepsiiRich_a(\rho)$, if there are a tied
  $(P,AI)$, a place $w\ni\ell$, a place $A$ over $w$, $O$ and a bridge $B_w$ inducing $A$,
  $\tau\colon E\to\Qlbar$, a jump multiset $J$ of $\rho|_{D_A}$ at $\tau$, and a lacunary set $J_0$
  of $a$ integers contained in $J$.
\end{itemize}
\formal{GlobalTheory.RoadWGeo.Lacunary}\formal{GlobalTheory.RoadWGeo.LacWtAt}\formal{UpperTheory.Gal.Int.Div.JumpsOcc}\formal{UpperTheory.Gal.Int.Div.HTRichEx}
\end{definition}

The remaining notions occur only in statements proved elsewhere and are used here only through
those statements: the twist $\stepsiitw(\rho,\chi)(g)=\chi(g)\rho(g)$ and a character $\chi$ viewed as
a $1\times1$ matrix representation $\chi^{(1)}$; the condition $\stepsiiPar_w(r;S,J)$ that $r$ has
Frobenius eigenvalues $S$ at $w\nmid\ell$, and is crystalline with jumps $J$ and Weil--Deligne
Frobenius eigenvalues $S$ at $w\mid\ell$; the variant $\stepsiiPhi^{\mathrm{sat}}_w(\rho;a,B)$ of
$\stepsiiPhi_w$ in which $\stepsiirt_R$ is replaced by the characteristic polynomials of the
unramified Weil--Deligne representation with Satake multiset $B$; Satake multisets of $U$ at $w$;
the multiset $q^{(a-1)/2}B=\{q^{(a-1)/2}\alpha:\alpha\in B\}$; Hecke characters $\chi$ of $E$
that are algebraic of infinity type $n$; and direct sums and conjugates of matrix representations.
Their definitions are recalled where the corresponding statements are proved.

\begin{definition}[The step statements]\label{def:pf:steps2:steps}
In the following list $E$, $\Omega$, $\ell$ and $\iota\colon\Qlbar\to\C$ range over number fields,
algebraic closures, primes and field isomorphisms; $H$ ranges over nice Hecke towers over $E$ with
$\theta\colon\Gamma\to\Qlbar$; ``$\rho$ admissible'' means that $\rho\colon G_E\to\GL_{dd}(\Qlbar)$ is
continuous and attached to $(H,\theta)$; and $w$ ranges over finite places of $E$.
$\sfname{Steps}$ is the conjunction of:
\begin{enumerate}[label=\textup{(S\arabic*)},leftmargin=3.2em,nosep]
\item (Existence.) For all $E,\Omega,\ell$ and nice $(H,\theta)$ there are $dd$ and a continuous
  semisimple $\rho\colon G_E\to\GL_{dd}(\Qlbar)$ attached to $(H,\theta)$.
\item (Eichler--Shimura away from $\ell$.) For $\rho$ admissible, $\ell\notin w$, $n\in\N$,
  $\vartheta\colon\N\to\Qlbar$ and $R\subset\C$: if $\stepsiiCong_w(H,\theta;n,\vartheta)$ holds and $R$
  is an $\iota$-root multiset for $(w,n,\vartheta)$, then $\stepsiiES_w(\rho,R)$.
\item (Crystallinity.) For $\rho$ admissible and $\ell\in w$: if $(H,\theta)$ has an onto cover over
  $\stepsiiOw$, then $\rho$ is crystalline at $w$.
\item (Hodge--Tate jumps.) For $\rho$ admissible and $\ell\in w$, $\stepsiiGood_w(H,\theta)$ implies
  $\stepsiiJ_w(\rho,\stepsiiHdg_{H,\theta})$.
\item (Crystalline Frobenius.) For $\rho$ admissible, $\ell\in w$, $n$, $\vartheta$ and $R$: if
  $\stepsiiCong_w(H,\theta;n,\vartheta)$ holds and $R$ is an $\iota$-root multiset for
  $(w,n,\vartheta)$, then $\stepsiiPhi_w(\rho,R)$.
\item (Cutting.) For every field $L$, nice $H$ over $E$, $\theta\colon\Gamma\to L$ and $w$:
  $\stepsiiGood^\circ_w(H,\theta)$ implies $\stepsiiGood_w(H,\theta)$, and
  $\stepsiiCong^\circ_w(H,\theta;n,\vartheta)$ implies $\stepsiiCong_w(H,\theta;n,\vartheta)$ for all
  $n,\vartheta$.
\item[(S7)--(S10)] (Twisting.) For continuous $\rho\colon G_E\to\GL_{dd}(\Qlbar)$ and
  $\chi\colon G_E\to\Qlbar^\times$, $w$, $c\in\C$ and $m\colon\Hom(E,\C)\to\Z$ such that
  $\stepsiiPar_w(\chi^{(1)};S,J)$ holds with $S=\{c\}$ and $J(\tau)=\{m(\iota\circ\tau)\}$: if
  $\ell\notin w$ then $\stepsiiES_w(\rho,R)$ implies $\stepsiiES_w(\stepsiitw(\rho,\chi),c\cdot R)$;
  if $\rho$ is crystalline at $w$ so is $\stepsiitw(\rho,\chi)$; $\stepsiiJ_w(\rho,Q)$ implies
  $\stepsiiJ_w(\stepsiitw(\rho,\chi),Q')$ with $Q'(\tau,j)=Q(\tau,j-m(\tau))$; and
  $\stepsiiPhi_w(\rho,R)$ implies $\stepsiiPhi_w(\stepsiitw(\rho,\chi),c\cdot R)$.
\item[(S11)] (Normalisation of Frobenius roots.) For $\rho$, $w$, $a\in\N$ and a multiset
  $B\subset\C^\times$ with $|B|=a$, if $\stepsiiPhi_w(\rho,R)$ holds for
  $R=\stepsiiNw^{(a-1)/2}B$, then $\stepsiiPhi^{\mathrm{sat}}_w(\rho;a,B)$.
\item[(S12)] (Satake multisets.) For $E$ CM, every Satake multiset of $U$ at $w$ has $a$ elements.
\item[(S13)] (Normalisation of jumps.) For $E$ CM, $U$, $\rho$, $w$ and a predicate $Q$: if
  $Q(\tau,j)$ for all $\tau\colon E\to\C$ with $|\iota^{-1}\tau(x)|<1$ for $x\in w$, all weights
  $\lambda$ of $U$ at $\tau$ and all $j\in J(\lambda)$, then $\stepsiiJ_w(\rho,Q)$ implies
  $\stepsiiW_w(U,\rho)$.
\item[(S14)] (Richness.) For $E$ CM, $a$, $dd$, $U$, $\rho\colon G_E\to\GL_{dd}(\Qlbar)$ and $v$ with
  $\ell\in v$: if $U$ has a lacunary weight at $v$ and $\stepsiiW_v(U,\rho)$, then
  $\stepsiiRich_a(\rho)$.
\item[(S15)] ($\ell$-adic characters.) There is a map $\mathcal L$ from pairs $(\chi,n)$, $\chi$ a Hecke
  character of $E$ and $n\colon\Hom(E,\C)\to\Z$, to homomorphisms $G_E\to\Qlbar^\times$ such that for
  $\chi$ algebraic of infinity type $n$ the character $\mathcal L(\chi,n)$ is continuous and satisfies
  $\stepsiiPar_w(\mathcal L(\chi,n)^{(1)};S,J)$ with $S=\{\chi(w)\}$ and $J(\tau)=\{-n(\iota\circ\tau)\}$
  at every $w$ where $\chi$ is unramified, and $\mathcal L(\chi,n)=\mathcal L(\chi',n')$ implies
  $\chi=\chi'$ for algebraic $(\chi,n)$, $(\chi',n')$.
\item[(S16)] (Semisimplicity of summands.) For every group and field: if a direct sum $r\oplus s$ is
  semisimple so is $r$; semisimplicity is preserved by conjugation and by twisting with a character.
\item[(S17)] (Continuity of summands.) For every topological group and topological field: if
  $r\oplus s$ is continuous so is $r$; continuity is preserved by conjugation; $\stepsiitw(r,\chi)$ and
  $\chi^{(1)}$ are continuous for continuous $r,\chi$.
\end{enumerate}
\formal{UpperTheory.Gal.Int.Geo.Steps}
\end{definition}

% ===========================================================================
\subsection{The main theorem}\label{sec:pf:steps2:assembly}

\begin{theorem}[Step lemmas]\label{thm:pf:steps2:main}
$\sfname{Steps}$ holds, i.e.\ the statements \textup{(S1)--(S17)} are true.
\formallabel{UpperTheory.Gal.Int.Geo.steps_holds}
\end{theorem}

The formal proof applies the following assembly lemma to the three results of
\cref{sec:pf:steps1}.

\begin{proposition}\label{prop:pf:steps2:stepsof}
If \textup{(S3)}, \textup{(S4)} and \textup{(S5)} hold, then $\sfname{Steps}$ holds.
\formallabel{UpperTheory.Gal.Int.Geo.steps_of}
\end{proposition}

\begin{proof}
(S3), (S4) and (S5) are the hypotheses. (S1) is \cref{thm:pf:steps2:repex}, (S2) is
\cref{thm:pf:steps2:esstep}, (S6) is \cref{thm:pf:steps2:cut} and (S14) is
\cref{thm:pf:steps2:richof}. The remaining statements are proved elsewhere:
\lref{UpperTheory.Gal.Int.Geo.twES} states (S7), that twisting by a character with parameters
$(\{c\},m)$ preserves the Eichler--Shimura condition away from $\ell$ with the roots multiplied by
$c$; \lref{UpperTheory.Gal.Int.Geo.twCrys} states (S8), that such twists preserve
crystallinity at $w$; \lref{UpperTheory.Gal.Int.Geo.twJ} states (S9), that they shift the jumps
by $m$; \lref{UpperTheory.Gal.Int.Geo.twPhi} states (S10), that they multiply the crystalline
Frobenius roots by $c$; \lref{UpperTheory.Gal.Int.Geo.phiOf} states (S11), that Frobenius roots
$\stepsiiNw^{(a-1)/2}B$ give the Satake form of the Frobenius condition;
\lref{UpperTheory.Gal.Int.Geo.satCard} states (S12), that Satake multisets of $U$ at $w$ have $a$
elements; \lref{UpperTheory.Gal.Int.Geo.jumpsOf} states (S13), that occurrence of all weight
jumps among the jumps of $\rho$ gives $\stepsiiW_w(U,\rho)$;
\lref{UpperTheory.Gal.TP.Sub.charEx} states (S15), the existence of the $\ell$-adic characters of
algebraic Hecke characters with their local properties (the construction of \cite{Serre1968});
\lref{UpperTheory.Gal.TP.Sub.sSLeft_holds} and \lref{UpperTheory.Gal.TP.Sub.contLeft_holds}
state (S16) and (S17), the semisimplicity and continuity of summands, conjugates and twists.
\end{proof}

\begin{proof}[Proof of \cref{thm:pf:steps2:main}]
Apply \cref{prop:pf:steps2:stepsof} to (S3), which is
\lref{UpperTheory.Gal.Seam.crysStep_holds} (for admissible $\rho$ and $\ell\in w$, an onto cover
over $\stepsiiOw$ makes $\rho$ crystalline at $w$); to (S4), which is
\lref{UpperTheory.Gal.Int.Geo.jumpStep_holds} (for admissible $\rho$ and $\ell\in w$,
$\stepsiiGood_w(H,\theta)$ implies that the Hodge degrees occur among the jumps of $\rho$); and to
(S5), which is \lref{UpperTheory.Ru.crystallineComparison.phiStep} (for admissible $\rho$ and
$\ell\in w$, a congruence datum and an $\iota$-root multiset $R$ give $\stepsiiPhi_w(\rho,R)$).
\end{proof}

\begin{remark}\label{rem:pf:steps2:scope}
All step statements concern representations attached to $(H,\theta)$; if the $\theta$-piece is zero,
then $dd=0$ and the local conclusions hold trivially. The non-degeneracy needed in the proof of
\cref{thm:geoone} enters through the Hodge-theoretic hypotheses there, via (S4) and (S14). The
conclusion of (S2) is the annihilation of $\rho(\sigma^{-1})$ by $\iota^{-1}\stepsiirt_R$, not an
equality of characteristic polynomials.
\end{remark}


% ===========================================================================
\subsection{Lacunary weights: statement (S14)}\label{sec:pf:steps2:rich}

\begin{lemma}[Bridges]\label{lem:pf:steps2:bridge}
Let $j\colon\Omega\to\Qlbar$ be a ring map and $A=\{x\in\Omega:|j(x)|\le1\}$. Let $K_w$ be the subfield
of $\Qlbar$ generated over $\Ql$ by $j(E)$ and $LL$ its normal closure over $\Ql$ in $\Qlbar$. Then
$K_w/\Ql$ is finite, and there is a bridge for $A$ with field $K_w$, embedding $j$, field $LL$ and
ring $O=\OO_{LL}$; it induces $A$.
\formallabel{GlobalTheory.PhtCompare.exists_bridge}
\end{lemma}

\begin{proof}
$K_w$ is generated over $\Ql$ by the images of a $\Q$-basis of the number field $E$, which are
algebraic over $\Ql$, so $K_w/\Ql$ is finite, and $LL$ is a finite extension containing $\tau(K_w)$
for every $\Ql$-embedding $\tau\colon K_w\to\Qlbar$. Clearly $j(E)\subseteq K_w$. The closure of $j(E)$
in $\Qlbar$ is a subfield containing $j(E)$ and the closure $\Ql$ of $\Q$, hence contains $K_w$. By
\lref{GlobalTheory.PlaceBridge.exists_mulEquiv}, which states that for $A$ and $K_w$ as here
there is an isomorphism of topological groups $D_A\cong\Gal(\Qlbar/K_w)$, $g\mapsto e(g)$, with
$e(g)(j(x))=j(g(x))$ for $x\in\Omega$, the remaining datum of a bridge exists; take for
$O\cong\OO_{LL}$ the identity. The bridge induces $A$ by the definition of $A$.
\end{proof}

\begin{theorem}[Richness]\label{thm:pf:steps2:richof}
Let $E$ be a CM field, $\iota\colon\Qlbar\to\C$ a field isomorphism, $a,dd\in\N$, $U$ a space of
functions on $\Uni_a(\A_{E^+})$, $\rho\colon G_E\to\GL_{dd}(\Qlbar)$ a homomorphism and $v$ a finite
place of $E$ with $\ell\in v$. If $U$ has a lacunary weight at $v$ and $\stepsiiW_v(U,\rho)$ holds,
then $\rho$ is Hodge--Tate rich of rank $a$.
\formallabel{UpperTheory.Gal.Int.Geo.richOf}
\end{theorem}

\begin{proof}
By hypothesis there are $\tau_c\colon E\to\C$ with $|\iota^{-1}\tau_c(x)|<1$ for $x\in v$ and a
weight $\lambda$ of $U$ at $\tau_c$ with $J(\lambda)$ lacunary. By
\lref{GlobalTheory.MemLoc.exists_locPrime_tied}, which states that for every prime $\ell$ there
is a tied period package with an $A_{\inf}$-structure, we may fix a tied $(P,AI)$. Put
$\tau_0=\iota^{-1}\circ\tau_c\colon E\to\Qlbar$, so that $\iota\circ\tau_0=\tau_c$ and
$|\tau_0(x)|<1$ for $x\in v$.

\emph{Step 1: for $x\in\OO_E$ one has $|\tau_0(x)|<1$ if and only if $x\in v$.} One direction is the
hypothesis. If $x\notin v$, maximality of $v$ gives $y\in\OO_E$ and $i\in v$ with $yx+i=1$. As $y$
is integral over $\Z$, $|\tau_0(y)|\le1$; hence $|\tau_0(yx)|<1$ and $|\tau_0(i)|<1$, and the
ultrametric inequality gives $1=|\tau_0(yx)+\tau_0(i)|<1$, a contradiction.

\emph{Step 2: a place and a bridge.} Since $\Omega$ is algebraic over $E$ and $\Qlbar$ is
algebraically closed, $\tau_0$ extends to a ring map $j\colon\Omega\to\Qlbar$ with $j|_E=\tau_0$. The
set $A=\{x\in\Omega:|j(x)|\le1\}$ is a valuation ring of $\Omega$ (the absolute value is
ultrametric and multiplicative). By \cref{lem:pf:steps2:bridge} there is a bridge $B_w$ for $A$, with
ring $O=\OO_{LL}$, field $K_w=\Ql(j(E))$ and embedding $j$, which induces $A$. The place $A$ lies
over $v$: for $x\in\OO_E$,
$v_A(x)<1$ if and only if $|j(x)|<1$, i.e.\ $|\tau_0(x)|<1$, i.e.\ $x\in v$ by Step~1. Finally
$\tau_0$ is induced via the identity of $K_w$, since $\tau_0(x)=j(x)$ for $x\in E$.

\emph{Step 3.} Apply $\stepsiiW_v(U,\rho)$ to $(P,AI)$, $A$, $O$, $B_w$, $\tau_0$, $LL'=K_w$,
$\tau'=\mathrm{id}$, the weight $\lambda$ of $U$ at $\iota\circ\tau_0=\tau_c$, and the multiset
$J=J(\lambda)$, whose image in $\C$ is $\{\frac{a-1}2-z:z\in\mathrm{inf}(\lambda)\}$ because
$\frac{a-1}2-(\lambda_k+\frac{a-1}2-k)=k-\lambda_k$ (this computation is
\lref{GlobalTheory.RoadWGeo.wtJumps_val}). We obtain a jump multiset $J'$ of $\rho|_{D_A}$ at
$\tau_0$ with $J(\lambda)\subseteq J'$. Hence $\rho$ is Hodge--Tate rich of rank $a$, with the
witnesses $(P,AI)$, $w=v$ (which contains $\ell$), $A$, $O$, $B_w$, $\tau_0$, $J'$ and
$J_0=J(\lambda)$, which has $a$ elements (\lref{GlobalTheory.RoadWGeo.card_wtJumps}, the
injectivity of $k\mapsto k-\lambda_k$) and is lacunary.
\end{proof}

% ===========================================================================
\subsection{Cutting covers: statement (S6)}\label{sec:pf:steps2:cut}

In this subsection $O$ is a discrete valuation ring with fraction field $k$, and a point of an
$O$-scheme $u\colon Y\to\Spec O$ is \emph{generic} if it lies over the generic point of $\Spec O$. A
morphism $j\colon Y'\to Y$ is \emph{a part} if it is an open immersion with closed image. For a part
$j$ and $s,t\colon T\to Y$, the \emph{restriction} of $(s,t)$ to $Y'$ is the open subscheme
$T'=s^{-1}(j(Y'))\cap t^{-1}(j(Y'))$ of $T$, with inclusion $i\colon T'\to T$, and the factorizations
$s',t'\colon T'\to Y'$ with $j\circ s'=s\circ i$ and $j\circ t'=t\circ i$.

\begin{lemma}\label{lem:pf:steps2:genfib}
Let $u\colon Y\to\Spec O$ be a morphism.
\begin{enumerate}[label=\textup{(\alph*)},nosep]
\item The projection $Y_k\to Y$ is a topological embedding whose image is the set of generic points.
\item If $u$ is flat, the generic points are dense in $Y$.
\item If $u$ is smooth, the local rings of $Y$ are domains.
\end{enumerate}
\formallabel{UpperTheory.Gal.Int.Geo.Cut.genRange}\formallabel{UpperTheory.Gal.Int.Geo.Cut.genDense}\formallabel{UpperTheory.Gal.Int.Geo.Cut.stalkDom}
\end{lemma}

\begin{proof}
(a) $\Spec k\to\Spec O$ is a localization, hence a preimmersion, and so is its base change; a point
$y$ is in the image if and only if $u(y)$ is in the image of $\Spec k$, i.e.\ is the zero ideal,
because $O\to k$ is injective. (b) Flat morphisms are generalizing: for $y$ in a nonempty open $W$,
the specialization $(0)\rightsquigarrow u(y)$ lifts to $y'\rightsquigarrow y$ with $u(y')=(0)$, and
$y'\in W$ since open sets are stable under generization. (c) A scheme smooth over a regular scheme
is regular, and regular local rings are domains \cite{Grothendieck1967}.
\end{proof}

\begin{lemma}\label{lem:pf:steps2:clopen}
Let $u\colon Y\to\Spec O$ be smooth and let $U\subseteq Y_k$ be open and closed. Then there is an
open and closed $V\subseteq Y$ whose preimage in $Y_k$ is $U$.
\formallabel{UpperTheory.Gal.Int.Geo.Cut.clopenExt}
\end{lemma}

\begin{proof}
\emph{Step 1: every point of $Y$ has an irreducible open neighbourhood.} $Y$ is locally noetherian.
Let $y\in W=\Spec R$ with $R$ noetherian; by \cref{lem:pf:steps2:genfib}(c) the local rings of $W$
are domains, so two distinct irreducible components of $W$ cannot meet (a common point would have
two minimal primes in its local ring). The finitely many components are therefore pairwise
disjoint, hence open, and the component containing $y$ is the required neighbourhood.

\emph{Step 2: a topological lemma.} Let $f\colon P\to X$ be a topological embedding with dense
image such that every point of $X$ has an open preirreducible neighbourhood, and let $U\subseteq P$
be open and closed. Write $U=f^{-1}(A_1)$ and $P\smallsetminus U=f^{-1}(A_2)$ with $A_1,A_2$ open.
For every open preirreducible $W\subseteq X$ one has $f^{-1}(W)\subseteq U$ or $f^{-1}(W)\subseteq
P\smallsetminus U$: otherwise $W\cap A_1$ and $W\cap A_2$ are nonempty, so $W\cap A_1\cap
A_2\ne\emptyset$ by preirreducibility, and by density it contains some $f(q)$, with $q\in
U\cap(P\smallsetminus U)$. Let $V$ be the union of the open preirreducible $W$ with
$f^{-1}(W)\subseteq U$. It is open. It is closed: if $x\notin V$, its preirreducible neighbourhood
$W_x$ satisfies $f^{-1}(W_x)\subseteq P\smallsetminus U$, and $W_x\cap V=\emptyset$, since a point of
$W_x\cap W$ with $f^{-1}(W)\subseteq U$ would give, by density, $f(q)\in W_x\cap W$ with $q\in U$ and
$q\notin U$. Finally $f^{-1}(V)=U$: the inclusion $\subseteq$ is clear, and for $q\in U$ a
preirreducible neighbourhood $W$ of $f(q)$ satisfies $f^{-1}(W)\subseteq U$, since $q\in
f^{-1}(W)\cap U$.

\emph{Step 3.} Apply Step~2 to $f\colon Y_k\to Y$, an embedding with dense image by
\cref{lem:pf:steps2:genfib}(a),(b) (smooth morphisms are flat), using Step~1.
\end{proof}

\begin{lemma}\label{lem:pf:steps2:covercut}
Let $C$ be a cover of $(H,\theta)$ over a ring $O\to k$, and let $V\subseteq Y$ be open and closed
with preimage in $Y_k$ equal to $\pi_k(Z)$. Then there is an onto cover $C'$ with model $V$, the
inclusion $j\colon V\to Y$ being a part, $u'=u\circ j$, $\pi=j\circ\pi'$, and the same $\Delta$ and
$(a_\delta)$.
\formallabel{UpperTheory.Gal.Int.Geo.Cut.coverCut}
\end{lemma}

\begin{proof}
Since $\pi_k(Z)$ lies in the preimage of $V$, $\pi(Z)\subseteq V$ and $\pi$ factors as $j\circ\pi'$.
The base change $m\colon V_k\to Y_k$ of $j$ is an open immersion with $\pi_k=m\circ\pi'_k$. Hence
$\pi'_k$ is finite and \'etale (as $\pi_k$ is, and $m$ is a separated \'etale monomorphism), and
the splitting of $\pi_k$ by $(a_\delta)$ is a splitting of $\pi'_k$, because
$Z\times_{V_k}Z\to Z\times_{Y_k}Z$ is an isomorphism ($m$ is a monomorphism). $V\to\Spec O$ is
projective (compose the closed immersion $V\to Y$ with an embedding of $Y$ in $\mathbb P^N_O$) and
smooth of the same relative dimension (an open subscheme). The conditions on $(a_\delta)$ and
$\theta$ do not involve $Y$. Finally $\pi'_k$ is surjective: a point $x$ of $V_k$ maps to a point of
$Y_k$ lying over $V$, hence in $\pi_k(Z)$, say $m(x)=\pi_k(z)=m(\pi'_k(z))$, and $m$ is injective.
\end{proof}

\begin{lemma}\label{lem:pf:steps2:cut}
Let $O$ be a discrete valuation ring with fraction field $k$ and $C$ a cover over $O$. There are an
onto cover $C'$ and a part $j\colon Y'\to Y$ with $u'=u\circ j$ and $\pi=j\circ\pi'$, such that every
generic point of $Y$ in $j(Y')$ lies in $\pi(Z)$.
\end{lemma}

\begin{proof}
$\pi_k(Z)$ is closed and open in $Y_k$ ($\pi_k$ is finite, hence closed, and \'etale, hence open).
By \cref{lem:pf:steps2:clopen} there is $V$ open and closed in $Y$ with preimage $\pi_k(Z)$, and
\cref{lem:pf:steps2:covercut} gives $C'$ and $j$. If $y\in V$ is generic, then $y$ is the image of a
point of $Y_k$ (\cref{lem:pf:steps2:genfib}(a)) lying in the preimage of $V$, i.e.\ in $\pi_k(Z)$;
hence $y\in\pi(Z)$.
\end{proof}

\begin{lemma}\label{lem:pf:steps2:restrict}
Let $C,C'$ be covers of $(H,\theta)$ over $O$ and $j\colon Y'\to Y$ a part with $\pi=j\circ\pi'$. Let
$s,t\colon T\to Y$, with restriction $(T',i,s',t')$.
\begin{enumerate}[label=\textup{(\alph*)},nosep]
\item If $(s,t)$ lies under $g$ with multiplicity $m$ (for $C$), then $s(x)\in\pi(Z)$ implies
  $t(x)\in\pi(Z)$ for $x\in T$, and $(s',t')$ lies under $g$ with multiplicity $m$ for $C'$.
\item If $s(x)\in j(Y')$ implies $t(x)\in j(Y')$ for all $x\in T$, and $s$ is finite \'etale of
  degree $d$, then so is $s'$.
\item Let $T$ be a $p$-correspondence on $Y$ such that $\beta_T(x)\in j(Y')$ for every $x\in T_Z$
  with $\alpha_T(x)\in j(Y')$ and $p$ invertible at $\alpha_T(x)$. Then the restriction
  $(T',s',t',d_T)$ of $(\alpha_T,\beta_T)$ is a $p$-correspondence on $Y'$; and if $T_0,\dots,T_n$
  satisfy this condition and $\stepsiiCR_{q,\varpi}(T)$, then $\stepsiiCR_{q,j^\ast\varpi}(T')$.
\end{enumerate}
\formallabel{UpperTheory.Gal.Int.Geo.Cut.liesCut}\formallabel{UpperTheory.Gal.Int.Geo.Cut.fEtCut}\formallabel{UpperTheory.Gal.Int.Geo.Cut.corrCut}\formallabel{UpperTheory.Gal.Int.Geo.Cut.ptIdCut}\formal{UpperTheory.Gal.Int.Geo.Cut.resEx}
\end{lemma}

\begin{proof}
(a) Let $\kappa\colon Y_g\to T$ be as in \cref{def:pf:steps2:cover}. If $s(x)=\pi(z)$, the fibre
product $Z\times_YT$ has a point over $(z,x)$; the morphism $Y_g\to Z\times_YT$ is finite \'etale of
degree $m>0$, hence surjective (every geometric point has $m$ lifts), so this point is the image of
some $y\in Y_g$, and then $\kappa(y)=x$ and $t(x)=\pi(t_g(y))$. In particular $\kappa(Y_g)\subseteq
T'$, as $s\kappa=\pi s_g$ and $t\kappa=\pi t_g$ take values in $\pi(Z)\subseteq j(Y')$; let
$\kappa'\colon Y_g\to T'$ be the factorization. The identities $s'\kappa'=\pi's_g$ and
$t'\kappa'=\pi't_g$ hold after composition with the monomorphism $j$. By the first assertion, the
second projection of $Z\times_{\pi,Y,s}T$ lands in $T'$, so the canonical map
$Z\times_{\pi',Y',s'}T'\to Z\times_{\pi,Y,s}T$ is an isomorphism compatible with the maps from $Y_g$;
therefore $Y_g\to Z\times_{Y'}T'$ is finite \'etale of degree $m$.

(b) Under the hypothesis $T'=s^{-1}(j(Y'))$, so $T'=T\times_{s,Y,j}Y'$ with $s'$ the base change of
$s$ along $j$. Finiteness and \'etaleness are stable under base change, and the lifts of a
geometric point $y'$ of $Y'$ along $s'$ correspond bijectively to the lifts of $j\circ y'$ along $s$.

(c) $T'$ is closed in $T_Z$ (the preimage of the closed set $j(Y')$ under $\alpha_T$ and $\beta_T$),
so $i$ is a closed immersion and $s'$ is finite, since $j\circ s'=\alpha_T\circ i$ is. Over
$Y'[1/p]=j^{-1}(Y[1/p])$ the morphism $s'$ is the restriction of $\alpha_T$ over $Y[1/p]$ to an open
subscheme, hence \'etale. If $y'\colon\Spec K\to Y'$ is a geometric point with $p\ne0$ in $K$, a lift
$z$ of $j\circ y'$ along $\alpha_T$ has $\alpha_T(z)\in j(Y')\cap Y[1/p]$, so $\beta_T(z)\in j(Y')$
and $z$ factors through $T'$; thus lifts of $y'$ along $s'$ and of $j\circ y'$ along $\alpha_T$
correspond, and their number is $d_T$.

For the congruence relation, let $V$, $y'\colon\Spec V\to Y'$, $\phi$ and bijections
$z'_a\colon I_a\to\stepsiiSec_{T'_a}(y')$ be as in \cref{def:pf:steps2:ptid}, the closed point of $y'$
being a generic point of $\{j^\ast\varpi=0\}$. Then the closed point of $j\circ y'$ is a generic point
of $\{\varpi=0\}$: a point $x'$ of $Y$ with $\varpi(x')=0$ specializing to a point of the open set
$j(Y')$ lies in $j(Y')$, and $j$ is a homeomorphism onto its image. Composition with $i_a$ maps
$\stepsiiSec_{T'_a}(y')$ bijectively onto $\stepsiiSec_{T_a}(j\circ y')$: it is injective as $i_a$ is a
monomorphism; and a section $z$ of $\alpha_{T_a}$ over $j\circ y'$ factors through $T'_a$, because
the generic point $\eta$ of $\Spec V$ maps to $Y[1/p]$ ($p\ne0$ in $V$), so that
$\beta_{T_a}(z(\eta))\in j(Y')$, and every point of $\Spec V$ specializes from $\eta$ while $j(Y')$ is
closed. Applying $\stepsiiCR_{q,\varpi}(T)$ to $j\circ y'$ and the bijections $i_a\circ z'_a$ gives
\eqref{eq:pf:steps2:altid} for the points $j\circ\bar v'_{a,i}$. Since composition with $j$ is
injective on $\kappa_V$-points and commutes with precomposition by $\Spec\phi$,
\eqref{eq:pf:steps2:altid} holds for the $\bar v'_{a,i}$.
\end{proof}

\begin{theorem}[Cutting]\label{thm:pf:steps2:cut}
Let $E$ be a number field, $L$ a field, $H$ a nice Hecke tower over $E$, $\theta\colon\Gamma\to L$ and
$w$ a finite place of $E$. Then $\stepsiiGood^\circ_w(H,\theta)$ implies $\stepsiiGood_w(H,\theta)$,
and $\stepsiiCong^\circ_w(H,\theta;n,\vartheta)$ implies $\stepsiiCong_w(H,\theta;n,\vartheta)$ for all
$n\in\N$ and $\vartheta\colon\N\to L$.
\formallabel{UpperTheory.Gal.Int.Geo.cutStep_holds}
\end{theorem}

\begin{proof}
Let $C$ be a cover over the discrete valuation ring $O=\stepsiiOw$ and let $C'$, $j$ be as in
\cref{lem:pf:steps2:cut}.

\emph{Good covers.} Let $C$ be good and $g\in\Gamma$, with $(T,s,t)$ such that $s$ is finite \'etale
of degree $d$, $us=ut$ and $(s,t)$ lies under $g$ with multiplicity $m$. We claim that
$s(x)\in j(Y')$ implies $t(x)\in j(Y')$ for $x\in T$. The set of $x$ with $s(x)\in j(Y')$ and
$t(x)\notin j(Y')$ is open ($j(Y')$ is open and closed). As $u\circ s$ is flat (\'etale followed by
smooth), if this set is nonempty it contains a point $x$ with $u(s(x))$ generic
(\cref{lem:pf:steps2:genfib}(b)); then $s(x)\in\pi(Z)$ by \cref{lem:pf:steps2:cut}, and
$t(x)\in\pi(Z)\subseteq j(Y')$ by \cref{lem:pf:steps2:restrict}(a), a contradiction. Let
$(T',s',t')$ be the restriction. By \cref{lem:pf:steps2:restrict}(b) $s'$ is finite \'etale of degree
$d$; $u's'=u't'$ follows from $us=ut$; and $(s',t')$ lies under $g$ with multiplicity $m$ for $C'$ by
\cref{lem:pf:steps2:restrict}(a). Thus $C'$ is an onto good cover.

\emph{Congruence data.} Let $p\in w$ be prime with $\stepsiiNw=p^f$, and let
$(\varpi,T,\gamma,m)$ witness that $C$ carries the congruence datum $(p,f,n,\vartheta)$. A point
$y\in Y$ at which $p$ is invertible is generic: $p$ lies in the maximal ideal of $\stepsiiOw$ (as
$p\in w$), so the prime $u(y)$, which does not contain $p$, is not maximal, hence is $(0)$.
Therefore, for $x\in(T_a)_Z$ with $\alpha_{T_a}(x)\in j(Y')$ and $p$ invertible at $\alpha_{T_a}(x)$,
the point $\alpha_{T_a}(x)$ is generic, hence in $\pi(Z)$ (\cref{lem:pf:steps2:cut}), and
$\beta_{T_a}(x)\in\pi(Z)\subseteq j(Y')$ by \cref{lem:pf:steps2:restrict}(a) applied to
$(\alpha_{T_a},\beta_{T_a})$, which lies under $\gamma_a$ with multiplicity $m_a$. By
\cref{lem:pf:steps2:restrict}(c) the restrictions $T'_a$ are $p$-correspondences on $Y'$ satisfying
$\stepsiiCR_{p^f,j^\ast u^\ast\varpi}(T')$, and $j^\ast u^\ast\varpi=u'^\ast\varpi$. The conditions
$u'\beta_{T'_a}=u'\alpha_{T'_a}$, $(\alpha_{T'_a},\beta_{T'_a})$ lying under $\gamma_a$ with
multiplicity $m_a$ (\cref{lem:pf:steps2:restrict}(a)), and $\theta(\gamma_a)=m_a\vartheta_a$ hold.
Hence the onto cover $C'$ carries the congruence datum $(p,f,n,\vartheta)$.
\end{proof}


% ===========================================================================
\subsection{Galois representations on eigen-pieces: statement (S1)}\label{sec:pf:steps2:repex}

\subsubsection{Linear algebra}

\begin{lemma}[Semisimplification]\label{lem:pf:steps2:grex}
Let $\Gamma'$ be a group, $k$ a field and $\rho_0\colon\Gamma'\to\GL_d(k)$ a homomorphism. There are
$P\in\GL_d(k)$ and a semisimple homomorphism $\rho\colon\Gamma'\to\GL_d(k)$ which is a graded part
of $P^{-1}\rho_0P$.
\formallabel{UpperTheory.Gal.Int.Geo.grEx}
\end{lemma}

\begin{proof}
For $c\colon\{1,\dots,d\}\to\N$ and $i\in\N$ put $\mathrm{low}(c,i)=\{v\in k^d:v_j=0\text{ whenever
}c(j)\ge i\}$, the span of the $e_j$ with $c(j)<i$. Call a sequence $(W_i)_{i\in\N}$ of
$\rho_1$-invariant subspaces a \emph{composition flag} for $\rho_1$ if for no $i$ there is a
$\rho_1$-invariant subspace strictly between $W_i$ and $W_{i+1}$.

(i) There is a monotone composition flag $(W_i)$ for $\rho_0$ with $W_0=0$ and $W_d=k^d$: build a
composition series by induction (each proper step increases the dimension) and pad it with
repetitions.

(ii) There are $P\in\GL_d(k)$ and $c$ with $P(\mathrm{low}(c,i))=W_i$ for all $i$: choose a basis
adapted to the flag, i.e.\ extend successively bases of $W_0\subseteq W_1\subseteq\cdots$, let $P$
have these vectors as columns and let $c(j)$ be the index $i$ with $Pe_j\in W_{i+1}\smallsetminus
W_i$.

(iii) Then $(\mathrm{low}(c,i))_i$ is a composition flag for $\rho_1=P^{-1}\rho_0P$, since $P^{-1}$
maps $\rho_0$-invariant subspaces bijectively onto $\rho_1$-invariant ones.

(iv) Invariance of all $\mathrm{low}(c,i)$ under $\rho_1(g)$ means that the entries $(i,j)$ of
$\rho_1(g)$ with $c(j)<c(i)$ vanish, i.e.\ $\rho_1(g)$ is block upper triangular for $c$. The
block-diagonal part of a product of block upper triangular matrices is the product of their
block-diagonal parts, so the block-diagonal part $\rho(g)$ of $\rho_1(g)$ is a homomorphism, and
$\rho$ is a graded part of $\rho_1$.

(v) $\rho$ is semisimple. The subspaces $B_i=\mathrm{span}\{e_j:c(j)=i\}$ are $\rho$-invariant and
$k^d=\bigoplus_iB_i$. If $U\subseteq B_i$ is $\rho$-invariant, then $\mathrm{low}(c,i)+U$ is
$\rho_1$-invariant (the action of $\rho_1$ on $\mathrm{low}(c,i+1)/\mathrm{low}(c,i)\cong B_i$ is the
block $\rho|_{B_i}$) and lies between $\mathrm{low}(c,i)$ and $\mathrm{low}(c,i+1)$; by (iii) it
equals one of them, so $U=0$ or $U=B_i$. A representation that is the sum of invariant subspaces
each of which is zero or irreducible is semisimple.
\end{proof}

\begin{lemma}[Matrices of stable subspaces]\label{lem:pf:steps2:matex}
Let a group $\Gamma'$ act linearly on an $L$-vector space $M$ by $a\colon\Gamma'\to\End_L(M)$ with
$a(1)=\mathrm{id}$ and $a(gh)=a(g)a(h)$, and let $U\subseteq M$ be an $a$-stable subspace with a
basis $b_1,\dots,b_d$. Then there is a homomorphism $\rho\colon\Gamma'\to\GL_d(L)$ with
$a(g)b_j=\sum_i\rho(g)_{ij}b_i$, and for every $(i,j)$ a linear form $\mu_{ij}$ on $M$ with
$\rho(g)_{ij}=\mu_{ij}(a(g)b_j)$ for all $g$. If $\rho_0$ is a matrix of a piece
(\cref{def:pf:steps2:piece}) and $P\in\GL_d(L)$, then $P^{-1}\rho_0P$ is a matrix of the same piece.
\formallabel{UpperTheory.Gal.Int.Geo.matEx_of}\formallabel{UpperTheory.Gal.Int.Geo.matConj_of}
\end{lemma}

\begin{proof}
Let $\rho(g)$ be the matrix of $a(g)|_U$ in the basis $b$; it is invertible with inverse
$\rho(g^{-1})$, and multiplicative because $a$ is. Let $\mu_{ij}$ be any extension to $M$ of the
$i$-th coordinate form on $U$. For the second assertion use the basis $b'_j=\sum_iP_{ij}b_i$.
\end{proof}

\subsubsection{Galois stability of the pieces}

\begin{lemma}\label{lem:pf:steps2:hopgal}
Let $H$ be a Hecke tower over a field $k$ in which $\ell\ne0$, $\Omega/k$ an extension, and
$g\in\Gamma$ with $s_g$ finite \'etale of some degree. Then $T_g\circ\sigma^\ast=\sigma^\ast\circ T_g$
on $\stepsiiH(Z_\Omega)$ for every $\sigma\in\Aut_k(\Omega)$. Consequently, if every $s_g$ is finite
\'etale of some degree, then every $\stepsiiPiece_\theta(H,\Omega,L)$ is stable under $\Aut_k(\Omega)$.
\formallabel{UpperTheory.Gal.Int.Geo.hopGal}\formallabel{UpperTheory.Gal.Int.Geo.pieceStab}
\end{lemma}

\begin{proof}
(1) $T_g$ is an operator of the correspondence $(s_{g,\Omega},t_{g,\Omega})$: by definition of $T_g$ it
suffices that some operator exists, and $s_{g,\Omega}$, being finite \'etale of constant degree, is
split by a finite \'etale Galois cover (\cref{lem:pf:steps2:galclos}(b)), over which $\stepsiitr_s\circ t^\ast$ satisfies the defining
identity (compare \cref{lem:pf:steps2:hopop}).
(2) If $T$ is an operator of $(s_{g,\Omega},t_{g,\Omega})$, so is $\sigma^\ast\circ
T\circ(\sigma^{-1})^\ast$: transport the Galois cover and the splitting by the automorphisms
$\sigma_{(-)}$, which commute with $s_{g,\Omega}$ and $t_{g,\Omega}$.
(3) Any two operators of $(s_{g,\Omega},t_{g,\Omega})$ coincide, by \cref{lem:pf:steps2:opuniq}.
By (1)--(3), $\sigma^\ast T_g(\sigma^{-1})^\ast=T_g$, which is the first claim. If $x\in
\stepsiiPiece_\theta$, then $T_g\sigma^\ast x=\sigma^\ast T_gx=\theta(g)\sigma^\ast x$ for every $g$.
\end{proof}

\subsubsection{Finiteness and continuity}

\begin{lemma}[Finiteness]\label{lem:pf:steps2:fin}
Let $k$ be an algebraically closed field with $\ell\ne0$ in $k$ and $X$ a smooth projective
$k$-scheme. Then every $H^n(X,\Ql)$ is finite-dimensional, there is $N$ with $H^n(X,\Ql)=0$ for $n<0$
and $n>N$, and hence $\stepsiiH(X)$ is finite-dimensional over $\Ql$.
\formallabel{UpperTheory.Gal.Int.Geo.degFin}\formallabel{UpperTheory.Gal.Int.Geo.htotFin}
\end{lemma}

\begin{proof}
$X$ is proper, noetherian and normal, so it is the disjoint union of finitely many open and closed
integral subschemes $Y_1,\dots,Y_r$ (its connected components), each smooth of some relative
dimension $d_s$ and projective (being closed in $X$). By Poincar\'e duality for smooth projective
integral varieties over algebraically closed fields
\cite[Exp.~XVIII]{ArtinGrothendieckVerdier1972}, \cite[Ch.~VI, \S11]{Milne1980}, which the
development derives from the package of \cref{prop:pf:steps2:etfound}(k): $H^n(Y_s,\Ql)$ is
finite-dimensional for all $n$, vanishes for $n>2d_s$, and there is an isomorphism
$\mathrm{tr}\colon H^{2d_s}(Y_s,\Ql)\cong\Ql$ for which the cup product pairings
$H^n\times H^{2d_s-n}\to\Ql$ are perfect. The cohomology of a finite disjoint union of open
subschemes is the product of the cohomologies at every finite level, hence after passing to the
inverse limit and tensoring with $\Ql$ (which commutes with finite products); so
$H^n(X,\Ql)\cong\prod_sH^n(Y_s,\Ql)$. This gives finiteness and vanishing for $n>2\max_sd_s$. For
$n<0$ already $H^n(X,\Z/\ell^m)=0$ for all $m$ (\cref{prop:pf:steps2:etfound}(j)). Finally a direct
sum of finite-dimensional spaces with finitely many nonzero terms is finite-dimensional.
\end{proof}

\begin{lemma}[Continuity]\label{lem:pf:steps2:cont}
Let $k$ be a field with $\ell\ne0$ in $k$, $\Omega$ an algebraic closure of $k$ and $X$ a proper
$k$-scheme. Then for every $v\in\stepsiiH(X_\Omega)$ and every linear form $\mu$ on
$\stepsiiH(X_\Omega)$ the map $\Aut_k(\Omega)\to\Ql$, $\sigma\mapsto\mu(\sigma^\ast v)$, is continuous.
\formallabel{UpperTheory.Gal.Int.Geo.totCont}
\end{lemma}

\begin{proof}
(a) \emph{Descent.} For $X$ quasi-compact and quasi-separated over $k$ and $\Omega/k$ algebraic,
every class in $H^n(X_\Omega,\Z/q(i))$ is the pull-back of a class on $X_L$ for some finite
subextension $L/k$. Indeed $X_\Omega$ is the limit of the $X_L$ along affine transition maps, and
\'etale cohomology with coefficients in $\Z/q(i)$ commutes with such limits of quasi-compact
quasi-separated schemes (\cref{lem:pf:steps2:limits}; the values of $\Z/q(i)$ on affine schemes
commute with filtered colimits of rings, \lref{GlobalTheory.T.tateTwist_finitary}).

(b) \emph{Open stabilizers.} For $x\in H^n(X_\Omega,\Zl(i))$ and $m\in\N$ the set of $\sigma$ with
$(\sigma^\ast x)_m=x_m$ is open. If $x_m$ comes from $X_L$, it is fixed by the open subgroup
$\Aut_L(\Omega)$, because $\sigma_X$ commutes with the projection $X_\Omega\to X_L$ for $\sigma$ fixing
$L$; the set in question is a union of cosets of this subgroup.

(c) \emph{Finite generation.} Since $X_\Omega$ is proper over the algebraically closed field $\Omega$,
the groups $H^n(X_\Omega,\Z/\ell^m(i))$ are finite (\lref{UpperTheory.Et.e4}, the finiteness of
\'etale cohomology with finite coefficients of proper schemes over algebraically closed fields in
which the coefficients are invertible); by \cref{lem:pf:steps2:adicfin} (which uses the Bockstein
sequences of \cref{lem:pf:steps2:bock}) the inverse limit $H^n(X_\Omega,\Zl(i))$ is a finitely
generated $\Zl$-module, and $H^n(X_\Omega,\Ql(i))$ is finite-dimensional.

(d) \emph{Continuity in each degree.} Apply \cref{lem:pf:steps2:abscont}(b) to the finitely generated
$\Zl$-module $M=H^n(X_\Omega,\Zl(i))$ (by (c)) with $K_m=\ker(M\to H^n(X_\Omega,\Z/\ell^m(i)))$, whose
intersection is $0$ and for which the sets in question are open by (b); a linear
form on $\Ql\otimes M$ restricts to a $\Zl$-linear map on $M$, and continuity extends to $\Ql\otimes
M$ by linearity.

(e) Write $v=\sum_nv_n$ with finitely many nonzero components; then
$\mu(\sigma^\ast v)=\sum_n\mu(\sigma^\ast v_n)$ is continuous by (d).
\end{proof}

\subsubsection{Proof of (S1)}

\begin{theorem}[Existence]\label{thm:pf:steps2:repex}
Let $E$ be a number field, $\Omega$ an algebraic closure of $E$, $\ell$ a prime, $H$ a nice Hecke tower
over $E$ and $\theta\colon\Gamma\to\Qlbar$. There are $dd\in\N$ and a continuous semisimple
homomorphism $\rho\colon G_E\to\GL_{dd}(\Qlbar)$ attached to $(H,\theta)$.
\formallabel{UpperTheory.Gal.Int.Geo.repEx_holds}\formallabel{UpperTheory.Gal.Int.Geo.repEx_core}
\end{theorem}

\begin{proof}
The field $\Omega$ is algebraically closed of characteristic $0$, and $\ell\ne0$ in $E$ and in
$\Omega$. The base change $Z_\Omega$ of the smooth projective $E$-scheme $Z$ is smooth and projective,
so $\stepsiiH(Z_\Omega)$ is finite-dimensional by \cref{lem:pf:steps2:fin}, and so is
$\stepsiiPiece_\theta\subseteq\Qlbar\otimes\stepsiiH(Z_\Omega)$. The group $G_E$ acts on
$\Qlbar\otimes\stepsiiH(Z_\Omega)$ by $\sigma\mapsto\sigma^\ast\otimes1$, a homomorphism
(\cref{def:pf:steps2:coh}), and $\stepsiiPiece_\theta$ is stable (\cref{lem:pf:steps2:hopgal}).
Choose a basis $b$ of $\stepsiiPiece_\theta$; \cref{lem:pf:steps2:matex} gives a matrix $\rho_0$ of
the piece and linear forms $\mu_{ij}$ with $\rho_0(\sigma)_{ij}=\mu_{ij}(\sigma^\ast b_j)$. By
\cref{lem:pf:steps2:grex} there are $P$ and a semisimple graded part $\rho$ of $P^{-1}\rho_0P$, and
$P^{-1}\rho_0P$ is again a matrix of the piece (\cref{lem:pf:steps2:matex}); hence $\rho$ is attached
to $(H,\theta)$.

Continuity: $q\colon Z\to\Spec E$ is proper (projective morphisms are proper), so by
\cref{lem:pf:steps2:cont} every function $\sigma\mapsto\nu(\sigma^\ast v)$ ($v\in\stepsiiH(Z_\Omega)$,
$\nu$ a $\Ql$-linear form) is continuous on $G_E$. Since $\Ql\to\Qlbar$ is continuous, the same holds
for $\sigma\mapsto\mu(\sigma^\ast v)$ with $v\in\Qlbar\otimes\stepsiiH(Z_\Omega)$ and $\mu$ a
$\Qlbar$-linear form (expand $v$ and $\mu$ in a $\Ql$-basis of $\stepsiiH(Z_\Omega)$). Hence the
entries of $\rho_0$ are continuous; the entries of $\rho$ are $0$ or entries of $P^{-1}\rho_0P$, which
are finite sums of products of constants and entries of $\rho_0$, hence continuous. A homomorphism to
$\GL_{dd}(\Qlbar)$ whose matrix entries are continuous is continuous, since the entries of
$\rho(\sigma)^{-1}=\rho(\sigma^{-1})$ are then continuous as well.
\end{proof}

\begin{remark}\label{rem:pf:steps2:corehyp}
The formal versions of \cref{lem:pf:steps2:fin,lem:pf:steps2:cont}, and the second formal theorem
tagged in \cref{thm:pf:steps2:repex}, take the package of \cref{prop:pf:steps2:etfound}(k) as a
hypothesis; the first theorem tagged in \cref{thm:pf:steps2:repex} supplies it.
\end{remark}

% ===========================================================================
\subsection{The Eichler--Shimura relation away from \texorpdfstring{$\ell$}{l}: statement (S2)}\label{sec:pf:steps2:es}

The proof of (S2) (\cref{thm:pf:steps2:esstep}) reduces, by linear algebra on the eigen-piece, to the
Eichler--Shimura relation for covers (\cref{prop:pf:steps2:towes}), which concerns a cover over a
discrete valuation ring carrying a congruence datum. This is reduced, by descent of the piece to the
generic fibre of the model, to the Eichler--Shimura relation for models
(\cref{prop:pf:steps2:esgen}) on the cohomology of a smooth proper scheme over a discrete valuation
ring with $p$-correspondences satisfying the pointwise congruence relation. The proof of the latter
occupies \cref{sec:pf:steps2:esgen}.

\subsubsection{Weighted traces and correspondences}

\begin{definition}[Weights, correspondences, cycles]\label{def:pf:steps2:weights}
A \emph{geometric point} of a scheme $X$ is a morphism $x\colon\Spec\Omega\to X$ with $\Omega$ an
algebraically closed field; $\mathrm{pt}(x)\in X$ denotes its image. A \emph{weight} on $Z$ is a
function $w$ assigning an integer $w(z)$ to every geometric point $z$ of $Z$. Let $p\colon Z\to S$ be
finite. A weight $w$ is a \emph{$p$-weight} if for every cartesian square formed by $u\colon U\to S$,
$r\colon Z_U\to U$ and $v\colon Z_U\to Z$, and every locally constant $\varphi\colon Z_U\to\Z$, there is
a locally constant $\psi\colon U\to\Z$ with $\psi(\mathrm{pt}(x))=\sum_{z:\,r\circ z=x}w(v\circ z)\,
\varphi(\mathrm{pt}(z))$ for every geometric point $x$ of $U$. A \emph{(weighted finite)
correspondence} $c$ from $S$ to $Y$ consists of a scheme $Z_c$, a finite morphism $p_c\colon Z_c\to
S$, a morphism $f_c\colon Z_c\to Y$ and a $p_c$-weight $w_c$; its \emph{cycle} is
\[
  \stepsiicyc_c(x,y)=\sum_{z:\ p_c\circ z=x,\ f_c\circ z=y}w_c(z)
\]
for geometric points $x$ of $S$ and $y$ of $Y$ with the same field. When $\ell$ is invertible on
$S$, the development defines trace maps $\stepsiitr_{p,w}\colon H^n(Z,\Ql(i))\to H^n(S,\Ql(i))$ for
$p$-weights $w$ (\cref{prop:pf:steps2:traces}), and the \emph{operator} of $c$ is
$\stepsiiop_c=\stepsiitr_{p_c,w_c}\circ f_c^\ast\colon H^n(Y,\Ql(i))\to H^n(S,\Ql(i))$.
\formal{UpperTheory.Gal.TranscendenceTraces.IsWt}\formal{UpperTheory.Gal.TranscendenceTraces.Corr.opQ}
\end{definition}

\begin{proposition}[Trace formalism]\label{prop:pf:steps2:traces}
Let $p\colon Z\to S$ be finite. In \textup{(b)--(g)} assume that $\ell$ (respectively $q$) is
invertible on $S$.
\begin{enumerate}[label=\textup{(\alph*)},nosep]
\item $p$-weights are stable under base change ($w\circ v$ is an $r$-weight), sums and integer
  multiples; the constant weight $1$ is a $p$-weight if $p$ is finite \'etale; a $p$-weight vanishing
  at all geometric points over a dense subset of $S$ vanishes; and if $S$ is locally normal (covered
  by open affines whose rings are integrally closed domains) there is a $p$-weight $w^{\mathrm{nor}}_p$
  which at every geometric point $z$ with $\mathrm{pt}(p\circ z)$ a maximal point of $S$ equals the
  length of the local ring of the fibre $Z\times_S\Spec\Omega$ at $z$.
\item (Sums.) For a cartesian square as in \cref{def:pf:steps2:weights} and a section $a$ of
  $\Z/q(i)$ over $Z_U$, the trace of $a$ is the section of $\Z/q(i)$ over $U$ whose pull-back to every
  geometric point $x$ of $U$ is $\sum_{r\circ z=x}w(v\circ z)\,z^\ast a$; this property characterizes
  the trace maps.
\item (Base change.) For a cartesian square $(v,p',p,g)$,
  $g^\ast\circ\stepsiitr_{p,w}=\stepsiitr_{p',w\circ v}\circ v^\ast$; for correspondences $c,c'$ with a
  cartesian square $(v,p_{c'},p_c,g)$, $f_c\circ v=g'\circ f_{c'}$ and $w_{c'}=w_c\circ v$, one has
  $g^\ast\circ\stepsiiop_c=\stepsiiop_{c'}\circ g'^\ast$.
\item (Composition.) For finite $r\colon W\to Z$ with an $r$-weight $w'$,
  $\stepsiitr_{p\circ r,\,w'\cdot(w\circ r)}=\stepsiitr_{p,w}\circ\stepsiitr_{r,w'}$.
\item (Degree.) If $\sum_{p\circ z=x}w(z)=d$ for every geometric point $x$ of $S$, then
  $\stepsiitr_{p,w}\circ p^\ast=d$.
\item (Splitting.) If $p$ is finite \'etale and $s_1,\dots,s_m\colon U\to Z$ with $p\circ s_j=x$ split
  $p$ over $x\colon U\to S$, then $x^\ast\circ\stepsiitr_{p,1}=\sum_js_j^\ast$.
\item (Cycles determine operators.) Let $c_1,\dots,c_r$ be correspondences from $S$ to $Y$, with
  morphisms $S\to B$ and $Y\to B$, the latter separated and locally of finite type, such that $f_{c_a}$
  and $p_{c_a}$ have the same composite to $B$; let $m_a\in\Z$ and let $D\subseteq S$ be dense. If
  $\sum_am_a\stepsiicyc_{c_a}(x,y)=0$ for all geometric points $x$ of $S$ with $\mathrm{pt}(x)\in D$
  and all $y$, then $\sum_am_a\stepsiiop_{c_a}=0$ on every $H^n(Y,\Ql(i))$.
\end{enumerate}
\formallabel{UpperTheory.Gal.TranscendenceTraces.all_holds}
\end{proposition}

This is a formalism of trace maps for finite morphisms with multiplicities. For flat quasi-finite
separated morphisms of finite presentation, and more generally for weighted quasi-finite separated
morphisms, trace morphisms of sheaves $f_!f^\ast F\to F$ that are compatible with base change and
composition and normalized by a degree condition are constructed in
\cite[Exp.~XVII, 6.2.3, 6.2.5]{ArtinGrothendieckVerdier1972}. There a weighting of $f\colon X\to S$ is an
integer function $n$ on the geometric points of $X$, depending only on their images, such that for
every \'etale $U\to S$ and every open and closed part $X_1$ of $X\times_SU$ finite over $U$ the function
$s\mapsto\sum_{x\in X_1(s)}n(x)$ is locally constant; by \cref{prop:pf:steps2:wts}(b) below, $p$-weights
are weightings in this sense. The development constructs the formalism from scratch; the proof is
given in \cref{sec:pf:steps2:trproof}. Statement (g) is the form in which the congruence relation
enters.
\begin{lemma}\label{lem:pf:steps2:corropeq}
Let $c$ be a correspondence from $S$ to $S$ with $p_c$ finite \'etale and $w_c=1$, and let $\ell$ be
invertible on $S$. If $T$ is an operator of the correspondence $(p_c,f_c)$
(\cref{def:pf:steps2:hop}), then $T=\bigoplus_n\stepsiiop_c$ on $\stepsiiH(S)$.
\formallabel{UpperTheory.Gal.Mod.corrOpEq}
\end{lemma}

\begin{proof}
Let $c'\colon W\to S$ be a Galois cover with group $G$ and $(\sigma_j)$ a family splitting $p_c$ over
$c'$ with $c'^\ast T(x)=\sum_j(f_c\sigma_j)^\ast x$. By \cref{prop:pf:steps2:traces}(f),
$c'^\ast\circ\stepsiitr_{p_c,1}=\sum_j\sigma_j^\ast$, so $c'^\ast\stepsiiop_c(x)=\sum_j\sigma_j^\ast
f_c^\ast x=c'^\ast T(x)$. Finally $c'^\ast$ is injective on $H^n(S,\Ql)$: by (f), applied to the
splitting $(a_g)_{g\in G}$ of $c'$ over itself, and by (e), $c'^\ast\stepsiitr_{c',1}=\sum_ga_g^\ast$
and $\stepsiitr_{c',1}c'^\ast=|G|$, so $c'^\ast$ has the left inverse $|G|^{-1}\stepsiitr_{c',1}$.
\end{proof}

\begin{lemma}\label{lem:pf:steps2:hopop}
Let $H$ be a Hecke tower over $k$, $\Omega/k$ an extension in which $\ell$ is invertible, and
$g\in\Gamma$ with $s_g$ finite \'etale of degree $d$. Then
$T_g=\bigoplus_n\stepsiitr_{s_{g,\Omega},1}\circ t_{g,\Omega}^\ast$.
\formallabel{UpperTheory.Gal.Mod.hopOp_holds}
\end{lemma}

\begin{proof}
$s_{g,\Omega}$ is finite \'etale of degree $d$ (base change), hence split over a finite \'etale
Galois cover $c\colon W\to Z_\Omega$ by $d$ morphisms $\sigma_j$ (\cref{lem:pf:steps2:galclos}(b)). By \cref{prop:pf:steps2:traces}(f),
$c^\ast\circ\stepsiitr_{s_{g,\Omega},1}\circ t^\ast_{g,\Omega}=\sum_j(t_{g,\Omega}\sigma_j)^\ast$, so
$T'=\bigoplus_n\stepsiitr_{s_{g,\Omega},1}t_{g,\Omega}^\ast$ is an operator of
$(s_{g,\Omega},t_{g,\Omega})$. Hence an operator exists, $T_g$ is one of them, and both $T_g$ and $T'$
equal $\bigoplus_n\stepsiiop_c$ for $c=(s_{g,\Omega},t_{g,\Omega},1)$ by \cref{lem:pf:steps2:corropeq}.
\end{proof}

\begin{lemma}\label{lem:pf:steps2:opcover}
Let $c,c'$ be correspondences from $S$ to $S'$ with $w_c=w_{c'}=1$, and $u\colon Z_c\to Z_{c'}$ finite
with $p_{c'}u=p_c$, $f_{c'}u=f_c$, such that every geometric point of $Z_{c'}$ has exactly $\delta$
lifts along $u$. If $\ell$ is invertible on $S$, then $\stepsiiop_c=\delta\cdot\stepsiiop_{c'}$.
\formallabel{UpperTheory.Gal.Mod.opCover}
\end{lemma}

\begin{proof}
By \cref{prop:pf:steps2:traces}(d), $\stepsiitr_{p_c,1}=\stepsiitr_{p_{c'},1}\circ\stepsiitr_{u,1}$, and
by (e), $\stepsiitr_{u,1}\circ u^\ast=\delta$. Hence
$\stepsiiop_c=\stepsiitr_{p_{c'},1}\stepsiitr_{u,1}u^\ast f_{c'}^\ast=\delta\,\stepsiiop_{c'}$.
\end{proof}

\subsubsection{Descent to the model and correspondences on the model}

Let $C$ be a cover of $(H,\theta)$ over $O\to k$ and $\Omega/k$ an extension. Write
$\pi_\Omega\colon Z_\Omega\to(Y_k)_\Omega$ for the base change of $\pi_k\colon Z\to Y_k$; it commutes
with the Galois actions, $\pi_\Omega^\ast\sigma^\ast=\sigma^\ast\pi_\Omega^\ast$ for
$\sigma\in\Aut_k(\Omega)$, as $\pi_\Omega\circ\sigma_Z=\sigma_{Y_k}\circ\pi_\Omega$. For a
$p$-correspondence $T$ on $Y$, a correspondence $c$ from $(Y_k)_\Omega$ to itself is a \emph{base
change of $T$ with weight $1$} if $w_c=1$ and there is $v\colon Z_c\to T_Z$ such that $(v,p_c)$ is a
pull-back of $\alpha_T$ along the projection $(Y_k)_\Omega\to Y$, $\beta_T\circ v$ equals the
projection composed with $f_c$, and $p_c$, $f_c$ have the same composite to $\Spec\Omega$.

\begin{lemma}\label{lem:pf:steps2:bcex}
If $p\ne0$ in $k$ and $T$ is a $p$-correspondence on $Y$ with $u\beta_T=u\alpha_T$, then $T$ has a base
change with weight $1$ to $(Y_k)_\Omega$.
\formallabel{UpperTheory.Gal.Mod.bcEx}
\end{lemma}

\begin{proof}
Let $Z_c=T_Z\times_Y(Y_k)_\Omega$ with its projections; $p_c$ is finite, and \'etale because $p$ is
invertible on $(Y_k)_\Omega$ and $\alpha_T$ is \'etale over $Y[1/p]$, so $1$ is a $p_c$-weight
(\cref{prop:pf:steps2:traces}(a)). The morphism $f_c$ is induced by $\beta_T$ and the projection to
$\Spec\Omega$, which is compatible because $u\beta_T=u\alpha_T$.
\end{proof}

\begin{lemma}\label{lem:pf:steps2:fixdesc}
Assume $\ell$ invertible in $\Omega$ and $L\supseteq\Ql$. Every $x\in\stepsiiPiece_\theta(H,\Omega,L)$
is of the form $(\pi_\Omega^\ast\otimes L)(y)$ with $y\in L\otimes\stepsiiH((Y_k)_\Omega)$.
\formallabel{UpperTheory.Gal.Mod.fixDesc}
\end{lemma}

\begin{proof}
The automorphisms $a_{\delta,\Omega}$ of $Z_\Omega$ split $\pi_\Omega$ over itself (base change of the
splitting of $\pi_k$), and $\pi_\Omega$ is finite \'etale. For $\delta\in\Delta$ choose $g$ with $s_g$
an isomorphism, $t_g=a_\delta s_g$ and $\theta(g)=1$. By \cref{lem:pf:steps2:hopop},
$T_g=\stepsiitr_{s_{g,\Omega},1}t_{g,\Omega}^\ast=\stepsiitr_{s_{g,\Omega},1}s_{g,\Omega}^\ast
a_{\delta,\Omega}^\ast$, and $\stepsiitr_{s_{g,\Omega},1}s_{g,\Omega}^\ast=\mathrm{id}$ for the
isomorphism $s_{g,\Omega}$ (\cref{prop:pf:steps2:traces}(e)); hence $T_g=a_{\delta,\Omega}^\ast$, and
$a_{\delta,\Omega}^\ast x=x$ because $x$ is an eigenvector of $T_g$ with eigenvalue $\theta(g)=1$. By
\cref{prop:pf:steps2:traces}(f), $\pi_\Omega^\ast\circ\stepsiitr_{\pi_\Omega,1}=\sum_\delta
a_{\delta,\Omega}^\ast$, so $y=|\Delta|^{-1}(\stepsiitr_{\pi_\Omega,1}\otimes L)(x)$ satisfies
$(\pi_\Omega^\ast\otimes L)(y)=x$ ($L$ has characteristic $0$).
\end{proof}

\begin{lemma}\label{lem:pf:steps2:liesop}
Assume $\ell$ invertible in $\Omega$. Let $T$ be a $p$-correspondence on $Y$ and $g\in\Gamma$ with
$s_g$ finite \'etale of some degree, such that $(\alpha_T,\beta_T)$ lies under $g$ with multiplicity
$m$, and let $c$ be a base change of $T$ with weight $1$ to $(Y_k)_\Omega$. Then
$T_g\circ\pi_\Omega^\ast=m\cdot\pi_\Omega^\ast\circ\bigoplus_n\stepsiiop_c$.
\formallabel{UpperTheory.Gal.Mod.liesOp}
\end{lemma}

\begin{proof}
Let $c_1$ be the pull-back of $c$ along $\pi_\Omega$: $Z_{c_1}=Z_c\times_{(Y_k)_\Omega}Z_\Omega$ with
$p_{c_1}$ the second projection, $f_{c_1}=f_c\circ\mathrm{pr}_1$ and weight $1$ (a $p_{c_1}$-weight by
\cref{prop:pf:steps2:traces}(a)). Let $c_2=(s_{g,\Omega},\pi_\Omega t_{g,\Omega},1)$. Pasting the
cartesian squares defining $c$ and $c_1$ identifies $Z_{c_1}$ with
$(Z\times_{\pi,Y,\alpha_T}T_Z)_\Omega$, and the morphism $\kappa$ of \cref{def:pf:steps2:cover}
induces $u\colon Y_{g,\Omega}\to Z_{c_1}$ with $p_{c_1}u=s_{g,\Omega}$ and $f_{c_1}u=\pi_\Omega
t_{g,\Omega}$ (the latter because $\beta_T\kappa=\pi t_g$); $u$ is the base change of $Y_g\to
Z\times_YT_Z$, which is finite \'etale of degree $m$. Hence $\stepsiiop_{c_2}=m\,\stepsiiop_{c_1}$ by
\cref{lem:pf:steps2:opcover}, and $\stepsiiop_{c_1}=\pi_\Omega^\ast\circ\stepsiiop_c$ by
\cref{prop:pf:steps2:traces}(c). By \cref{lem:pf:steps2:hopop},
$T_g\pi_\Omega^\ast=\stepsiitr_{s_{g,\Omega},1}t_{g,\Omega}^\ast\pi_\Omega^\ast=\stepsiiop_{c_2}=
m\,\pi_\Omega^\ast\stepsiiop_c$.
\end{proof}

\subsubsection{The Eichler--Shimura relations for models and for covers}

For $q,j\in\N$ put $\stepsiico(q,j)=(-1)^jq^{\binom j2}$.

\begin{proposition}[Eichler--Shimura relation for models]\label{prop:pf:steps2:esgen}
Let $\ell$ and $p$ be primes and $f,n\in\N$. Let $O$ be a discrete valuation ring with uniformizer
$\varpi$ in which $p$ is not a unit; $K$ and $\kappa$ fields with ring maps $O\to K$, $O\to\kappa$
such that $p\ne0$ in $K$, $\varpi\mapsto0$ in $\kappa$, and $a^{p^f}=a$ for all $a\in\kappa$; $\Omega$ an
algebraically closed extension of $K$ and $V\subseteq\Omega$ a valuation ring containing the image of
$O$; $\bar k$ an algebraically closed extension of $\kappa$ in which $\ell$ is invertible, with a
surjective ring map $r\colon V\to\bar k$ compatible with $O\to\kappa\to\bar k$. Let $u\colon Y\to\Spec
O$ be smooth and proper, $T_0,\dots,T_n$ $p$-correspondences on $Y$ with $u\beta_{T_j}=u\alpha_{T_j}$
and $\stepsiiCR_{p^f,u^\ast\varpi}(T)$, and $c_j$ base changes of $T_j$ with weight $1$ to
$Y_\Omega:=(Y_K)_\Omega$. Let $\sigma\in\Aut_K(\Omega)$ with $\sigma(V)\subseteq V$, and $k\in\Z$.
Then on $H^k(Y_\Omega,\Ql)$:
\begin{enumerate}[label=\textup{(\alph*)},nosep]
\item if $r(\sigma(v))=r(v)$ for all $v\in V$, then $\sigma^\ast=\mathrm{id}$;
\item if $r(\sigma(v))=r(v)^{p^f}$ for all $v\in V$, then
  $\sum_{j=0}^n\stepsiico(p^f,j)\,\stepsiiop_{c_j}\circ(\sigma^{-1})^{\ast(n-j)}=0$.
\end{enumerate}
\formallabel{UpperTheory.Gal.Mod.esGen_all}\formallabel{UpperTheory.Gal.Mod.esGen}
\end{proposition}

The proof is given in \cref{sec:pf:steps2:esgen}.

\begin{proposition}[Eichler--Shimura relation for covers]\label{prop:pf:steps2:towes}
Let $\ell,p$ be primes and $f,n\in\N$. Let $O$ be a discrete valuation ring in which $p$ is not a
unit and whose residue field satisfies $a^{p^f}=a$; $K$ a field with a ring map $O\to K$ and $p\ne0$
in $K$; $\Omega$ an algebraically closed extension of $K$ and $V\subseteq\Omega$ a valuation ring
containing the image of $O$, such that the maximal ideal of $O$ maps into the maximal ideal of $V$
and $\ell$ is a unit of $V$. Let $L\supseteq\Ql$ be a field, $H$ a nice Hecke tower over $K$,
$\theta\colon\Gamma\to L$, $C$ a cover of $(H,\theta)$ over $O$ carrying a congruence datum
$(p,f,n,\vartheta)$, and assume that $\stepsiiPiece_\theta=\stepsiiPiece_\theta(H,\Omega,L)$ is stable
under $\Aut_K(\Omega)$. Let $\sigma\in\Aut_K(\Omega)$ with $\sigma(V)\subseteq V$ and
$x\in\stepsiiPiece_\theta$.
\begin{enumerate}[label=\textup{(\alph*)},nosep]
\item If $v_V(\sigma(v)-v)<1$ for all $v\in V$, then $\sigma^\ast x=x$.
\item If $v_V(\sigma(v)-v^{p^f})<1$ for all $v\in V$, then
  $\sum_{j=0}^n\stepsiico(p^f,j)\vartheta_j\,(\sigma^{-1})^{\ast(n-j)}x=0$.
\end{enumerate}
\formallabel{UpperTheory.Gal.Mod.towES_holds}\formallabel{UpperTheory.Gal.Mod.towES}
\end{proposition}

\begin{proof}
Let $(\varpi,T,\gamma,m)$ witness the congruence datum. Since $\ell$ is a unit of $V$, it is
invertible in $\Omega$. By \cref{lem:pf:steps2:fixdesc}, $x=\pi_\Omega^\ast y$ (we omit $\otimes L$
from the notation). By \cref{lem:pf:steps2:bcex} choose base changes $c_j$ of $T_j$ with weight $1$ to
$Y_\Omega$. The model $u$ is smooth and proper (projective). The map $O\to V$ is local, so it
induces a map of residue fields $\kappa_O\to\kappa_V$; the residue field $\kappa_V$ is algebraically
closed, as $V$ is a valuation ring with algebraically closed fraction field, and $\ell$ is invertible
in it. Thus \cref{prop:pf:steps2:esgen} applies with $\kappa=\kappa_O$, $\bar k=\kappa_V$ and $r$ the
residue map; its hypothesis $\varpi\mapsto0$ in $\kappa_O$ holds as $\varpi$ is not a unit. The
conditions on $\sigma$ in (a), (b) translate into $r(\sigma v)=r(v)$, respectively
$r(\sigma v)=r(v)^{p^f}$.

(a) By \cref{prop:pf:steps2:esgen}(a), $\sigma^\ast=\mathrm{id}$ on $\stepsiiH(Y_\Omega)$; hence
$\sigma^\ast x=\sigma^\ast\pi_\Omega^\ast y=\pi_\Omega^\ast\sigma^\ast y=\pi_\Omega^\ast y=x$.

(b) Put $G=(\sigma^{-1})^\ast$ on $L\otimes\stepsiiH(Z_\Omega)$, $G'=(\sigma^{-1})^\ast$ on
$L\otimes\stepsiiH(Y_\Omega)$, $P=\pi_\Omega^\ast$, $A_j=T_{\gamma_j}$, $B_j=\bigoplus_n\stepsiiop_{c_j}$
and $c'_j=\stepsiico(p^f,j)$. Then $PG'=GP$; $A_jP=m_jPB_j$ with $m_j\ne0$
(\cref{lem:pf:steps2:liesop}; every $s_g$ is finite \'etale of some degree since $H$ is nice);
$\sum_jc'_jB_jG'^{\,n-j}=0$ (\cref{prop:pf:steps2:esgen}(b), in each degree); $G$ preserves
$\stepsiiPiece_\theta$; and $A_j$ acts on $\stepsiiPiece_\theta$ by $\theta(\gamma_j)=m_j\vartheta_j$.
Since $G^iPy\in\stepsiiPiece_\theta$ for all $i$,
\[
  m_j\vartheta_j\,G^{n-j}Py=A_jG^{n-j}Py=A_jPG'^{\,n-j}y=m_jPB_jG'^{\,n-j}y,
\]
so $\vartheta_jG^{n-j}Py=PB_jG'^{\,n-j}y$ ($L$ has characteristic $0$), and summing with the
coefficients $c'_j$ gives $\sum_jc'_j\vartheta_jG^{n-j}x=P\bigl(\sum_jc'_jB_jG'^{\,n-j}y\bigr)=0$.
\end{proof}

\begin{lemma}\label{lem:pf:steps2:placedata}
Let $w$ be a finite place of $E$ and $A\subset\Omega$ a valuation ring over $w$. Then $\stepsiiOw$
maps into $A$ and its maximal ideal into the maximal ideal of $A$; the residue field of $\stepsiiOw$
satisfies $a^{\stepsiiNw}=a$; an integer $m\notin w$ is a unit of $A$, and an integer $m\in w$ is not a
unit of $\stepsiiOw$; every $s\in I_A$ satisfies $s(A)\subseteq A$ and $v_A(s(v)-v)<1$ for $v\in A$; and
every arithmetic Frobenius element $\sigma$ for $(A,q)$ satisfies $\sigma(A)\subseteq A$ and
$v_A(\sigma(v)-v^q)<1$ for $v\in A$.
\formallabel{UpperTheory.Gal.Mod.placeData}
\end{lemma}

\begin{proof}
These are direct consequences of the definitions: $A\cap E=\stepsiiOw$ because $A$ lies over $w$; the
residue field $\OO_E/w$ of $\stepsiiOw$ is finite with $\stepsiiNw$ elements; $m\notin w$ means
$v_A(m)=1$; elements of $I_A$ and Frobenius elements lie in $D_A$, the stabilizer of $A$, and the
congruences are their definitions.
\end{proof}

\begin{lemma}\label{lem:pf:steps2:grpass}
If $\rho$ is a graded part of $\rho_0$ and $Q$ is a polynomial with $Q(\rho_0(g))=0$, then
$Q(\rho(g))=0$.
\formallabel{UpperTheory.Gal.Mod.grPass}
\end{lemma}

\begin{proof}
$\rho_0(g)$ is block upper triangular for some $c$ and $\rho(g)$ is its block-diagonal part. The
block-diagonal part is multiplicative and additive on block upper triangular matrices, so the
block-diagonal part of $Q(\rho_0(g))=0$ is $Q(\rho(g))$.
\end{proof}

\begin{theorem}[Eichler--Shimura relation away from $\ell$]\label{thm:pf:steps2:esstep}
Let $E$ be a number field, $\Omega$ an algebraic closure, $\ell$ a prime, $\iota\colon\Qlbar\to\C$ a
field isomorphism, $H$ a nice Hecke tower over $E$, $\theta\colon\Gamma\to\Qlbar$, and
$\rho\colon G_E\to\GL_{dd}(\Qlbar)$ continuous and attached to $(H,\theta)$. Let $w$ be a finite place
with $\ell\notin w$, $n\in\N$, $\vartheta\colon\N\to\Qlbar$ and $R$ a finite multiset in $\C$. If
$\stepsiiCong_w(H,\theta;n,\vartheta)$ holds and $R$ is an $\iota$-root multiset for $(w,n,\vartheta)$,
then $\stepsiiES_w(\rho,R)$: for every place $A$ of $\Omega$ over $w$, $\rho(I_A)=1$ and
$(\iota^{-1}\stepsiirt_R)(\rho(\sigma^{-1}))=0$ for every arithmetic Frobenius element $\sigma$ for
$(A,\stepsiiNw)$.
\formallabel{UpperTheory.Gal.Mod.esStep_holds}\formallabel{UpperTheory.Gal.Mod.esStep}
\end{theorem}

\begin{proof}
Let $\rho_0$ be a matrix of the piece, with basis $b$, of which $\rho$ is a graded part, and let $C$,
$p$, $f$ witness $\stepsiiCong_w(H,\theta;n,\vartheta)$: $C$ is a cover over $\stepsiiOw$, $p\in w$ is
prime, $\stepsiiNw=p^f$, and $C$ carries the congruence datum $(p,f,n,\vartheta)$. The piece is
$G_E$-stable, since $\sigma^\ast b_j=\sum_i\rho_0(\sigma)_{ij}b_i$. Fix a place $A$ over $w$. By
\cref{lem:pf:steps2:placedata}, \cref{prop:pf:steps2:towes} applies with $O=\stepsiiOw$, $K=E$, $V=A$
and $L=\Qlbar$: $p$ is not a unit of $\stepsiiOw$, $p\ne0$ in $E$, the residue field satisfies
$a^{p^f}=a$, the maximal ideal maps into that of $A$, and $\ell$ is a unit of $A$ as $\ell\notin w$.
For an endomorphism $\phi$ of the piece with matrix $M$ in the basis $b$ and a polynomial $Q$, the
matrix of $Q(\phi)$ is $Q(M)$; so polynomial identities for $\sigma^\ast$ on the piece are identities
for $\rho_0(\sigma)$.

If $s\in I_A$, then $s(A)\subseteq A$ and $v_A(s(v)-v)<1$, so $s^\ast$ is the identity on the piece by
\cref{prop:pf:steps2:towes}(a). Thus $\rho_0(s)-1=0$, and $\rho(s)=1$ by \cref{lem:pf:steps2:grpass}.

If $\sigma$ is an arithmetic Frobenius element for $(A,\stepsiiNw)$, then
\cref{prop:pf:steps2:towes}(b) gives, on the piece,
\[
  \sum_{j=0}^n\stepsiico(\stepsiiNw,j)\vartheta_j(\sigma^{-1})^{\ast(n-j)}=0,
\]
i.e.\ $\stepsiiHk_{\stepsiiNw,n,\vartheta}(\rho_0(\sigma^{-1}))=0$, and hence
$\stepsiiHk_{\stepsiiNw,n,\vartheta}(\rho(\sigma^{-1}))=0$ by \cref{lem:pf:steps2:grpass}. As $R$ is
an $\iota$-root multiset, $\stepsiiHk_{\stepsiiNw,n,\vartheta}=\vartheta_0\cdot\iota^{-1}\stepsiirt_R$
with $\vartheta_0\ne0$, so $(\iota^{-1}\stepsiirt_R)(\rho(\sigma^{-1}))=0$.
\end{proof}

\begin{remark}\label{rem:pf:steps2:esunused}
In this proof the hypothesis that the cover in $\stepsiiCong_w(H,\theta;n,\vartheta)$ is onto is not
used, nor is the continuity of $\rho$. The conclusion concerns $\rho(\sigma^{-1})$, i.e.\ geometric
Frobenius elements, and asserts annihilation by $\iota^{-1}\stepsiirt_R$.
\end{remark}


\subsubsection{Proof of the Eichler--Shimura relation for models}\label{sec:pf:steps2:esgen}

\paragraph{Specialization.}
A local ring $A$ is \emph{strictly henselian} if it is henselian with separably closed residue field,
and \emph{absolutely integrally closed} if every monic polynomial of positive degree over $A$ has a
root in $A$.

\begin{proposition}[Specialization]\label{prop:pf:steps2:specgood}
Let $O$ be a ring, $K$, $\kappa$ fields with ring maps from $O$, $\Omega$ an algebraically closed
extension of $K$, $V\subseteq\Omega$ a valuation ring containing the image of $O$, $\bar k$ an
algebraically closed extension of $\kappa$ and $r\colon V\to\bar k$ a surjective ring map compatible
with $O\to\kappa\to\bar k$. Let $u\colon Y\to\Spec O$ be smooth and proper, and let
$i_Y\colon Y_{\bar k}:=(Y_\kappa)_{\bar k}\to Y_V$ and $j_Y\colon Y_\Omega:=(Y_K)_\Omega\to Y_V$ be the
natural morphisms. Then:
\begin{enumerate}[label=\textup{(\alph*)},nosep]
\item for every $q$ invertible in $\bar k$ and every $i$, $i_Y^\ast$ and $j_Y^\ast$ are bijective on
  $H^n(-,\Z/q(i))$ for all $n$;
\item for every prime $\ell$ invertible in $\bar k$ and all $n,i$ there is a $\Ql$-linear isomorphism
  $\stepsiisp\colon H^n(Y_{\bar k},\Ql(i))\to H^n(Y_\Omega,\Ql(i))$ with
  $\stepsiisp\circ i_Y^\ast=j_Y^\ast$, such that $\stepsiisp\circ\tau^\ast=\sigma^\ast\circ\stepsiisp$
  whenever $\sigma\in\Aut_K(\Omega)$, $\sigma(V)\subseteq V$, $\tau\in\Aut_\kappa(\bar k)$ and
  $r\circ\sigma|_V=\tau\circ r$.
\end{enumerate}
\formallabel{UpperTheory.Gal.Mod.specGood}\formallabel{UpperTheory.Gal.Mod.spSpecIf}
\end{proposition}

\begin{proof}
The kernel of $r$ is the maximal ideal of $V$, so $\bar k$ is the residue field of $V$ and $i_Y$ is
the inclusion of the closed fibre of $Y_V$, while $j_Y$ is the inclusion of the generic fibre. The
valuation ring $V$ has algebraically closed fraction field $\Omega$; hence it is absolutely
integrally closed (roots in $\Omega$ of monic polynomials over $V$ are integral over $V$, hence lie
in $V$) and strictly henselian (it is henselian, and its residue field is algebraically closed).
Since $Y_V$ is proper over $V$, the proper base change theorem
(\lref{UpperTheory.Et.proper_base_change}, which states that for a strictly henselian local ring
$A$ with $q\in A^\times$ and $X$ proper over $A$, restriction to the closed fibre is bijective on
$H^\ast(-,\Z/q(i))$) gives (a) for $i_Y$; since $Y_V$ is smooth over $V$,
\cref{prop:pf:steps2:etfound}(b) gives (a) for $j_Y$; $q$ is invertible in $V$ because it is
invertible in $\bar k$.

For (b) define $\stepsiisp$ at finite level as $j_Y^\ast\circ(i_Y^\ast)^{-1}$; these maps are
compatible with reduction and pass to the inverse limit and to $\Ql$-coefficients, giving an
isomorphism with $\stepsiisp\circ i_Y^\ast=j_Y^\ast$. If $\sigma$ and $\tau$ are as stated, the ring
endomorphism $\sigma|_V$ induces $\sigma_V\colon Y_V\to Y_V$ with $\sigma_V\circ j_Y=j_Y\circ\sigma_Y$
and, because $r\circ\sigma|_V=\tau\circ r$, $\sigma_V\circ i_Y=i_Y\circ\tau_Y$. Hence
$\stepsiisp(\tau^\ast i_Y^\ast z)=\stepsiisp(i_Y^\ast\sigma_V^\ast z)=j_Y^\ast\sigma_V^\ast z=
\sigma^\ast j_Y^\ast z=\sigma^\ast\stepsiisp(i_Y^\ast z)$, and $i_Y^\ast$ is surjective.
\end{proof}

\paragraph{Frobenius.}
For a scheme $X$ with $p=0$ on $X$ and $f\in\N$, the absolute Frobenius $F_X\colon X\to X$ is the
identity on the underlying space and $a\mapsto a^{p^f}$ on the structure sheaf.

\begin{lemma}\label{lem:pf:steps2:frob}
\begin{enumerate}[label=\textup{(\alph*)},nosep]
\item If $p=0$ on $X$ and $m$ is invertible on $X$, then $F_X^\ast$ is the identity on
  $H^n(X,\Z/m(0))$ for all $n$; hence $F_X^\ast$ is the identity on $H^n(X,\Ql)$ if $\ell$ is
  invertible on $X$.
\item Let $\kappa\subseteq\bar k$ be fields with $a^{p^f}=a$ on $\kappa$, $X$ a $\kappa$-scheme with $p=0$
  on $X_{\bar k}$, and $\tau\in\Aut_\kappa(\bar k)$ with $\tau(a)=a^{p^f}$. There is a $\bar
  k$-morphism $\stepsiiFr\colon X_{\bar k}\to X_{\bar k}$ with $\tau_X\circ\stepsiiFr=F_{X_{\bar k}}$
  such that for every field $\Omega'$ with an endomorphism $\phi(a)=a^{p^f}$ and every
  $x\colon\Spec\Omega'\to X_{\bar k}$, the projections of $\stepsiiFr\circ x$ and $x\circ\Spec\phi$ to
  $X$ coincide.
\end{enumerate}
\formallabel{UpperTheory.Gal.Mod.frobTors}\formallabel{UpperTheory.Gal.Mod.frobQ}\formallabel{UpperTheory.Gal.Mod.frEx}
\end{lemma}

\begin{proof}
(a) This is the classical triviality of the absolute Frobenius on \'etale cohomology
\cite[Exp.~XV, \S1]{Grothendieck1977}. The formal proof uses the trace formalism. Absolute Frobenius
morphisms commute with all morphisms of schemes on which $p=0$, so the maps $F_U^\ast$ on the terms of
the resolution of $\Z/m(0)$, over the squares $(u,\mathrm{id},u)$ over $\mathrm{id}_X$, form a family
over $\mathrm{id}_X$ in the sense of \cref{sec:pf:steps2:trproof}, which on the self square is
$F_X^\ast$. It sums the constant weight $1$: for a section $a$ of $\Z/m(0)$ over $U$ and a geometric
point $x$ of $U$, $x^\ast F_U^\ast a=(F_U\circ x)^\ast a=(x\circ F_\Omega)^\ast a=F_\Omega^\ast x^\ast a=
x^\ast a$, because sections of $\Z/m(0)$ over $\Spec\Omega$ are fixed by $F_\Omega^\ast$ (they are
determined by their values at the single point, which $F_\Omega$ does not move; formally this is
checked using a primitive $m$-th root of unity). By \cref{lem:pf:steps2:trchar} and
\cref{prop:pf:steps2:trtors}(c), $F_X^\ast=\stepsiitr_{\mathrm{id},1}=\mathrm{id}$ on $H^n(X,\Z/m(0))$.
The statement for $\Ql$ follows levelwise. (b) $\stepsiiFr$ is the relative Frobenius
$F_X\times\mathrm{id}_{\bar k}$; then $\tau_X\circ\stepsiiFr=F_X\times\tau=F_{X_{\bar k}}$ as $\tau$ is
the $p^f$-power map, and the last property follows from the description of $F_X$ on points.
\end{proof}

\paragraph{Models over discrete valuation rings.}

\begin{lemma}\label{lem:pf:steps2:basefacts}
Let $O$ be a discrete valuation ring with uniformizer $\varpi$, $p\ne0$ a non-unit of $O$, and
$u\colon Y\to\Spec O$ smooth.
\begin{enumerate}[label=\textup{(\alph*)},nosep]
\item $Y$ is locally normal; every generic point of an irreducible component of $Y$ lies in
  $Y[1/p]$; every point of $Y$ specializes from a point of $Y[1/p]$; and $Y[1/p]\subseteq
  Y[1/\varpi]$.
\item For a field $\bar k$ with a ring map $O\to\bar k$ killing $\varpi$, the set of points of
  $Y_{\bar k}$ whose image in $Y$ is a generic point of $\{u^\ast\varpi=0\}$ is dense in $Y_{\bar k}$.
\item For every $p$-correspondence $T$ on $Y$, the normalized weight $w^{\mathrm{nor}}_{\alpha_T}$
  (\cref{prop:pf:steps2:traces}(a)) is an $\alpha_T$-weight equal to $1$ at every geometric point $z$
  of $T_Z$ at which $p\ne0$.
\end{enumerate}
\formallabel{UpperTheory.Gal.Mod.baseFacts}\formallabel{UpperTheory.Gal.Mod.denseGen}\formallabel{UpperTheory.Gal.Mod.nwtOne}
\end{lemma}

\begin{proof}
(a) Smooth over a normal ring, $Y$ is normal, and it is locally noetherian, hence locally normal.
$u$ is flat, so generizations lift: a generic point of $Y$ maps to the generic point of $\Spec O$,
where $p\ne0$, and every $y$ has a generization over the generic point of $\Spec O$. Finally
$p=e\varpi^s$ with $e$ a unit and $s\ge1$, so $p$ invertible at a point forces $\varpi$ invertible
there.
(b) The map $O\to\bar k$ factors through the residue field $O/\varpi$; the special fibre $Y_{O/\varpi}$
is smooth over a field, so the points of $Y_{\bar k}$ lying over generic points of $Y_{O/\varpi}$ are
dense, and the generic points of $Y_{O/\varpi}=\{u^\ast\varpi=0\}$ are the generic points of
$\{u^\ast\varpi=0\}$ in the sense of \cref{def:pf:steps2:ptid}.
(c) By (a) and \cref{prop:pf:steps2:traces}(a), $w^{\mathrm{nor}}_{\alpha_T}$ is an $\alpha_T$-weight which
at points over generic points of $Y$ is the multiplicity of the fibre. Over $Y[1/p]$ the morphism
$\alpha_T$ is finite \'etale; normalized weights are compatible with flat base change along maps
whose generic points are dense, and the multiplicity of an \'etale morphism is $1$. Since the generic
points of $Y$ lie in $Y[1/p]$, the restriction of $w^{\mathrm{nor}}_{\alpha_T}$ over $Y[1/p]$ is the
constant weight $1$.
\end{proof}

\paragraph{Pointwise cycles.}
Let $T=(T_0,\dots,T_n)$ be $p$-correspondences on $Y$ and $w_j$ weights on $(T_j)_Z$. For a geometric
point $x$ of $Y$ over a field $\Omega$, an endomorphism $\phi$ of $\Omega$ and $q\in\N$, we say that the
\emph{cycle identity holds for $(q,\phi,x)$} if for every geometric point $y$ of $Y$ over $\Omega$
\begin{equation}\label{eq:pf:steps2:cycid}
  \sum_{j=0}^n\stepsiico(q,j)\sum_{z}w_j(z)=0,
\end{equation}
the inner sum being over the geometric points $z$ of $(T_j)_Z$ over $\Omega$ with $\alpha_{T_j}\circ
z=x$ and $F_\phi^{\,n-j}(\beta_{T_j}\circ z)=y$, where $F_\phi$ is precomposition with $\Spec\phi$.

\begin{lemma}\label{lem:pf:steps2:wtval}
Let $T$ be a $p$-correspondence on $Y$ and $w$ an $\alpha_T$-weight equal to $1$ at geometric points
where $p\ne0$. Let $V$ be a valuation ring with algebraically closed fraction field $K_V$ and
algebraically closed residue field $\kappa_V$, with $p\ne0$ in $V$, and $y\colon\Spec V\to Y$. Then
$\stepsiiSec_T(y)$ is finite, and for every $z_0\colon\Spec\kappa_V\to T_Z$ lying over the closed point
of $y$, $w(z_0)$ equals the number of $s\in\stepsiiSec_T(y)$ whose reduction is $z_0$.
\formallabel{UpperTheory.Gal.Mod.wtVal_holds}\formallabel{UpperTheory.Gal.Mod.wtVal}
\end{lemma}

\begin{proof}
Let $r\colon W=T_Z\times_{\alpha_T,Y,y}\Spec V\to\Spec V$ be the base change, a finite morphism. By the
valuative criterion for the finite (hence proper) morphism $r$, restriction to the generic point is a
bijection from sections of $r$ to $K_V$-points of $W$ over the generic point, and the latter set is
finite. The scheme $W=\Spec A$ is affine with $A$ a finite $V$-algebra, and its points over the closed
point of $\Spec V$ are the maximal ideals of $A$, each with exactly one $\kappa_V$-point. By
\cref{lem:pf:steps2:hens} there is an idempotent $e\in A$ outside the maximal ideal of the point of
$z_0$ and inside all other maximal ideals; so $U=\{e\ne0\}$ is an open and closed subset of $W$
containing the point $z_0$ (viewed in $W$) and no other $\kappa_V$-point over the closed point. Apply
the definition of an $\alpha_T$-weight to the base change
square of $W$ and the locally constant function $\varphi=\mathbf 1_U$: there is a locally constant
$\psi$ on $\Spec V$ with $\psi(\mathrm{pt}(x))=\sum_{z\mapsto x}w(z)\mathbf 1_U(z)$. As $\Spec V$ is
connected, $\psi$ is constant. At the closed point the sum is $w(z_0)$. At the generic point all
weights are $1$ ($p\ne0$ in $K_V$), so the sum is the number of $K_V$-points of $W$ in $U$, i.e.\ the
number of sections $s$ of $r$ whose generic point lies in $U$. A section is a morphism from the
connected $\Spec V$, so its generic point lies in the open and closed set $U$ if and only if its
closed point does, i.e.\ if and only if its reduction is $z_0$.
\end{proof}

\begin{proposition}\label{prop:pf:steps2:cycall}
Let $p$ be a prime, $f,n\in\N$, $Y$ a scheme, $\varpi\in\Gamma(Y,\OO_Y)$, and $T_0,\dots,T_n$
$p$-correspondences on $Y$ with weights $w_j$ as in \cref{lem:pf:steps2:wtval}, such that
$\stepsiiCR_{p^f,\varpi}(T)$ holds, every point of $Y$ specializes from a point of $Y[1/p]$, and
$Y[1/p]\subseteq Y[1/\varpi]$. Then for every algebraically closed field $\Omega$ and every geometric
point $x\colon\Spec\Omega\to Y$ such that $\mathrm{pt}(x)$ is a generic point of $\{\varpi=0\}$, there is
an endomorphism $\phi$ of $\Omega$ with $\phi(a)=a^{p^f}$ such that the cycle identity holds for
$(p^f,\phi,x)$.
\formallabel{UpperTheory.Gal.Mod.cycAll_holds}\formallabel{UpperTheory.Gal.Mod.cycAll}
\end{proposition}

\begin{proof}
\emph{Step 1: from multisets to weighted cycles.} Let $V$ be a valuation ring with algebraically
closed fraction field and residue field, $p\ne0$ in $V$, $y\colon\Spec V\to Y$ with closed point a
generic point of $\{\varpi=0\}$, and $\psi(a)=a^{p^f}$ on $\kappa_V$. Apply the congruence relation
with $I_j=\stepsiiSec_{T_j}(y)$ (finite by \cref{lem:pf:steps2:wtval}) and the identity bijections,
and take the multiplicity of an arbitrary $\kappa_V$-point $y_0$ on both sides of
\eqref{eq:pf:steps2:altid}: the result is $\sum_j\stepsiico(p^f,j)N_j=0$, where $N_j$ is the number of
sections $s$ of $\alpha_{T_j}$ over $y$ with $F_\psi^{\,n-j}(\beta_{T_j}\circ\bar s)=y_0$. Grouping the
sections by their reductions $z$ and using \cref{lem:pf:steps2:wtval}, $N_j=\sum_zw_j(z)$, the sum
over the $\kappa_V$-points $z$ of $(T_j)_Z$ over the closed point of $y$ with
$F_\psi^{\,n-j}(\beta_{T_j}\circ z)=y_0$. Thus the cycle identity holds for $(p^f,\psi,\bar y)$, $\bar
y$ the closed point of $y$.

\emph{Step 2: lifting $x$.} Choose $\eta\in Y[1/p]$ specializing to $\mathrm{pt}(x)$. There are a
valuation ring $V$ with algebraically closed fraction field, a morphism $y\colon\Spec V\to Y$ sending
the generic point to $\eta$ and the closed point to $\mathrm{pt}(x)$, an algebraically closed field
$\Omega'$ and embeddings $a\colon\Omega\to\Omega'$, $b\colon\kappa_V\to\Omega'$ with $x\circ\Spec a=\bar
y\circ\Spec b$: take a valuation ring of an algebraic closure of the residue field of $\eta$
dominating the image of the local ring of $Y$ at $\mathrm{pt}(x)$, and a common algebraically closed
extension of $\Omega$ and $\kappa_V$ over the residue field of $\mathrm{pt}(x)$. Moreover $p\ne0$ in
$V$ (as $p$ is invertible at $\eta$), $\kappa_V$ is algebraically closed, and $p=0$ at
$\mathrm{pt}(x)$, because $\mathrm{pt}(x)\notin Y[1/\varpi]\supseteq Y[1/p]$. Hence $\Omega$,
$\Omega'$ and $\kappa_V$ have characteristic $p$ and carry the endomorphisms $\phi,\phi',\psi$ given by
$a\mapsto a^{p^f}$.

\emph{Step 3: change of field.} If $a\colon\Omega\to\Omega'$ is an embedding of algebraically closed
fields with $a\circ\phi=\phi'\circ a$, then the cycle identity holds for $(q,\phi,x)$ if and only if it
holds for $(q,\phi',x\circ\Spec a)$: composition with $\Spec a$ is injective on $\Omega$-points and
maps the geometric fibres of the finite morphisms $\alpha_{T_j}$ over $x$ bijectively onto those over
$x\circ\Spec a$, and weights depend only on the image point (\cref{prop:pf:steps2:traces}(a) and the
definition of $p$-weights), so the sums in \eqref{eq:pf:steps2:cycid} correspond.

\emph{Conclusion.} By Step~1 the cycle identity holds for $(p^f,\psi,\bar y)$; by Step~3 applied to
$b$ it holds for $(p^f,\phi',\bar y\circ\Spec b)=(p^f,\phi',x\circ\Spec a)$; by Step~3 applied to $a$ it
holds for $(p^f,\phi,x)$.
\end{proof}

\begin{proposition}\label{prop:pf:steps2:espt}
Let $\ell,p$ be primes, $f,n\in\N$, and $Y$, $\varpi$, $T_j$, $w_j$ as in \cref{prop:pf:steps2:cycall}.
Let $s\colon S\to Y$ and $b_S\colon S\to B$ with $b_S$ separated and locally of finite type, $\ell$
invertible on $S$, such that the set of points of $S$ mapping to generic points of $\{\varpi=0\}$ is
dense and geometric points of $S$ are determined by their images in $Y$ and in $B$. Let
$\stepsiiFr\colon S\to S$ satisfy $b_S\stepsiiFr=b_S$ and $s\circ\stepsiiFr\circ x=F_\phi(s\circ x)$ for
every geometric point $x$ of $S$ over a field with an endomorphism $\phi(a)=a^{p^f}$. Let $c_j$ be
correspondences from $S$ to $S$ that are pull-backs of $(\alpha_{T_j},\beta_{T_j},w_j)$ along $s$
(cartesian squares $v_j\colon Z_{c_j}\to(T_j)_Z$ with $sf_{c_j}=\beta_{T_j}v_j$,
$b_Sf_{c_j}=b_Sp_{c_j}$ and $w_{c_j}=w_j\circ v_j$). Then for all $k,i$,
\[
  \sum_{j=0}^n\stepsiico(p^f,j)\,\stepsiiop_{c_j}\circ(\stepsiiFr^\ast)^{n-j}=0
  \quad\text{on }H^k(S,\Ql(i)).
\]
\formallabel{UpperTheory.Gal.Mod.esPt_holds}\formallabel{UpperTheory.Gal.Mod.esPt}
\end{proposition}

\begin{proof}
Let $c''_j=(p_{c_j},\stepsiiFr^{\,n-j}\circ f_{c_j},w_{c_j})$. Then
$\stepsiiop_{c''_j}=\stepsiiop_{c_j}\circ(\stepsiiFr^\ast)^{n-j}$, and
$b_S\circ\stepsiiFr^{\,n-j}f_{c_j}=b_Sp_{c_j}$. By \cref{prop:pf:steps2:traces}(g), applied with $Y=S$
and the dense set $D$ of the hypothesis, it suffices to show
$\sum_j\stepsiico(p^f,j)\stepsiicyc_{c''_j}(x,y)=0$ for geometric points $x,y$ of $S$ over $\Omega$ with
$\mathrm{pt}(x)\in D$. If $b_S\circ y\ne b_S\circ x$, all cycles vanish since
$b_Sf_{c''_j}=b_Sp_{c''_j}$. Otherwise take $\phi$ from \cref{prop:pf:steps2:cycall} for $s\circ x$.
Composition with $v_j$ is a bijection from the geometric points $z$ of $Z_{c_j}$ with $p_{c_j}z=x$ and
$\stepsiiFr^{\,n-j}f_{c_j}z=y$ onto the geometric points $z'$ of $(T_j)_Z$ with $\alpha_{T_j}z'=s\circ x$
and $F_\phi^{\,n-j}(\beta_{T_j}z')=s\circ y$: the squares are cartesian,
$s\circ\stepsiiFr^{\,n-j}\circ t=F_\phi^{\,n-j}(s\circ t)$, and a geometric point of $S$ is determined by
its images in $Y$ and $B$. The weights correspond, so the required identity is
\eqref{eq:pf:steps2:cycid} for $(s\circ x,s\circ y)$.
\end{proof}

\begin{lemma}\label{lem:pf:steps2:corrpull}
Let $T$ be a $p$-correspondence on $Y$ over $\Spec O$ ($u\beta_T=u\alpha_T$) with an
$\alpha_T$-weight $w$. For every cartesian square with $s\colon S\to Y$, $b_S\colon S\to B$ over
$\beta\colon B\to\Spec O$, the pull-back of $(\alpha_T,\beta_T,w)$ along $s$ is a correspondence $c$ as
in \cref{prop:pf:steps2:espt}. If $s'=s\circ i$ for a morphism $i\colon S'\to S$ over a morphism of
bases, $\ell$ is invertible on $S$, and $c'$ is the pull-back along $s'$, then
$i^\ast\circ\stepsiiop_c=\stepsiiop_{c'}\circ i^\ast$.
\formallabel{UpperTheory.Gal.Mod.corrPull}\formallabel{UpperTheory.Gal.Mod.corrComm}
\end{lemma}

\begin{proof}
The first assertion is the existence of fibre products together with \cref{prop:pf:steps2:traces}(a)
for the weight; $f_c$ is induced by $\beta_T$ and $b_S$, using $u\beta_T=u\alpha_T$. The second
follows from \cref{prop:pf:steps2:traces}(c), as $c'$ is the pull-back of $c$ along $i$.
\end{proof}

\begin{proof}[Proof of \cref{prop:pf:steps2:esgen}]
Write $Y_{\bar k}=(Y_\kappa)_{\bar k}$, $Y_V$ and $Y_\Omega=(Y_K)_\Omega$, and let $\stepsiisp$ be the
isomorphism of \cref{prop:pf:steps2:specgood}(b) in degree $k$ (with $i=0$).

(a) If $r\circ\sigma|_V=r$, apply \cref{prop:pf:steps2:specgood}(b) with $\tau=1$:
$\sigma^\ast\stepsiisp=\stepsiisp$, and $\stepsiisp$ is surjective, so $\sigma^\ast=\mathrm{id}$.

(b) Since $p$ is not a unit of the discrete valuation ring $O$, $p\in(\varpi)$, so $p=0$ in $\kappa$ and
in $\bar k$. Let $\tau\in\Aut_\kappa(\bar k)$ be $a\mapsto a^{p^f}$ (an automorphism of the perfect
field $\bar k$, trivial on $\kappa$ as $a^{p^f}=a$ there); then $r\circ\sigma|_V=\tau\circ r$, so
$\stepsiisp\circ\tau^\ast=\sigma^\ast\circ\stepsiisp$. Let $\stepsiiFr$ be the relative Frobenius of
\cref{lem:pf:steps2:frob}(b) on $Y_{\bar k}$. By \cref{lem:pf:steps2:frob}(a),
$\stepsiiFr^\ast\circ\tau^\ast=(\tau_Y\circ\stepsiiFr)^\ast=F_{Y_{\bar k}}^\ast=\mathrm{id}$, so
$\stepsiiFr^\ast=(\tau^{-1})^\ast$ and
\begin{equation}\label{eq:pf:steps2:spfr}
  \stepsiisp\circ\stepsiiFr^\ast=\stepsiisp\circ(\tau^{-1})^\ast=(\sigma^{-1})^\ast\circ\stepsiisp .
\end{equation}
By \cref{lem:pf:steps2:basefacts}(a),(c) the weights $w_j=w^{\mathrm{nor}}_{\alpha_{T_j}}$ satisfy the
hypotheses of \cref{prop:pf:steps2:cycall}, and every point of $Y$ specializes from $Y[1/p]\subseteq
Y[1/\varpi]$. Let $c'_j$ be the pull-backs of $(\alpha_{T_j},\beta_{T_j},w_j)$ to $Y_{\bar k}$
(\cref{lem:pf:steps2:corrpull}). \Cref{prop:pf:steps2:espt} applies to $s\colon Y_{\bar k}\to Y$,
$B=\Spec\bar k$ and $\stepsiiFr$: $\ell$ is invertible on $Y_{\bar k}$; the density hypothesis is
\cref{lem:pf:steps2:basefacts}(b); geometric points of the fibre product $Y_{\bar k}=Y\times_{\Spec
O}\Spec\bar k$ are determined by their two projections; and $\stepsiiFr$ is a $\bar k$-morphism acting
on geometric points as in \cref{lem:pf:steps2:frob}(b). Hence
\begin{equation}\label{eq:pf:steps2:esspecial}
  \sum_j\stepsiico(p^f,j)\,\stepsiiop_{c'_j}\circ(\stepsiiFr^\ast)^{n-j}=0\quad\text{on }H^k(Y_{\bar
  k},\Ql).
\end{equation}
Let $c^V_j$ be the pull-backs to $Y_V$. By \cref{lem:pf:steps2:corrpull},
$i_Y^\ast\stepsiiop_{c^V_j}=\stepsiiop_{c'_j}i_Y^\ast$ and $j_Y^\ast\stepsiiop_{c^V_j}=\stepsiiop_{c_j}j_Y^\ast$;
for the latter note that $c_j$ is the pull-back of $T_j$ to $Y_\Omega$ with weight $1$, which agrees
with $w_j$ pulled back, since $p\ne0$ at all geometric points of $Y_\Omega$
(\cref{lem:pf:steps2:basefacts}(c)). As $i_Y^\ast$ is surjective on $H^k(-,\Ql)$ (it is bijective at
every finite level by \cref{prop:pf:steps2:specgood}(a)), we get
$\stepsiiop_{c_j}\circ\stepsiisp=\stepsiisp\circ\stepsiiop_{c'_j}$. Combining this with
\eqref{eq:pf:steps2:spfr}, for every $y\in H^k(Y_{\bar k},\Ql)$
\[
  \sum_j\stepsiico(p^f,j)\stepsiiop_{c_j}(\sigma^{-1})^{\ast(n-j)}\stepsiisp(y)
  =\stepsiisp\Bigl(\sum_j\stepsiico(p^f,j)\stepsiiop_{c'_j}(\stepsiiFr^\ast)^{n-j}y\Bigr)=0
\]
by \eqref{eq:pf:steps2:esspecial}, and $\stepsiisp$ is surjective.
\end{proof}


% ===========================================================================
\subsection{Proof of the trace formalism}\label{sec:pf:steps2:trproof}

In this subsection we prove \cref{prop:pf:steps2:traces}. The construction has three layers: the
theory of weights, which is purely geometric (\cref{prop:pf:steps2:wts}); the trace on sections of
$\Z/q(i)$, defined by the sum formula (\cref{lem:pf:steps2:trsec}); and the passage to cohomology,
through families of maps over all base changes of $p$ (\cref{lem:pf:steps2:fam,lem:pf:steps2:trchar}).
The properties of the traces are then proved with torsion coefficients
(\cref{prop:pf:steps2:trtors}) and passed to $\Zl$ and $\Ql$ coefficients
(\cref{prop:pf:steps2:tradic}).

\emph{Conventions.} Throughout, $p\colon Z\to S$ is finite. A \emph{square over $p$} is a cartesian
square formed by a morphism $u\colon U\to S$, a morphism $r\colon Z_U\to U$ and a morphism $v\colon
Z_U\to Z$ with $p\circ v=u\circ r$ (any cartesian square, not only the chosen fibre product); $r$ is
finite. A morphism from a square $(u',r',v')$ to $(u,r,v)$ is a pair $(g,k)$ with $g\colon U'\to U$,
$k\colon Z_{U'}\to Z_U$, $ug=u'$, $vk=v'$ and $rk=gr'$. The \emph{self square} is
$(\mathrm{id}_S,p,\mathrm{id}_Z)$. For a morphism $r\colon V\to U$ and a morphism $x\colon\Spec
\Omega\to U$ from the spectrum of a field, $r^{-1}(x)$ is the set of $z\colon\Spec\Omega\to V$ with
$r\circ z=x$. Besides the operations on weights of \cref{def:pf:steps2:weights} we use the
\emph{push-forward} $(h_\ast w)(z')=\sum_{z\in h^{-1}(z')}w(z)$ along a finite $h\colon Z\to Z'$; the
\emph{product} $(w'\cdot w)(y)=w'(y)\,w(r\circ y)$ of weights $w'$ on $W$ and $w$ on $Z$, for
$r\colon W\to Z$; and the \emph{multiplicity} $\stepsiim_p(z)$ of a geometric point $z$ of $Z$ with
field $\Omega$, the length of the local ring, at the point defined by $z$, of the geometric fibre
$Z\times_{S,p\circ z}\Spec\Omega$. A point of a scheme is \emph{maximal} if it has no generization
other than itself. A $p$-weight $w$ is \emph{normalized} if $w(z)=\stepsiim_p(z)$ for every geometric
point $z$ with $\mathrm{pt}(p\circ z)$ maximal in $S$; $w^{\mathrm{nor}}_p$ denotes a normalized
$p$-weight if one exists and the weight $0$ otherwise.

\subsubsection{Geometric points and weights}

\begin{lemma}[Geometric points]\label{lem:pf:steps2:gpts}
Let $r\colon V\to U$ be a morphism of schemes.
\begin{enumerate}[label=\textup{(\alph*)},nosep]
\item Every point of a scheme is $\mathrm{pt}(x)$ for a geometric point $x$.
\item If $r$ is finite, then $r^{-1}(x)$ is finite for every $x\colon\Spec\Omega\to U$, $\Omega$ a field.
\item If $r$ is integral, $x$ is a geometric point of $U$ and $r(v)=\mathrm{pt}(x)$, then there is
  $z\in r^{-1}(x)$ with $\mathrm{pt}(z)=v$.
\item If geometric points $z,z'$ of a scheme, with fields $\Omega,\Omega'$, have the same image, there
  are an algebraically closed field $L$ and morphisms $a\colon\Spec L\to\Spec\Omega$,
  $a'\colon\Spec L\to\Spec\Omega'$ with $z\circ a=z'\circ a'$.
\item If $\Omega$ is algebraically closed, $r\colon V\to\Spec\Omega$ is finite and $z,z'\colon\Spec
  L\to V$ satisfy $r\circ z=r\circ z'$ and $\mathrm{pt}(z)=\mathrm{pt}(z')$, then $z=z'$.
\item \textup{(Isolation.)} For finite $p\colon Z\to S$ and $z\in Z$ there are a square $(u,r,v)$ over $p$
  with $u$ \'etale, an open and closed $W\subseteq Z_U$ and $y\in W$ with $v(y)=z$, such that $y$ is
  the only point of $W$ in its fibre over $U$.
\item The maximal points of a scheme are dense; a point of an open subscheme is maximal there if and
  only if it is maximal in the scheme; the only maximal point of an integral scheme is its generic
  point.
\item \textup{(Common image.)} Let $S\to B$ and $Y\to B$ be morphisms, the second separated and locally
  of finite type, and let $p_a\colon Z_a\to S$ \textup{(}finite\textup{)} and $f_a\colon Z_a\to Y$,
  $a=1,\dots,r$, have the same composites to $B$. There are a finite $p_G\colon G\to S$, a morphism
  $f_G\colon G\to Y$ and finite $h_a\colon Z_a\to G$ with $p_Gh_a=p_a$, $f_Gh_a=f_a$, such that two
  morphisms $T\to G$ with the same composites with $p_G$ and with $f_G$ are equal.
\end{enumerate}
\formallabel{UpperTheory.Gal.TranscendenceTraces.ptEx_holds}\formallabel{UpperTheory.Gal.TranscendenceTraces.ptLift_holds}\formallabel{UpperTheory.Gal.TranscendenceTraces.ptCommon_holds}\formallabel{UpperTheory.Gal.TranscendenceTraces.ptOne_holds}\formallabel{UpperTheory.Gal.TranscendenceTraces.isolate_holds}\formallabel{UpperTheory.Gal.TranscendenceTraces.genBasic_holds}\formallabel{UpperTheory.Gal.TranscendenceTraces.imageFin_holds}
\end{lemma}

\begin{proof}
Finite morphisms are affine: over an affine $\Spec A\to U$ the base change of a finite $r$ is $\Spec B$
with $B$ a finite $A$-algebra. (a) Take an algebraic closure of the residue field. (b) This is also
\lref{UpperTheory.Gal.TranscendenceTraces.fibFin_holds} (fibres of finite morphisms over spectra of
fields are finite); directly, $r^{-1}(x)$ is
the set of $\Omega$-algebra maps $C\to\Omega$, $C=B\otimes_A\Omega$ finite over $\Omega$; such a map is
determined by its kernel, a maximal ideal of the artinian ring $C$, and there are finitely many.
(c) The residue field of $v$ is algebraic over that of $r(v)$, so the embedding of the latter into
the algebraically closed field of $x$ extends. (d) $\Omega\otimes_{\kappa}\Omega'$, $\kappa$ the residue
field of the common image, is a nonzero ring; take an algebraic closure of its quotient by a maximal
ideal. (e) $z$ and $z'$ factor through the local ring of $V$ at their common image, an artinian local
$\Omega$-algebra with residue field $\Omega$; an $\Omega$-algebra map from it to $L$ kills the maximal
ideal and is then determined. (f) This is the local structure of quasi-finite morphisms
\cite[18.12.1]{Grothendieck1967}, applied to the base change of $p$ over an affine open
neighbourhood of $p(z)$. (g) A point $s$ has a maximal generization (a minimal prime of
$\OO_{S,s}$), and every open set containing $s$ contains its generizations; the other two
statements are clear. (h) Let $P=Y\times_BS$, separated and locally of finite type over $S$, and
$k\colon\coprod_aZ_a\to P$ the morphism induced by the $(f_a,p_a)$. Its composite with $P\to S$ is
finite, so $k$ is finite \cite[6.1.5]{Grothendieck1961ega2}. Let $G\subseteq P$ be the scheme-theoretic
image of $k$: the morphism $\coprod_aZ_a\to G$ is finite, and $G\to S$ is integral ($\OO_G$ embeds
into the finite $\OO_S$-algebra $k_\ast\OO$) and of finite type, hence finite. Since $G\to P$ is a
closed immersion, hence a monomorphism, the last assertion holds.
\end{proof}

\begin{proposition}[Weights]\label{prop:pf:steps2:wts}
Let $p\colon Z\to S$ be finite and $w$ a $p$-weight.
\begin{enumerate}[label=\textup{(\alph*)},nosep]
\item \textup{(Operations.)} For a cartesian square $(v,p',p,g)$, $w\circ v$ is a $p'$-weight; for finite
  $r\colon W\to Z$ and an $r$-weight $w'$, $w'\cdot w$ is a $(p\circ r)$-weight; sums and integer
  multiples of $p$-weights are $p$-weights, and constants are weights for isomorphisms; for finite
  $h\colon Z'\to Z$ and a $(p\circ h)$-weight $w''$, $h_\ast w''$ is a $p$-weight.
\item \textup{(Points.)} $w(z)$ depends only on $\mathrm{pt}(z)$.
\item \textup{(Uniqueness.)} If $D\subseteq S$ is dense and $w(z)=0$ whenever $\mathrm{pt}(p\circ z)\in D$,
  then $w=0$. Hence two $p$-weights agreeing over $D$ are equal.
\item \textup{(Degree.)} If $S$ is connected, $p_\ast w$ is constant.
\item \textup{(Multiplicities.)} $\stepsiim_{p'}=\stepsiim_p\circ v$ for cartesian squares $(v,p',p,g)$;
  $\stepsiim_{p\circ j}=\stepsiim_p\circ j$ for open immersions $j$; $\stepsiim_p=1$ if $p$ is \'etale. If
  $p$ is flat, $\sum_{z\in p^{-1}(x)}\stepsiim_p(z)$ is the rank of $p$ at $\mathrm{pt}(x)$; if $p$ is
  flat of finite presentation, $\stepsiim_p$ is a $p$-weight.
\item \textup{(Normalized weights.)} If $S$ is locally normal, $w^{\mathrm{nor}}_p$ is a normalized
  $p$-weight. A scheme all of whose local rings are integrally closed domains, and which is integral or
  locally noetherian, is locally normal.
\item Let $(v,p',p,g)$ be cartesian with $p'$ flat of finite presentation, such that the points of $S'$
  mapping to maximal points of $S$ are dense, and let $w^{\mathrm{nor}}_p$ be normalized. Then
  $w^{\mathrm{nor}}_p\circ v=\stepsiim_{p'}$. In particular $w^{\mathrm{nor}}_p=\stepsiim_p$ if $p$ is flat
  of finite presentation.
\item If $S$ is integral and locally normal and $g\colon\Spec K\to S$ hits the generic point, then
  $p_\ast w^{\mathrm{nor}}_p$ is the constant $\dim_K\Gamma(Z\times_S\Spec K,\OO)$.
\end{enumerate}
\formallabel{UpperTheory.Gal.TranscendenceTraces.weights_holds}\formallabel{UpperTheory.Gal.TranscendenceTraces.wtPt_holds}\formallabel{UpperTheory.Gal.TranscendenceTraces.wtVanish_holds}\formallabel{UpperTheory.Gal.TranscendenceTraces.wtUnique_holds}\formallabel{UpperTheory.Gal.TranscendenceTraces.wtFlat_holds}\formallabel{UpperTheory.Gal.TranscendenceTraces.normSpec_holds}\formallabel{UpperTheory.Gal.TranscendenceTraces.wtNormal_holds}\formallabel{UpperTheory.Gal.TranscendenceTraces.nwtDeg_holds}
\end{proposition}

\begin{proof}
(a) In a cartesian square the fibres $r^{-1}(x)$ of the two vertical maps over corresponding
geometric points are in bijection, and pasting cartesian squares gives cartesian squares. A square
over $p'$ pasted with $(v,p',p,g)$ is a square over $p$, and the condition for $w\circ v$ is the
condition for $w$ on the pasted square. For the product, given a square $(u,\cdot,\cdot)$ over
$p\circ r$ with $W_U\to Z_U\to U$, apply the $r$-weight condition to $\varphi$ on $W_U$ (over the
square $Z_U\to Z$) and then the $p$-weight condition to the resulting locally constant function on
$Z_U$; the sums combine because $(p\circ r)^{-1}(x)=\coprod_{y\in p^{-1}(x)}r^{-1}(y)$. Sums and
multiples: add or multiply the functions $\psi$; for an isomorphism every fibre is a single point. For
$h_\ast w''$, apply the condition for $w''$ to $\varphi\circ h_U$ and regroup the sum over the fibres of
$h_U$.

(b) Let $z$ have field $\Omega$ and let $a\colon\Spec L\to\Spec\Omega$ with $L$ algebraically closed.
Consider the square over $p$ with $u=p\circ z$; its $V=Z\times_S\Spec\Omega$ is finite over
$\Spec\Omega$, hence a finite discrete space, and $z$ defines a section $z_0$ of $r\colon V\to\Spec
\Omega$. Let $\varphi$ be the indicator function of $\mathrm{pt}(z_0)$ and $\psi$ the function given by
the weight condition. Evaluating at the geometric points $\mathrm{id}$ and $a$ of $\Spec\Omega$, which
have the same image, and using (e) of \cref{lem:pf:steps2:gpts} (the only point of $r^{-1}(\mathrm{id})$,
resp.\ $r^{-1}(a)$, at $\mathrm{pt}(z_0)$ is $z_0$, resp.\ $z_0\circ a$), we get
$w(z)=\psi=w(z\circ a)$. For $z,z'$ with the same image, \cref{lem:pf:steps2:gpts}(d) gives $z\circ
a=z'\circ a'$, so $w(z)=w(z\circ a)=w(z'\circ a')=w(z')$.

(c) Let $z$ be a geometric point; take $(u,r,v)$, $W$ and $y$ as in \cref{lem:pf:steps2:gpts}(f) for
$\mathrm{pt}(z)$, let $\varphi$ be the indicator function of $W$ (locally constant) and $\psi$ given by the
weight condition. At a geometric point $x$ of $U$ with $\mathrm{pt}(x)=r(y)$, the points of $r^{-1}(x)$
lying in $W$ all have image $y$, so by (b) $\psi(r(y))=w(z)\cdot N$, where $N>0$ is the number of
$z'\in r^{-1}(x)$ with $\mathrm{pt}(z')=y$ (positive by \cref{lem:pf:steps2:gpts}(c)). The set
$\{\psi=\psi(r(y))\}$ is open, so its image under the open map $u$ is an open neighbourhood of
$p(\mathrm{pt}(z))$, which meets $D$ at some $u(t)$. At a geometric point over $t$ all terms of the sum
defining $\psi(t)$ vanish by hypothesis, so $\psi(r(y))=\psi(t)=0$ and $w(z)=0$. The second statement
follows by applying the first to the difference.

(d) The weight condition for the self square and $\varphi=1$ gives a locally constant $\psi$ with
$\psi(\mathrm{pt}(x))=(p_\ast w)(x)$; on a connected space it is constant.

(e) The geometric fibres of $p'$ are those of $p$, and those of $p\circ j$ are open subschemes of those
of $p$; the geometric fibres of an \'etale morphism are disjoint unions of copies of $\Spec\Omega$. If
$V=\Spec C$ is finite over $\Spec\Omega$, $\Omega$ algebraically closed, then $C$ is the product of its
localizations at its finitely many maximal ideals, each an artinian local $\Omega$-algebra with residue
field $\Omega$, so $\dim_\Omega C$ is the sum of their lengths; the points of $r^{-1}(\mathrm{id})$
correspond to the maximal ideals. Applied to the geometric fibre this gives the rank formula. If $p$ is
flat of finite presentation, we verify the weight condition on indicator functions of open and closed
sets, which suffices: near a point of $U$, a locally constant $\varphi$ takes finitely many values on
the quasi-compact preimage of an affine neighbourhood and is a combination of indicator functions of
open and closed sets there, and a function which near every point agrees with a locally constant one
satisfying the sum formula is itself such a function. For open and closed $W\subseteq Z_U$, the
restriction $W\to U$ of $r$ is finite flat of finite presentation, so its rank is locally constant,
and by the rank formula (for $W\to U$, with the first two statements) it equals
$\sum_{z\in r^{-1}(x),\,\mathrm{pt}(z)\in W}\stepsiim_p(v\circ z)$.

(f) \emph{Affine case.} Let $A$ be an integrally closed domain with fraction field $K$ and $B$ a
finite $A$-algebra, $p\colon\Spec B\to\Spec A$. There is a finite normal extension $L/K$ such that,
with $A_L$ the integral closure of $A$ in $L$, the set $J$ of $A$-algebra maps $\sigma\colon B\to A_L$ is
finite and, for every field $\Omega$, every ring map $\phi\colon A_L\to\Omega$ injective on $A$ and every
ring map $z\colon B\to\Omega$ agreeing with $\phi$ on $A$, there is exactly one $\sigma\in J$ with
$\phi\circ\sigma=z$ (take $L$ normal over $K$ containing the residue fields of the finitely many points
of $\Spec(B\otimes_AK)$). Let $\pi\colon T=\Spec A_L\to\Spec A$, integral and surjective, with the
sections $\sigma_j\colon T\to\Spec B$ ($j\in J$), and let $c_j\in\Spec B$ be the image under $\sigma_j$
of the generic point of $T$. Any two ring maps $A_L\to\Omega$ agreeing on $A$ differ by an
$A$-automorphism of $A_L$ \cite[Ch.~V, \S2, no.~2, Thm.~2]{Bourbaki1985}; hence for geometric points
$x,x'$ of $T$ with $\pi x=\pi x'$ there is a permutation $e$ of $J$ with $x\sigma_j=x'\sigma_{e(j)}$
and $c_{e(j)}=c_j$. Let $m_j=\stepsiim_p(z)$ for any geometric point $z$ with $\mathrm{pt}(z)=c_j$ (this
does not depend on $z$: base change to the residue field of $p(c_j)$, over which $p$ is flat, and
apply (e), (b) and (a)). Define $w(z)=\sum_{j:\,x\sigma_j=z}m_j$, where $x$ is any geometric point
of $T$ with $\pi x=p z$; by the symmetry just proved this is independent of $x$. It is a $p$-weight:
for a square $(u,r,v)$ and locally constant $\varphi$ on $Z_U$, the function
$\chi(t)=\sum_jm_j\varphi(\tau_j(t))$ on $T_U=T\times_SU$ (with $\tau_j\colon T_U\to Z_U$ induced by
$\sigma_j$) is locally constant and, by the definition of $w$ and (c) of \cref{lem:pf:steps2:gpts},
equals $\sum_{z\in r^{-1}(x)}w(vz)\varphi(\mathrm{pt}(z))$ at any geometric point over $t$, $x$ its image
in $U$; it is therefore constant on the fibres of $T_U\to U$, which is integral and surjective, hence
closed and surjective, hence a quotient map, so $\chi$ descends to a locally constant $\psi$ on $U$
with the required property. If $p z$ lies over the generic point of $\Spec A$, then $x$ is injective
on $A$, so there is exactly one $j$ with $x\sigma_j=z$, and $c_j=\mathrm{pt}(z)$; thus
$w(z)=\stepsiim_p(z)$, and $w$ is normalized.

\emph{Gluing.} If $u\colon U\to S$ is an open immersion and $w$ is a normalized $p$-weight, then
$w\circ v$ is a normalized weight for the restriction $r\colon Z_U\to U$, by (a), the first statement
of (e) and \cref{lem:pf:steps2:gpts}(g). Hence, given two open squares with normalized weights, their
restrictions to the intersection of the two opens are normalized weights there, which agree over the
dense set of maximal points and therefore everywhere, by (c). If $S$ is locally normal, cover it by
affine opens $\Spec A$ with $A$ an integrally closed domain; over each, $p$ is $\Spec B\to\Spec A$ with
$B$ finite, and the affine case gives a normalized weight. By the agreement just proved these glue to a
weight $w$ on $Z$; being a weight is local on $S$ for open covers (the weight condition can be checked
near each point of $U$, as in the proof of (e)), so $w$ is a normalized $p$-weight. Hence
$w^{\mathrm{nor}}_p$ is a normalized $p$-weight. For the second statement: on an integral scheme the ring
of an affine open is a domain, integrally closed because its localizations are
\cite[Prop.~5.13]{AtiyahMacdonald1969}; on a locally noetherian scheme with local domains, a
noetherian affine open has finitely many irreducible components, which are disjoint (a point lies on
only one, its local ring having one minimal prime), hence open, and the component through a given point
is an integral affine open subscheme, to which the integral case applies.

(g) Both sides are $p'$-weights, by (a) and (e); they agree at the geometric points over points of $S'$
mapping to maximal points of $S$, by normality and the first statement of (e); so they are equal by
(c). In the particular case, $\stepsiim_p$ is a normalized $p$-weight by (e), so a normalized weight
exists and $w^{\mathrm{nor}}_p$ is normalized; it agrees with $\stepsiim_p$ over the dense set of maximal
points, hence everywhere by (c).

(h) $S$ is connected, so $p_\ast w^{\mathrm{nor}}_p=d$ is constant by (d). Evaluate at $y\circ g$, where
$y\colon\Spec\Omega\to\Spec K$ is an algebraic closure: at points over the generic point
$w^{\mathrm{nor}}_p=\stepsiim_p$, so $d=\sum_{z\in p^{-1}(yg)}\stepsiim_p(z)=\sum_{z'\in p'^{-1}(y)}
\stepsiim_{p'}(z')=\dim_K\Gamma(Z\times_S\Spec K,\OO)$, by (e) for the flat morphism $p'$ over the field
$K$.
\end{proof}

\subsubsection{Traces on sections and on cohomology}

For $U$ with $q$ invertible on $U$ and a root $\zeta\in\Gamma(U,\OO)$ of the $q$-th cyclotomic
polynomial, let $\zeta^{(i)}$ be the section of $\Z/q(i)$ over $U$ defined by $\zeta$.

\begin{lemma}[Traces on sections]\label{lem:pf:steps2:trsec}
Let $q\ge1$ and $i\in\Z$.
\begin{enumerate}[label=\textup{(\alph*)},nosep]
\item Two sections of $\Z/q(i)$ over $U$ with the same pull-back to every geometric point of $U$ are
  equal.
\item If $\zeta$ is as above, then for every section $y$ of $\Z/q(i)$ over $U$ there is a locally
  constant $\varphi\colon U\to\Z$ with $x^\ast y=\varphi(\mathrm{pt}(x))\,x^\ast\zeta^{(i)}$ for all
  geometric points $x$, and every locally constant $\varphi$ arises in this way.
\item Let $p\colon Z\to S$ be finite, $w$ a $p$-weight and $q$ invertible on $S$. For every square
  $(u,r,v)$ over $p$ and every section $a$ of $\Z/q(i)$ over $Z_U$ there is a unique section
  $\stepsiitr^0(a)$ of $\Z/q(i)$ over $U$ with
  $x^\ast\stepsiitr^0(a)=\sum_{z\in r^{-1}(x)}w(v\circ z)\,z^\ast a$ for every geometric point $x$ of
  $U$. The maps $\stepsiitr^0$ are additive and natural with respect to morphisms of squares.
\end{enumerate}
\formallabel{UpperTheory.Gal.TranscendenceTraces.ptInj_holds}\formallabel{UpperTheory.Gal.TranscendenceTraces.secFun_holds}\formallabel{UpperTheory.Gal.TranscendenceTraces.funSec_holds}\formallabel{UpperTheory.Gal.TranscendenceTraces.sumEx_holds}\formallabel{UpperTheory.Gal.TranscendenceTraces.trSum_holds}
\end{lemma}

\begin{proof}
(a) and (b) are properties of the twisted constant sheaves: \'etale locally $\Z/q(i)$ is the constant
sheaf $\Z/q$ with generator $\zeta^{(i)}$, sections of a constant sheaf are locally constant functions,
and a locally constant function is determined by its values at the points, all of which are images of
geometric points (\cref{lem:pf:steps2:gpts}(a)). (c) Uniqueness follows from (a). Existence: since $q$
is invertible on $U$, there is an \'etale covering $\{U_k\to U\}$ such that every $U_k$ carries a root
of the $q$-th cyclotomic polynomial (adjoin one; the extension is finite \'etale and surjective). On
$U_k$, write the pull-back of $a$ to $Z_{U_k}$ as $\varphi$ by (b), apply the weight condition to the
square $Z_{U_k}\to U_k$ (a square over $p$) to get a locally constant $\psi$, and let $b_k$ be the
section attached to $\psi$ by (b); then $b_k$ satisfies the sum formula over $U_k$, because
$z^\ast a=\varphi(\mathrm{pt}(z))\,x^\ast\zeta^{(i)}$ for $z\in r^{-1}(x)$. The sum formula is
compatible with pull-back along morphisms of squares (the fibres correspond bijectively), so by
uniqueness the $b_k$ agree on overlaps and glue to a section over $U$, which satisfies the sum formula
at every geometric point (a geometric point of $U$ lifts to every \'etale covering, since \'etale
coverings of the spectrum of an algebraically closed field have sections). Additivity and naturality
follow from uniqueness, using that the fibres are finite.
\end{proof}

A \emph{family over $p$} from a complex $I$ to a complex $J$ of sheaves of modules on the big \'etale
site consists of homomorphisms $\tau_A\colon I^s(Z_U)\to J^s(U)$, for all squares $A=(u,r,v)$ over $p$
and all $s$, natural with respect to morphisms of squares and commuting with the differentials. A
\emph{family of maps on sections} from $F$ to $G$ is the same without degrees and differentials; for
instance, $\stepsiitr^0$ of \cref{lem:pf:steps2:trsec}(c) is a family of maps on sections from $\Z/q(i)$
to itself. For a sheaf $F$ we write $F\to I_F$ for the injective resolution used in
\cref{def:pf:steps2:coh}. A family $\tau$ from $I_F$ to $I_G$ \emph{lies over} a family $t$ of maps on
sections if $\tau$ and $t$ commute with the augmentations $F\to I_F^0$ and $G\to I_G^0$.

\begin{lemma}[Families]\label{lem:pf:steps2:fam}
\begin{enumerate}[label=\textup{(\alph*)},nosep]
\item Every class in $H^n(X,F)$ is the class $[a]$ of a cocycle $a\in I_F^n(X)$; $[a]=[a']$ if and only
  if $a-a'$ is a coboundary; $[a]+[a']=[a+a']$, $f^\ast[a]=[f^\ast a]$, and a morphism of complexes
  $\rho\colon I_F\to I'$ \textup{(}compatible with restrictions\textup{)} maps $[a]$ to $[\rho(a)]$.
\item A family $\tau$ over $p$ from $I_F$ to $I_G$ induces homomorphisms $\tau_\ast\colon H^n(Z,F)\to
  H^n(S,G)$ with $\tau_\ast[a]=[\tau_{\mathrm{self}}(a)]$.
\item Let $p$ be finite and $F,G$ sheaves of $\Z/Q$-modules. Every family $t$ of maps on sections from
  $F$ to $G$ lies under a family from $I_F$ to $I_G$, and if $\tau,\tau'$ both lie over $t$, then for
  every square $A$ and cocycle $a$, $\tau_A(a)-\tau'_A(a)$ is a coboundary; in particular
  $\tau_\ast=\tau'_\ast$.
\end{enumerate}
\formallabel{UpperTheory.Gal.TranscendenceTraces.rpBasic_holds}\formallabel{UpperTheory.Gal.TranscendenceTraces.rpOp_holds}\formallabel{UpperTheory.Gal.TranscendenceTraces.fL_holds}\formallabel{UpperTheory.Gal.TranscendenceTraces.opEq_holds}\formallabel{UpperTheory.Gal.TranscendenceTraces.phiEx_holds}\formallabel{UpperTheory.Gal.TranscendenceTraces.locFinSec_holds}
\end{lemma}

\begin{proof}
(a) By definition $H^n(X,F)$ is the group of homotopy classes of maps from the free chain complex of
the constant simplicial presheaf $X$ to $I_F[n]$; this chain complex is a resolution of the free
presheaf on $X$, so by Yoneda's lemma such homotopy classes correspond to the cohomology classes of
the complex $I_F(X)$, compatibly with sums, pull-backs and maps of complexes. (b) The family acts
componentwise on the complexes computing cohomology, commuting with the differentials; the formula
follows from (a).

(c) Let $\Phi(F)$ be the sheaf associated with the presheaf $U\mapsto\bigoplus_{u\colon U\to S}
F(Z\times_{S,u}U)$. By the universal properties of direct sums and sheafification, morphisms
$\Phi(F)\to G$ correspond bijectively and naturally to families of maps on sections from $F$ to $G$;
hence families over $p$ from $I_F$ to $I_G$ correspond to morphisms of complexes $\Phi(I_F)\to I_G$,
and lying over $t$ means lifting the morphism $\Phi(F)\to G$ attached to $t$. The functor $\Phi$ is
exact. The only point is: if $F_1\to F_2\to F_3$ is exact and $b$ is a section of $F_2$ over $Z_U$
mapping to $0$ in $F_3$, then $b$ lifts to $F_1$ after pull-back to $Z_{U'}$ for the members $U'$ of
an \'etale covering of $U$. By exactness $b$ lifts on an \'etale surjective $W\to Z_U$, and it suffices
to find, for each point $x$ of $U$ (we may assume $U$ affine), an \'etale $U'\to U$ hit by $x$ such
that $W$ has a section over $Z_{U'}$. Pass to the strict henselization $O$ of $U$ at $x$, which is
ind-\'etale over $U$: a finite $O$-algebra is a finite product of strictly henselian local rings, and
an \'etale faithfully flat algebra over a strictly henselian local ring has a section
\cite[Ch.~I, \S4]{Milne1980}; since the algebras involved are of finite presentation,
the section is defined over some \'etale neighbourhood of $x$ \cite[8.8.2]{Grothendieck1967}. (This is
the exactness of push-forward along finite morphisms, cf.\ \cite[Exp.~VIII, 5.5]{ArtinGrothendieckVerdier1972}.)
Consequently $\Phi(I_F)$ is a resolution of $\Phi(F)$, and by the comparison theorem for maps from a
resolution to a complex of injectives, the morphism $\Phi(F)\to G\to I_G^0$ extends to a morphism of
complexes $\Phi(I_F)\to I_G$, unique up to homotopy. A homotopy between two extensions gives, on a
cocycle $a$ over a square $A$, an element $h$ with $\tau_A(a)-\tau'_A(a)=dh$; the last statement follows
from (b) and (a).
\end{proof}

\begin{lemma}[Characterization of the trace]\label{lem:pf:steps2:trchar}
Let $p\colon Z\to S$ be finite, $w$ a $p$-weight and $q$ invertible on $S$. Define
$\stepsiitr_{p,w}\colon H^n(Z,\Z/q(i))\to H^n(S,\Z/q(i))$ as $\tau_\ast$ for a family $\tau$ lying over
$\stepsiitr^0$. Say that a family $\tau$ from $I_{\Z/q(i)}$ to itself \emph{sums $w$} if for every square
$A$ and section $a$ of $\Z/q(i)$ over $Z_U$, $\tau_A$ maps the image of $a$ in $I^0(Z_U)$ to the image
of a section satisfying the sum formula of \cref{lem:pf:steps2:trsec}(c). Then $\stepsiitr_{p,w}$ is
well defined, and $\tau_\ast=\stepsiitr_{p,w}$ for every family $\tau$ which sums $w$.
\formallabel{UpperTheory.Gal.TranscendenceTraces.trChar_holds}
\end{lemma}

\begin{proof}
A family lying over $\stepsiitr^0$ exists and $\tau_\ast$ does not depend on it, by
\cref{lem:pf:steps2:fam}(c). A family which sums $w$ lies over $\stepsiitr^0$, by the uniqueness in
\cref{lem:pf:steps2:trsec}(c).
\end{proof}

\begin{proposition}[Traces with torsion coefficients]\label{prop:pf:steps2:trtors}
Let $q\ge1$, $n,i\in\Z$, all bases on which traces are taken having $q$ invertible, and all weights
being weights for the indicated morphisms.
\begin{enumerate}[label=\textup{(\alph*)},nosep]
\item \textup{(Base change.)} For a cartesian square $(v,p',p,g)$ with $p$ finite,
  $g^\ast\circ\stepsiitr_{p,w}=\stepsiitr_{p',w\circ v}\circ v^\ast$.
\item \textup{(Composition.)} $\stepsiitr_{p\circ r,\,w'\cdot w}=\stepsiitr_{p,w}\circ\stepsiitr_{r,w'}$.
\item \textup{(Linearity.)} $\stepsiitr_{p,w+w'}=\stepsiitr_{p,w}+\stepsiitr_{p,w'}$,
  $\stepsiitr_{p,dw}=d\,\stepsiitr_{p,w}$, and $\stepsiitr_{\mathrm{id},d}=d$.
\item \textup{(Push-forward.)} For finite $h\colon Z\to Z'$ and $p\colon Z'\to S$,
  $\stepsiitr_{p\circ h,w}\circ h^\ast=\stepsiitr_{p,h_\ast w}$.
\item \textup{(Degree.)} If $p_\ast w=d$ is constant, then $\stepsiitr_{p,w}\circ p^\ast=d$.
\item \textup{(Relations.)} If $\sum_am_aw_a(z)=0$ whenever $\mathrm{pt}(p\circ z)$ lies in a dense
  subset $D$ of $S$, then $\sum_am_a\stepsiitr_{p,w_a}=0$.
\item \textup{(Cycles.)} The statement of \cref{prop:pf:steps2:traces}(g) holds with $\Z/q(i)$
  coefficients.
\end{enumerate}
\formallabel{UpperTheory.Gal.TranscendenceTraces.torsion_holds}\formallabel{UpperTheory.Gal.TranscendenceTraces.trBC_holds}\formallabel{UpperTheory.Gal.TranscendenceTraces.trComp_holds}\formallabel{UpperTheory.Gal.TranscendenceTraces.trLin_holds}\formallabel{UpperTheory.Gal.TranscendenceTraces.trPush_holds}\formallabel{UpperTheory.Gal.TranscendenceTraces.trDeg_holds}\formallabel{UpperTheory.Gal.TranscendenceTraces.trRel_holds}\formallabel{UpperTheory.Gal.TranscendenceTraces.corrEq_holds}
\end{proposition}

\begin{proof}
Each of (a)--(d) follows the same pattern: from families summing the given weights we build a
family summing the weight on the left-hand side, related to the right-hand side by restriction or
composition, and conclude with \cref{lem:pf:steps2:trchar} and \cref{lem:pf:steps2:fam}(a),(b).
Let $\tau$ be a family summing $w$ (e.g.\ one lying over $\stepsiitr^0$).

(a) A square over $p'$ pasted with $(v,p',p,g)$ is a square over $p$; restricting $\tau$ to these
squares gives a family $\tau'$ over $p'$, which sums $w\circ v$ (\cref{prop:pf:steps2:wts}(a) and the
bijection of fibres) and satisfies $\tau'_{\mathrm{self}}(v^\ast a)=g^\ast\tau_{\mathrm{self}}(a)$ by
naturality. Hence for $x=[a]$: $g^\ast\stepsiitr_{p,w}(x)=[g^\ast\tau_{\mathrm{self}}(a)]=
[\tau'_{\mathrm{self}}(v^\ast a)]=\stepsiitr_{p',w\circ v}(v^\ast x)$.

(b) For a square over $p\circ r$ with base $U$, the base changes $W_U\to Z_U\to U$ form a square over
$r$ (with base $Z_U$) and one over $p$; composing $\tau'$ (summing $w'$) on the first with $\tau$ on the
second gives a family over $p\circ r$, which sums $w'\cdot w$ because
$(p\circ r)^{-1}(x)=\coprod_{y\in p^{-1}(x)}r^{-1}(y)$; on the self square it is
$\tau_{\mathrm{self}}\circ\tau'_{\mathrm{self}}$.

(c) The sum of two families sums $w+w'$, $d\tau$ sums $dw$, and $d$ times the identity family over
$\mathrm{id}_S$ sums the constant $d$.

(d) For a square over $p$ with base $U$, pull back along $h_U\colon Z_U\to Z'_U$ and apply $\tau$ (a
family over $p\circ h$ summing $w$) to the square $Z_U\to U$; this is a family over $p$, which sums
$h_\ast w$ by regrouping the sum over the fibres of $h_U$, and on the self square it is
$a\mapsto\tau_{\mathrm{self}}(h^\ast a)$.

(e) By (d) with $h=p$ and $p=\mathrm{id}_S$, and (c): $\stepsiitr_{p,w}\circ
p^\ast=\stepsiitr_{\mathrm{id},p_\ast w}=d$.

(f) By (c), $\sum_am_a\stepsiitr_{p,w_a}=\stepsiitr_{p,\sum_am_aw_a}$; the weight $\sum_am_aw_a$ is a
$p$-weight vanishing over $D$, hence $0$ by \cref{prop:pf:steps2:wts}(c), and $\stepsiitr_{p,0}=0$ by
(c).

(g) This is the argument given after \cref{prop:pf:steps2:traces}: \cref{lem:pf:steps2:gpts}(h) gives
$G$, $p_G$, $f_G$, $h_a$; by (d) and \cref{prop:pf:steps2:wts}(a),
$\stepsiiop_{c_a}=\stepsiitr_{p_G,w'_a}\circ f_G^\ast$ with $w'_a=h_{a\ast}w_{c_a}$; since $(p_G,f_G)$ is
injective on geometric points, $w'_a(g)=\stepsiicyc_{c_a}(p_G\circ g,f_G\circ g)$; so the weights
$w'_a$ satisfy the hypothesis of (f).
\end{proof}

For a prime $\ell$ invertible on $S$, the reduction maps $H^n(-,\Z/\ell^{k+1}(i))\to
H^n(-,\Z/\ell^k(i))$ are induced by a map of resolutions over the reduction of coefficients; we define
$\stepsiitr_{p,w}$ on $H^n(Z,\Zl(i))$ levelwise and on $H^n(Z,\Ql(i))$ by extension of scalars, which
is legitimate by (a) of the following proposition. The \emph{composite} of correspondences $c$ from $S$
to $Y$ and $d$ from $Y$ to $X$ is $e=(Z_e,p_c\circ b,f_d\circ a,(w_d\circ a)\cdot w_c)$, where
$Z_e=Z_c\times_YZ_d$ (fibre product of $f_c$ and $p_d$) with projections $b\colon Z_e\to Z_c$ and
$a\colon Z_e\to Z_d$.

\begin{proposition}[Traces with $\ell$-adic coefficients]\label{prop:pf:steps2:tradic}
Let $\ell$ be a prime invertible on the bases on which traces are taken.
\begin{enumerate}[label=\textup{(\alph*)},nosep]
\item For $q\mid Q$ with $Q$ invertible, the traces commute with the reduction maps
  $H^n(-,\Z/Q(i))\to H^n(-,\Z/q(i))$. Hence the levelwise traces define $\Zl$-linear maps on
  $H^n(-,\Zl(i))$.
\item Statements \textup{(a)}, \textup{(b)}, \textup{(e)} and \textup{(g)} of
  \cref{prop:pf:steps2:trtors} hold for the $\Ql$-linear traces, and $\stepsiitr_{\mathrm{id},d}=d$.
\item For correspondences $c'$, $c$ with a cartesian square $(v,p_{c'},p_c,g)$, $f_c\circ
  v=g'\circ f_{c'}$ and $w_{c'}=w_c\circ v$, one has $g^\ast\circ\stepsiiop_c=\stepsiiop_{c'}\circ
  g'^\ast$.
\item The composite $e$ of $c$ and $d$ satisfies $\stepsiicyc_e(x,y)=\sum_t\stepsiicyc_c(x,t)\,
  \stepsiicyc_d(t,y)$ and, if $\ell$ is invertible on $S$ and $Y$, $\stepsiiop_e=\stepsiiop_c\circ
  \stepsiiop_d$.
\item For $f\colon S\to Y$ and $d\in\Z$, the graph correspondence $(S,\mathrm{id},f,d)$ has cycle $d$ at
  the pairs $(x,f\circ x)$ and $0$ elsewhere, and operator $d\,f^\ast$.
\item \textup{(Splitting.)} Statement \textup{(f)} of \cref{prop:pf:steps2:traces} holds.
\end{enumerate}
\formallabel{UpperTheory.Gal.TranscendenceTraces.adic_holds}\formallabel{UpperTheory.Gal.TranscendenceTraces.adic}\formallabel{UpperTheory.Gal.TranscendenceTraces.trRed_holds}\formallabel{UpperTheory.Gal.TranscendenceTraces.trHl_holds}\formallabel{UpperTheory.Gal.TranscendenceTraces.corrCompQ_holds}\formallabel{UpperTheory.Gal.TranscendenceTraces.corrGraphQ_holds}\formallabel{UpperTheory.Gal.TranscendenceTraces.trSplitQ_holds}
\end{proposition}

\begin{proof}
(a) Let $\rho$ be the reduction map of resolutions and $f$ the reduction $\Z/Q(i)\to\Z/q(i)$ on
sections, which commutes with restrictions. By the sum formula and \cref{lem:pf:steps2:trsec}(a),
$f\circ\stepsiitr^0_Q=\stepsiitr^0_q\circ f$; call this family of maps on sections $t$. View
$\Z/q(i)$ as a sheaf of $\Z/Q$-modules and compose with the comparison map $\kappa$ from the
cohomology of $\Z/q(i)$ as a sheaf of $\Z/q$-modules to its cohomology as a sheaf of $\Z/Q$-modules,
which is induced by a map of resolutions and is bijective (\cref{prop:pf:steps2:etfound}(j)). The
families $\kappa\rho\tau_Q$ and $\kappa\tau_q\rho$ (with $\tau_Q,\tau_q$ lying over the traces on
sections) both lie over $t$, so they induce the same map on cohomology by \cref{lem:pf:steps2:fam}(c);
since $\kappa$ is injective, $\rho\circ\stepsiitr=\stepsiitr\circ\rho$. The levelwise traces of a
compatible system are therefore compatible, and $\Zl$-linearity follows from additivity.

(b) Apply \cref{prop:pf:steps2:trtors} levelwise and extend scalars.

(c) By (b), (a) of \cref{prop:pf:steps2:trtors}: $g^\ast\stepsiitr_{p_c,w_c}f_c^\ast=
\stepsiitr_{p_{c'},w_{c'}}v^\ast f_c^\ast=\stepsiitr_{p_{c'},w_{c'}}f_{c'}^\ast g'^\ast$.

(d) A geometric point of $Z_e$ is a pair of geometric points of $Z_c$ and $Z_d$ with the same image
$t$ in $Y$; summing $w_d(a\circ\cdot)\,w_c(b\circ\cdot)$ over those with given images $x$, $y$ gives the
cycle formula. For the operators, the square $(a,b,p_d,f_c)$ is cartesian, so by base change and
composition (\cref{prop:pf:steps2:trtors}(a),(b) with $\Ql$ coefficients)
$\stepsiiop_e=\stepsiitr_{p_c,w_c}\stepsiitr_{b,w_d\circ a}a^\ast f_d^\ast=
\stepsiitr_{p_c,w_c}f_c^\ast\stepsiitr_{p_d,w_d}f_d^\ast=\stepsiiop_c\stepsiiop_d$.

(e) The only geometric point over $x$ is $x$ itself, and $\stepsiitr_{\mathrm{id},d}=d$.

(f) By base change ((b) and \cref{prop:pf:steps2:trtors}(a), with the weight $1$, which is a weight for
finite \'etale morphisms by \cref{prop:pf:steps2:wts}(e)) we may replace $p$ by its base change $p_U$
over $x$, which is split by the induced sections $s'_j$. Compare the correspondence $(Z_U,p_U,
\mathrm{id},1)$ from $U$ to $Z_U$ with the sum of the graph correspondences $(U,\mathrm{id},s'_j,1)$:
their cycles at $(x_0,y)$ are $1$ if $p_U\circ y=x_0$ and $0$ otherwise, resp.\ the number of $j$ with
$s'_j\circ x_0=y$, and these agree because the $s'_j$ identify $\coprod_jU$ with $Z_U$. By
\cref{prop:pf:steps2:trtors}(g) (with $B=U$, $Z_U\to U$ and $\mathrm{id}_U$ as structure maps, and $D=U$) and (e), $\stepsiitr_{p_U,1}=\sum_js'^\ast_j$, and
composing with the projection to $Z$ gives the claim.
\end{proof}

\begin{proof}[Proof of \cref{prop:pf:steps2:traces}]
(a) is \cref{prop:pf:steps2:wts}(a),(c),(e),(f); (b) is \cref{lem:pf:steps2:trsec}(c) and
\cref{lem:pf:steps2:trchar}; (c) is \cref{prop:pf:steps2:tradic}(b),(c); (d), (e) and (g) are
\cref{prop:pf:steps2:tradic}(b); (f) is \cref{prop:pf:steps2:tradic}(f).
\end{proof}

% ===========================================================================
\subsection{Auxiliary results}\label{sec:pf:steps2:aux}

This subsection proves auxiliary results used above and in \cref{sec:pf:steps1}: the existence of
residue data (used in the crystalline comparison), the splitting of finite algebras over valuation
rings, Galois closures and the uniqueness of operators of correspondences, and the finiteness and
continuity statements for $\ell$-adic cohomology used in \cref{lem:pf:steps2:cont}.

\subsubsection{Residue data}

Let $\ell$ be a prime, $C=\C_\ell$ the completion of $\Qlbar$, $\OO_C$ its valuation ring and
$\OO_C^\flat=\varprojlim_{x\mapsto x^\ell}\OO_C/\ell$ its tilt, with the multiplicative map
$x\mapsto x^\sharp$, $\OO_C^\flat\to\OO_C$. We use two standard facts: $x^\sharp\equiv x_0\pmod{\ell\OO_C}$,
where $x_0\in\OO_C/\ell$ is the component of index $0$; and if $k$ is a perfect ring of characteristic
$\ell$ and $g\colon k\to\OO_C/\ell$ a ring map, then $x\mapsto(g(x^{1/\ell^m}))_{m\ge0}$ is a ring map
$s\colon k\to\OO_C^\flat$ with $s(x)_0=g(x)$.

\begin{lemma}[Residue data]\label{lem:pf:steps2:resdata}
Let $K/\Ql$ be a finite extension inside $\Qlbar$. There are a perfect field $k$ of characteristic
$\ell$, a surjective ring map $\rho\colon\OO_K\to k$ and a ring map $s\colon k\to\OO_C^\flat$ with
$|s(\rho(x))^\sharp-x|<1$ for all $x\in\OO_K$.
\formallabel{CrystallineSemistableComparison.resEx}\formallabel{CrystallineSemistableComparison.exists_pow_mem_span}\formallabel{CrystallineSemistableComparison.exists_tilt_section}\formallabel{CrystallineSemistableComparison.norm_lt_one_of_mem_max}
\end{lemma}

\begin{proof}
Let $k$ be the residue field of the discrete valuation ring $\OO_K$, finite and hence perfect, and
$\rho$ the reduction map.

\emph{Step 1.} There is $n$ with $y^{\ell^n}\in\ell\OO_C$ for all $y$ in the maximal ideal
$\mathfrak m_K$. Indeed, let $\varpi$ be a uniformizer; $|\varpi|<1$, so $|\varpi|^n<|\ell|$ for some $n$.
Every $y\in\mathfrak m_K$ is $u\varpi$ with $u\in\OO_K$, so $|y^{\ell^n}|\le|\varpi|^{\ell^n}\le|\varpi|^n
\le|\ell|$, and an element of $\OO_C$ of absolute value at most $|\ell|$ lies in $\ell\OO_C$.

\emph{Step 2.} The map $\OO_K\to\OO_C/\ell$, $y\mapsto y^{\ell^n}\bmod\ell$, is a ring map (a power of
the Frobenius of the $\mathbb F_\ell$-algebra $\OO_C/\ell$ composed with the inclusion) which kills
$\mathfrak m_K$ by Step 1, so it induces $\psi\colon k\to\OO_C/\ell$. Let $g=\psi\circ\phi^{-n}$, where
$\phi$ is the Frobenius of $k$ (bijective as $k$ is perfect), and let $s\colon k\to\OO_C^\flat$ be the
ring map with $s(x)_0=g(x)$. Given $x\in\OO_K$, choose $y\in\OO_K$ with $\rho(y)=\phi^{-n}(\rho(x))$,
i.e.\ $\rho(y^{\ell^n}-x)=0$. Then $s(\rho(x))^\sharp\equiv g(\rho(x))=\psi(\rho(y))\equiv
y^{\ell^n}\pmod\ell$, so $s(\rho(x))^\sharp-x=\ell a+m$ with $a\in\OO_C$ and $m=y^{\ell^n}-x\in\mathfrak
m_K$.

\emph{Step 3.} $|\ell a|\le|\ell|<1$, and $|m|<1$ because an element of $\OO_K$ of absolute value $1$
is a unit. By the ultrametric inequality $|s(\rho(x))^\sharp-x|\le\max(|\ell a|,|m|)<1$.
\end{proof}

\subsubsection{Finite algebras over valuation rings}

\begin{lemma}[Idempotents]\label{lem:pf:steps2:hens}
Let $V$ be a valuation ring whose fraction field is algebraically closed, $A$ a finite $V$-algebra
and $\mathfrak m$ a maximal ideal of $A$. There is an idempotent $e\in A$ with $e\notin\mathfrak m$ and
$e\in\mathfrak m'$ for every maximal ideal $\mathfrak m'\ne\mathfrak m$.
\formallabel{UpperTheory.Gal.Mod.hensSplit}
\end{lemma}

\begin{proof}
Maximal ideals of $A$ contain the image of the maximal ideal $\mathfrak m_V$ of $V$, since $A$ is
integral over the local ring $V$.

\emph{Step 1.} For every $a\in A$ there is an idempotent $e_a\notin\mathfrak m$ such that for every
maximal $\mathfrak m'$ with $e_a\notin\mathfrak m'$ some $\beta\in V$ satisfies
$a-\beta\in\mathfrak m\cap\mathfrak m'$. Let $F$ be a monic polynomial over $V$ with $F(a)=0$. Since
the fraction field of $V$ is algebraically closed and $V$ is integrally closed, $F=\prod_{\alpha\in
R}(X-\alpha)$ for a multiset $R$ of elements of $V$, and $\prod_{\alpha\in R}(a-\alpha)=0$. Let $g$
(resp.\ $h$) be the product of the factors $a-\alpha$ lying in $\mathfrak m$ (resp.\ not in $\mathfrak
m$). Then $gh=0$, $h\notin\mathfrak m$, hence $g\in\mathfrak m$. If $a-\alpha\in\mathfrak m$ and
$a-\beta\notin\mathfrak m$, then $\alpha-\beta$ is a unit of $V$ (otherwise
$a-\beta=(a-\alpha)+(\alpha-\beta)\in\mathfrak m$); so the factors of $g$ and of $h$ are pairwise
comaximal, and $ug+vh=1$ for some $u,v\in A$. Put $e_a=vh$: then $e_a(1-e_a)=vh\,ug=0$, and
$1-e_a=ug\in\mathfrak m$, so $e_a\notin\mathfrak m$. If $e_a\notin\mathfrak m'$, then $h\notin\mathfrak
m'$; since $\prod_\alpha(a-\alpha)=0\in\mathfrak m'$, some factor $a-\alpha$ lies in $\mathfrak m'$, and
it is a factor of $g$, so $a-\alpha\in\mathfrak m\cap\mathfrak m'$.

\emph{Step 2.} Let $t_1,\dots,t_r$ generate $A$ as a $V$-module and $e=\prod_ie_{t_i}$, an idempotent
not in the prime $\mathfrak m$. Let $\mathfrak m'\ne\mathfrak m$ be maximal and suppose $e\notin
\mathfrak m'$. Then each $e_{t_i}\notin\mathfrak m'$, so each $t_i$ is congruent to an element of $V$
modulo $\mathfrak m\cap\mathfrak m'$, and hence so is every element of $A$. If $x\in\mathfrak m$ and
$x-\beta\in\mathfrak m\cap\mathfrak m'$ with $\beta\in V$, then $\beta\in\mathfrak m$, so $\beta$ is not
a unit, $\beta\in\mathfrak m_V$ and $\beta\in\mathfrak m'$; hence $x\in\mathfrak m'$. Thus
$\mathfrak m\subseteq\mathfrak m'$ and $\mathfrak m=\mathfrak m'$, a contradiction.
\end{proof}

\subsubsection{Galois closures and operators of correspondences}

A morphism $s\colon Y\to Z$ is \emph{finite \'etale of degree $d$} if it is finite and \'etale and
every geometric fibre $s^{-1}(x)$ has exactly $d$ elements.

\begin{lemma}[Galois closures]\label{lem:pf:steps2:galclos}
Let $s\colon Y\to Z$ be finite \'etale and $d\in\N$.
\begin{enumerate}[label=\textup{(\alph*)},nosep]
\item There are a finite \'etale $c\colon W\to Z$ and $\sigma_1,\dots,\sigma_d\colon W\to Y$ with
  $s\sigma_j=c$ and $j\mapsto\sigma_j\circ x$ injective for every $x\colon\Spec\Omega\to W$
  \textup{(}$\Omega$ a field\textup{)}, universal for these properties: for every $z\colon U\to Z$ and
  $y_1,\dots,y_d\colon U\to Y$ with $sy_j=z$ and $j\mapsto y_j\circ x$ injective for all such $x$, there
  is a unique $w\colon U\to W$ with $cw=z$ and $\sigma_jw=y_j$.
\item If $s$ is finite \'etale of degree $d$, then there are a finite group $G$ acting on $W$ such that
  $c$ is a Galois cover \textup{(}\cref{def:pf:steps2:hop}\textup{)}, and the $\sigma_j$ split $s$ over
  $c$.
\end{enumerate}
\formallabel{UpperTheory.Gal.G3.pOrdEx}\formallabel{UpperTheory.Gal.G3.pOrdStep}\formallabel{UpperTheory.Gal.G3.pGalClosD}\formallabel{UpperTheory.Gal.Mod.galEx}
\end{lemma}

\begin{proof}
(a) Induction on $d$; for $d=0$ take $W=Z$. Given $(W,c,\sigma_1,\dots,\sigma_d)$, let
$P=Y\times_ZW$ with projections $\pi_Y,\pi_W$ and $\varphi_j=(\pi_Y,\sigma_j\pi_W)\colon P\to
Y\times_ZY$. The diagonal of the \'etale separated morphism $s$ is an open and closed immersion
\cite[17.4.2]{Grothendieck1967}, so the union of the preimages under the $\varphi_j$ of the diagonal is open
and closed in $P$; let $W'$ be its complement, $c'=c\pi_W$, $\sigma'_j=\sigma_j\pi_W$ ($j\le d$) and
$\sigma'_{d+1}=\pi_Y$. A morphism $x\colon\Spec\Omega\to P$ factors through $W'$ if and only if
$\sigma_j\pi_Wx\ne\pi_Yx$ for all $j$; this gives the injectivity for $W'$ and, with the universal
property of $W$ and of the fibre product, the universal property of $W'$ (a morphism from $U$ factors
through the open subscheme $W'$ if and only if all its points do, which is checked on
$\Spec\Omega$-valued points). $W'\to W$ is finite \'etale, being the restriction of the base change
$P\to W$ of $s$ to an open and closed subscheme.

(b) For $g$ in the symmetric group $G=S_d$, the universal property applied to
$(c,\sigma_{g(1)},\dots,\sigma_{g(d)})$ gives $a(g)\colon W\to W$ with $ca(g)=c$ and
$\sigma_ja(g)=\sigma_{g(j)}$; by uniqueness $a(1)=\mathrm{id}$ and $a(gh)=a(h)\circ a(g)$. The $\sigma_j$
split $s$ over $c$ by the criterion \lref{UpperTheory.Gal.G3.pSplitsPts} (finitely many
sections of a finite \'etale $s$ over $c$ split $s$ if for every geometric point $x$ of $W$ and every
geometric point $y$ of $Y$ with $sy=cx$ there is exactly one $j$ with $\sigma_jx=y$): the $d$ points
$\sigma_jx$ are distinct and lie in the fibre of $s$ over $cx$, which has exactly $d$ elements. The
morphism $c$ is surjective: over a geometric point $x$ of $Z$ the $d$ elements of $s^{-1}(x)$, in any
order, give by the universal property a geometric point of $W$ over $x$. Finally, the morphism
$\coprod_{g\in G}W\to W\times_ZW$ given by $(a(g),\mathrm{id})$ on the $g$-th summand is a morphism of
schemes \'etale over $W$ which is bijective on geometric points (a geometric point of $W\times_ZW$ is a
pair of orderings of the same fibre, and exactly one $g$ carries the second to the first), hence an
isomorphism by \lref{UpperTheory.Gal.G3.pIsoPts} (an \'etale morphism inducing bijections on
geometric points is an isomorphism). So $c$ is a Galois cover with group $G$.
\end{proof}

\begin{lemma}[Uniqueness of Hecke operators]\label{lem:pf:steps2:opuniq}
Let $k$ be a field with $\ell\ne0$ in $k$, $H$ a Hecke tower over $k$, $\Omega/k$ an extension and
$g\in\Gamma$ with $s_g$ finite \'etale of some degree. Every operator of the correspondence
$(s_{g,\Omega},t_{g,\Omega})$ equals the Hecke operator $T_g$.
\formallabel{UpperTheory.Gal.Int.Geo.opUniq_of}\formallabel{UpperTheory.Gal.G3.pOpEx}\formallabel{UpperTheory.Gal.G3.pOpEt}
\end{lemma}

\begin{proof}
The base change $s_{g,\Omega}$ is finite \'etale of the same degree and $\ell$ is invertible on
$Z_\Omega$. By \cref{lem:pf:steps2:galclos}(b), $s_{g,\Omega}$ is split over a Galois cover, so (as in
\cref{lem:pf:steps2:hopop}) an operator exists and $T_g$ is an operator. By
\cref{lem:pf:steps2:corropeq}, applied to the correspondence $(Y_{g,\Omega},s_{g,\Omega},t_{g,\Omega},
1)$ from $Z_\Omega$ to $Z_\Omega$, every operator of $(s_{g,\Omega},t_{g,\Omega})$ equals
$\bigoplus_n\stepsiiop_c$; in particular it equals $T_g$. The formal proof runs the same argument in
an axiomatic setting: \lref{UpperTheory.Gal.G3.pOpTr} states that, for a cohomology theory in
which pull-back along Galois covers is injective and finite \'etale morphisms have traces satisfying
the splitting identity, every operator of a correspondence $(s,t)$ equals the trace along $s$
composed with $t^\ast$; \'etale $\ell$-adic cohomology over a base on which $\ell$ is invertible
satisfies these hypotheses (\lref{UpperTheory.Gal.crystallineComparison.Br.etGood}), and
splittings are preserved by base change (\lref{UpperTheory.Gal.crystallineComparison.Br.bcSplits}).
\end{proof}

\subsubsection{Finiteness and continuity of \texorpdfstring{$\ell$}{l}-adic cohomology}

For a space $E$ (a simplicial presheaf as in \cref{def:pf:steps2:coh}) we say that $\ell$ is
invertible on $E$ if it is invertible on every scheme $U$ with $E_0(U)\ne\emptyset$. For $x\in
H^n(E,\Zl(i))$ we write $x_k\in H^n(E,\Z/\ell^k(i))$ for its component of level $k$.

\begin{lemma}[Limits]\label{lem:pf:steps2:limits}
Let $F$ be a sheaf of modules on the big \'etale site such that $A\mapsto F(\Spec A)$ commutes with
filtered colimits of rings, and let $X=\varprojlim X_i$ be a cofiltered limit of quasi-compact
quasi-separated schemes with affine transition maps. Then every class in $H^n(X,F)$ is the pull-back
of a class on some $X_i$, and a class on $X_i$ whose pull-back to $X$ vanishes has vanishing
pull-back to some $X_j$, $j\to i$.
\formallabel{GlobalTheory.T.contCone_of_isLimit_of_qcqs_affine_transitions}
\end{lemma}

\begin{proof}
This is the continuity of \'etale cohomology \cite[Exp.~VII, 5.7]{ArtinGrothendieckVerdier1972},
applied to the restrictions of $F$ to the $X_i$; the hypothesis on $F$ ensures that the restriction
of $F$ to $X$ is the colimit of the pull-backs of these restrictions (an \'etale $X$-scheme of finite
presentation descends to some $X_i$ \cite[8.8.2]{Grothendieck1967}, and the values of $F$ on affine
schemes commute with the corresponding colimits of rings).
\end{proof}

\begin{lemma}[Bockstein sequences]\label{lem:pf:steps2:bock}
Let $\ell$ be invertible on $E$, $i\in\Z$ and $a,b\in\N$.
\begin{enumerate}[label=\textup{(\alph*)},nosep]
\item Write $H^n_c=H^n(E,\Z/\ell^c(i))$. There are homomorphisms $u_a\colon H^n_a\to H^n_{a+b}$ and
  $\delta_a\colon H^n_b\to H^{n+1}_a$ such that
  \[
    \cdots\to H^n_a\xrightarrow{u_a}H^n_{a+b}\xrightarrow{\mathrm{red}}H^n_b\xrightarrow{\delta_a}
    H^{n+1}_a\to\cdots
  \]
  is exact, $u_a\circ\mathrm{red}=\ell^b$ on $H^n_{a+b}$, and $u_a$, $\delta_a$ commute
  with the reductions from level $a+1$ to level $a$.
\item If moreover all $H^n(E,\Z/\ell^m(i))$ are finite and $x\in H^n(E,\Zl(i))$ has $x_b=0$, then
  $x=\ell^by$ for some $y\in H^n(E,\Zl(i))$.
\end{enumerate}
\formallabel{LateralTheory.Fin.Bock.etale_bockstein_sequence_adic_finiteness}\formallabel{UpperTheory.Gal.Int.Geo.exists_eq_pow_smul_of_lev_eq_zero}
\end{lemma}

\begin{proof}
(a) The sequence of sheaves $0\to\Z/\ell^a(i)\xrightarrow{\ell^b}\Z/\ell^{a+b}(i)\to\Z/\ell^b(i)\to0$
is exact (\'etale locally, where a primitive $\ell^{a+b}$-th root of unity exists, it is the sequence
of constant sheaves $0\to\Z/\ell^a\to\Z/\ell^{a+b}\to\Z/\ell^b\to0$). Viewing it as a sequence of
sheaves of $\Z/\ell^{a+b}$-modules, its long exact cohomology sequence
(\cref{prop:pf:steps2:etbasic}(f)), combined with the invariance of cohomology under restriction of
scalars (\cref{prop:pf:steps2:etbasic}(g)) to identify the outer terms with the cohomology of
$\Z/\ell^a(i)$ and $\Z/\ell^b(i)$, gives the exact sequence. The composite of reduction and
multiplication by $\ell^b$ on sheaves is multiplication by $\ell^b$, and the compatibilities with
reduction follow from the naturality of the long exact sequence, using
\cref{prop:pf:steps2:etbasic}(h) to compare maps of sheaves through their effect on sections and the
transitivity of reductions (\cref{prop:pf:steps2:etfound}(j)).

(b) For $a\ge0$ let $Y_a=\{y\in H^n_a:u_a(y)=x_{a+b}\}$. It is finite, and nonempty by
exactness since $\mathrm{red}(x_{a+b})=x_b=0$; reduction maps $Y_{a+1}$ to $Y_a$ by the compatibility in
(a). An inverse limit of nonempty finite sets is nonempty, so there is a compatible family
$y=(y_a)\in H^n(E,\Zl(i))$ with $y_a\in Y_a$. Then $(\ell^by)_{a+b}=\ell^by_{a+b}=u_a(y_a)=x_{a+b}$ for
all $a$, while for $k<b$ both $x_k$ (a reduction of $x_b=0$) and $\ell^by_k$ vanish. Hence $x=\ell^by$.
\end{proof}

\begin{lemma}[$\ell$-adic finiteness]\label{lem:pf:steps2:adicfin}
Let $\ell$ be a prime, $E$ a space and $n,i\in\Z$. If all $H^n(E,\Z/\ell^k(i))$ are finite and every
$x\in H^n(E,\Zl(i))$ with $x_k=0$ is divisible by $\ell^k$ in $H^n(E,\Zl(i))$, then $H^n(E,\Zl(i))$ is
a finitely generated $\Zl$-module. In particular this holds if all $H^n(E,\Z/\ell^k(i))$ are finite
and $\ell$ is invertible on $E$.
\formallabel{LateralTheory.EWeil.AdicFinPf.AdicFin_of}
\end{lemma}

\begin{proof}
Let $M=H^n(E,\Zl(i))$. The image $S$ of $M$ in level $1$ is finite; choose $g_t\in M$ with
$(g_t)_1=t$ for $t\in S$. Every $x\in M$ is congruent modulo $\ell M$ to some $g_t$ (apply the
hypothesis with $k=1$ to $x-g_t$, $t=x_1$); iterating, $x\in\sum_t\Zl g_t+\ell^kM$ for every $k$, so the
set $T_k$ of $c\in(\Z/\ell^k)^S$ with $x_k=\sum_tc_t(g_t)_k$ is nonempty (the level-$k$ component of
$\ell^kM$ vanishes). The sets $T_k$ are finite and reduction of coefficients maps $T_{k+1}$ to $T_k$;
an inverse limit of nonempty finite sets is nonempty, so there is a compatible family $(c^{(k)})$, which
defines $\alpha\in\Zl^S$ with $\alpha\bmod\ell^k=c^{(k)}$. Then $x$ and $\sum_t\alpha_tg_t$ have the
same component in every level, so they are equal. Hence the $g_t$ generate $M$. The last statement
follows from \cref{lem:pf:steps2:bock}(b).
\end{proof}

\begin{lemma}[Continuity]\label{lem:pf:steps2:abscont}
Let $M$ be a finitely generated $\Zl$-module and $K_0\supseteq K_1\supseteq\cdots$ submodules with
$\bigcap_kK_k=0$.
\begin{enumerate}[label=\textup{(\alph*)},nosep]
\item For every $N$ there is $k$ with $K_k\subseteq\ell^NM$.
\item Let a topological group $G$ act $\Zl$-linearly on $M$ such that $\{g:gx-x\in K_k\}$ is open for
  all $x\in M$ and $k$. Then $g\mapsto\varphi(gx)$ is continuous for every $\Zl$-linear
  $\varphi\colon M\to\Ql$ and $x\in M$.
\end{enumerate}
\formallabel{LateralTheory.EWeil.ChevPf.Chev_of}\formallabel{LateralTheory.EWeil.AbsContPf.AbsCont_of}
\end{lemma}

\begin{proof}
(a) Choose a surjection $\pi\colon\Zl^r\to M$ and let $V_k=\pi^{-1}(K_k)$, a decreasing sequence of
finitely generated, hence compact, submodules of $\Zl^r$ with $\bigcap_kV_k=\ker\pi$. The set
$U=\ker\pi+\ell^N\Zl^r$ is open and contains $\bigcap_kV_k$, so by compactness $V_k\subseteq U$ for some
$k$, and $K_k=\pi(V_k)\subseteq\pi(U)=\ell^NM$.

(b) $\varphi$ is bounded on $M$, say $|\varphi|\le C$: it is bounded on a finite set of generators and
$|\varphi(\lambda m+\mu m')|\le\max(|\varphi(m)|,|\varphi(m')|)$ for $\lambda,\mu\in\Zl$. Let $g_0\in G$
and $\varepsilon>0$; choose $N$ with $\ell^{-N}C<\varepsilon$ and $k$ as in (a). With $y=g_0x$, the set
of $g$ with $(gg_0^{-1})y-y\in K_k$ is an open neighbourhood of $g_0$, and for such $g$,
$gx-g_0x\in K_k\subseteq\ell^NM$, so $|\varphi(gx)-\varphi(g_0x)|\le\ell^{-N}C<\varepsilon$.
\end{proof}

% ===========================================================================
\subsection{Foundational results proved in the formal development}\label{sec:pf:steps2:found}

The results collected here are used in the proofs above and in several proofs of
\cref{sec:pf:steps1}. The development proves all of them from first principles (several with formal
proofs of $10^3$ to $10^4$ Lean lines). They are of three kinds. (1) Classical theorems: for these we
give a precise reference and, where the formal statement is a special case or a slight reformulation
of the published one, the deduction. (2) Properties of the construction of cohomology in
\cref{def:pf:steps2:coh}, and other statements specific to the development with short proofs: these
are proved below. (3) Internal steps of the formal proofs of the crystalline comparison theorem and of
the base change of towers, which are specific to the development and have long formal proofs: these
are stated with their formal tags, but their proofs are not converted here; they are listed in
\cref{rem:pf:steps2:deferred}. Throughout, $q\ge1$, $i\in\Z$, and cohomology is as in \cref{def:pf:steps2:coh}; ``$f$
induces bijections'' means that $f^\ast$ is bijective on $H^n(-,\Z/q(i))$ for all $n$.

\begin{proposition}[\'Etale cohomology]\label{prop:pf:steps2:etfound}
\begin{enumerate}[label=\textup{(\alph*)},leftmargin=2.2em]
\item (Affine analogue of proper base change.) If $(A,I)$ is a henselian pair and $q\in A^\times$,
  then $\Spec(A/I)\to\Spec A$ induces bijections \cite{Gabber1994}.
\item (Smooth base change, generic fibre.) If $T$ is an absolutely integrally closed local domain with
  $q\in T^\times$ and $X\to\Spec T$ is smooth, then the inclusion of the generic fibre
  $X\times_T\mathrm{Frac}(T)\to X$ induces bijections.
\item (The relative line.) If $T$ is as in (b) and $B$ is an \'etale $T[x]$-algebra, then
  $\Spec(\mathrm{Frac}(T)\otimes_TB)\to\Spec B$ induces bijections.
\item (Homotopy invariance.) If $q\in A^\times$, then $\Spec A[x]\to\Spec A$ induces bijections
  \cite[Exp.~XV, 2.2]{ArtinGrothendieckVerdier1972}; the development first proves this for strictly
  henselian $A$.
\item (Laurent polynomials.) Let $A$ be strictly henselian with $q\in A^\times$, and
  $A_x=A[x,x^{-1}]$. Then $H^n(\Spec A_x,\Z/q(i))=0$ for $n\ge2$, and
  $H^1(\Spec A_x,\Z/q(i))\cong\Z/q$.
\item (Fields.) Let $K$ be a field with separable closure $K^s$, $q\in K^\times$, $\zeta\in K^s$ a
  primitive $q$-th root of unity, and $M$ a discrete $\Z/q[\Gal(K^s/K)]$-module free of rank one with
  basis $b$ such that $\sigma b=a^ib$ whenever $\sigma\zeta=\zeta^a$. Then $H^n(\Spec K,\Z/q(i))\cong H^n(\Gal(K^s/K),M)$ for all $n\ge0$
  \cite[Exp.~VIII, 2.1--2.3]{ArtinGrothendieckVerdier1972}.
\item (Tsen.) If $k$ is algebraically closed and $K$ is algebraic over $k(t)$, then
  $\mathrm{cd}_p\Gal(K^s/K)\le1$ for every prime $p$, and $H^n(\Spec K,\Z/q(i))=0$ for $n\ge2$ if $q\in
  K^\times$.
\item (Affine curves.) If $k$ is algebraically closed, $q\in k^\times$ and $B$ is a Dedekind domain of
  finite type over $k$, then $H^n(\Spec B,\Z/q(i))=0$ for $n\ge2$
  \cite[Exp.~IX, 5.7]{ArtinGrothendieckVerdier1972}.
\item (Kummer theory.) There is a monomorphism $\kappa\colon\Z/q(1)\to\mathbb G_m$ with $q\kappa=0$,
  sending the generator attached to a primitive root $\zeta$ to $\zeta$, such that over every scheme on
  which $q$ is invertible every $q$-torsion section of $\mathbb G_m$ comes from $\Z/q(1)$ and every
  section of $\mathbb G_m$ has $q$-th roots locally; its boundary maps $H^n(-,\mathbb G_m)\to
  H^{n+1}(-,\Z/q(1))$ are natural, fit into the long exact sequence, and agree up to a sign with the
  Kummer class of units ($n=0$) and the first Chern class of line bundles ($n=1$)
  \cite[Ch.~II, 2.18(b); Ch.~III, \S4]{Milne1980}. If the Kummer class of a unit $w$ vanishes, $w$ is a $q$-th power of
  a unit. If $q\in A^\times$, every unit of $A$ is a $q$-th power and $\Pic(A)$ has no $q$-torsion,
  then $H^1(\Spec A,\Z/q(1))=0$.
\item (Formal properties.) Cohomology vanishes in negative degrees, and $H^0$ is the group of global
  sections; restriction of scalars along a ring map does not change cohomology; reduction of
  coefficients $\Z/Q\to\Z/q$ ($q\mid Q$) is transitive and compatible with cup products, units and
  first Chern classes; cup products are associative and graded commutative; the cohomology of
  $\Spec\varinjlim A_j$, for a filtered colimit of rings and a sheaf commuting with filtered colimits,
  is the colimit of the cohomologies.
\item (The package used for duality.) Statements (h), (c), (d), (e) and the projective bundle formula
  hold simultaneously; the last is \lref{GlobalTheory.EtC.projectiveBundleFormula_vectorBundle_any_space},
  which states that for a vector bundle of rank $N$ on any space the classes of the cohomology of the
  base cupped with the powers $c_1^j$, $0\le j<N$, of the tautological line bundle give an
  isomorphism with the cohomology of the projective bundle. From this package and the properties
  above the development derives Poincar\'e duality and finiteness for smooth projective varieties
  (\cref{lem:pf:steps2:fin}).
\end{enumerate}
\formallabel{UpperTheory.Et.affine_proper_base_change}\formallabel{UpperTheory.Et.etale_base_change_of_duality}\formallabel{UpperTheory.Cb.etale_cohomology_ind_etale_base_change}\formallabel{GlobalTheory.T.etale_cohomology_ind_etale_base_change}\formallabel{GlobalTheory.EtC.HomotopyInv.homotopy_invariance_etale}\formallabel{GlobalTheory.T.homotopyInvariance_tateTwist_strictlyHenselian}\formallabel{GlobalTheory.T.homotopy_invariance_etale_alt}\formallabel{GlobalTheory.EtC.laurentPolynomial_cohomology_strictlyHenselian_vanishing}\formallabel{GlobalTheory.T.multiplicativeGroup_cohomology_strictlyHenselian}\formallabel{GlobalTheory.T.etale_cohomology_locally_acyclic}\formallabel{GlobalTheory.T.etaleCohomology_field_eq_galoisCohomology}\formallabel{GlobalTheory.T.cdp_galois_le_one_of_isAlgebraic_ratFunc}\formallabel{GlobalTheory.T.etale_cohomology_vanishing_algebraic_over_ratFunc}\formallabel{GlobalTheory.EtC.etC_c2}\formallabel{GlobalTheory.T.subsingleton_Ht_two_of_isDedekindDomain_of_isAlgClosed}\formallabel{GlobalTheory.T.smooth_affine_curve_etale_cohomologicalDimension_le_one}\formallabel{GlobalTheory.T.kummer_sequence_etale_tateTwist}\formallabel{GlobalTheory.T.exists_kummer_boundary_etale_cohomology}\formallabel{UpperTheory.exists_root_of_kummer_class_eq_zero}\formallabel{GlobalTheory.T.kummer_etale_vanishing}\formallabel{GlobalTheory.T.etaleCohomology_negative_vanishing_degree_zero}\formallabel{GlobalTheory.EtC.a8}\formallabel{GlobalTheory.EtC.RedCoef.a11}\formallabel{GlobalTheory.T.tateTwistCupProduct_assoc}\formallabel{GlobalTheory.T.tateTwistCupProduct_graded_comm}\formallabel{GlobalTheory.T.etaleCohomology_filteredColimit_rings_finitary}\formallabel{UpperTheory.Gal.KunnethLefschetz.core_holds}
\end{proposition}

The formal tags of \cref{prop:pf:steps2:etfound} are listed in the order of the items (a)--(k); some items are proved by two or three formal theorems with the same statement.

\begin{proof}
For a scheme $X$ the cohomology of \cref{def:pf:steps2:coh} with coefficients in $\Z/q(i)$ agrees with
the cohomology of the small \'etale site of $X$ \cite[Ch.~III, \S3]{Milne1980}; we use this
tacitly to apply results stated for the small site. Recall that $\Z/q(i)$ is isomorphic to $\Z/q$ over
a scheme carrying a primitive $q$-th root of unity, e.g.\ a strictly henselian local scheme with $q$
invertible.

(a) This is the case of the torsion sheaf $\Z/q(i)$ of Gabber's affine analogue of the proper base
change theorem \cite[Thm.~1]{Gabber1994}.

(b) An absolutely integrally closed local domain $T$ is strictly henselian (every monic polynomial over
$T$ is a product of linear factors), and so is its localization $T_{\mathfrak p}$ at every prime;
the fraction field $K$ of $T$ is algebraically closed, so $\eta=\Spec K$ is a geometric point, and
$q$-th roots of unity exist in $T$. Let $j\colon X_\eta\to X$ be the inclusion of the generic fibre
(quasi-compact and quasi-separated) and $\Lambda=\Z/q$. Let $\bar x$ be a geometric point of $X$
over $\mathfrak p\in\Spec T$; the strict localization of $\Spec T$ at $\mathfrak p$ is
$\Spec T_{\mathfrak p}$, with generic point $\eta$. By \cite[Exp.~VIII, 5.2]{ArtinGrothendieckVerdier1972},
$(R^mj_\ast\Lambda)_{\bar x}=H^m(X_{(\bar x)}\times_{\Spec T_{\mathfrak p}}\eta,\Lambda)$, where
$X_{(\bar x)}$ is the strict localization of $X$ at $\bar x$. Smooth morphisms are locally acyclic
\cite[Exp.~XV, 2.1]{ArtinGrothendieckVerdier1972}, so this group is $\Lambda$ for $m=0$ and $0$ for
$m>0$. Hence $\Lambda\to Rj_\ast\Lambda$ is an isomorphism, and the Leray spectral sequence of $j$
shows that $H^n(X,\Lambda)\to H^n(X_\eta,\Lambda)$ is bijective.

(c) $\Spec B\to\Spec T[x]\to\Spec T$ is smooth, and $\Spec(\mathrm{Frac}(T)\otimes_TB)$ is its generic
fibre; apply (b).

(d) The morphism $\mathbb A^1_A\to\Spec A$ is acyclic for coefficients prime to the residue
characteristics \cite[Exp.~XV, 2.2]{ArtinGrothendieckVerdier1972}, which gives the statement by the
Leray spectral sequence.

(e) Let $S=\Spec A$, $X=\mathbb A^1_S$, $Y\subseteq X$ the zero section and $U=X\setminus Y=\Spec A_x$.
By (d) and since strictly local schemes have no higher cohomology
\cite[Exp.~VIII, 4.7]{ArtinGrothendieckVerdier1972}, $H^m(X,\Lambda)=H^m(S,\Lambda)=0$ for $m>0$ and
$H^0(X,\Lambda)=\Lambda$. By relative cohomological purity for the smooth $S$-pair $(X,Y)$ of
codimension $1$ \cite[Exp.~XVI, 3.7]{ArtinGrothendieckVerdier1972}, the cohomology of $X$ with supports
in $Y$ is $H^m_Y(X,\Lambda)=H^{m-2}(S,\Lambda(-1))$, i.e.\ $\Lambda(-1)(S)\cong\Z/q$ for $m=2$ and $0$
otherwise. The exact sequence $H^m(X)\to H^m(U)\to H^{m+1}_Y(X)\to H^{m+1}(X)$ gives
$H^1(U,\Lambda)\cong\Z/q$ and $H^m(U,\Lambda)=0$ for $m\ge2$. The formal proof follows a different,
elementary route.

(f) Sheaves on the small \'etale site of $\Spec K$ are the discrete $\Gal(K^s/K)$-modules, and their
cohomology is Galois cohomology \cite[Exp.~VIII, 2.1--2.3]{ArtinGrothendieckVerdier1972}; $\Z/q(i)$
corresponds to the module $M$ of the statement.

(g) The Galois group of $K$ is a closed subgroup of that of $k(t)$ (restriction to the separable
closure of $k(t)$ inside $K^s$). By Tsen's theorem $k(t)$ is a $C_1$ field and hence
$\mathrm{cd}_p\le1$ for $p\ne\mathrm{char}\,k$ \cite[II.3.2, II.3.3]{Serre1997}, while
$\mathrm{cd}_p\le1$ for $p=\mathrm{char}\,k$ holds for every field of characteristic $p$
\cite[II.2.2]{Serre1997}; closed subgroups have smaller cohomological dimension
\cite[I.3.3]{Serre1997}. The vanishing follows by (f).

(h) $\Spec B$ is a point or an affine smooth curve over the algebraically closed field $k$;
\cite[Exp.~IX, 5.7]{ArtinGrothendieckVerdier1972} gives the vanishing for $n\ge2$.

(i) Identify $\Z/q(1)$ with $\mu_q$; the Kummer sequence $0\to\mu_q\to\mathbb G_m\xrightarrow{q}
\mathbb G_m\to0$ is exact on the \'etale site when $q$ is invertible
\cite[Ch.~II, 2.18(b)]{Milne1980}, and $H^1(X,\mathbb G_m)=\Pic(X)$ \cite[Ch.~III, 4.9]{Milne1980}. The
long exact sequence gives the two last statements: a unit with vanishing Kummer class is in the image
of multiplication by $q$ on $H^0(X,\mathbb G_m)$; and $H^1(X,\mu_q)$ is an extension of $\Pic(X)[q]$ by
$\Gamma(X,\OO)^\times/q$. The comparisons of the boundary maps with the development's Kummer classes and
first Chern classes are properties of its definitions (both are defined through this boundary).

(j) By \cref{lem:pf:steps2:fam}(a), $H^n(X,F)=H^n(I_F(X))$, which vanishes for $n<0$, and $H^0(X,F)=
\ker(I_F^0(X)\to I_F^1(X))=F(X)$ by left exactness of sections. Restriction of scalars along
$R\to S$: for an injective sheaf $I$ of $S$-modules, $\Hom_R(G,I)=\Hom_S(S\otimes_RG,I)$, and
$S\otimes_R-$ carries the free resolutions computing the cohomology of a simplicial presheaf
(\cref{lem:pf:steps2:fam}(a)) to free resolutions; hence $I$ is acyclic as a sheaf of $R$-modules, and
an injective resolution of a sheaf of $S$-modules computes its cohomology as a sheaf of $R$-modules.
The reduction maps $\Z/Q\to\Z/q$ on cohomology are induced by maps of resolutions lifting the
reduction of sheaves; their transitivity and their compatibility with cup products, units and first
Chern classes follow from the uniqueness up to homotopy of such lifts and from the definitions of
these classes by explicit cocycles. Cup products are defined with the Alexander--Whitney diagonal of
the simplicial chains; associativity follows from its coassociativity, and graded commutativity from a
chain homotopy between the diagonal and its composite with the swap of factors (with the Koszul sign),
which the development constructs explicitly; cf.\ \cite[Ch.~V, \S1]{Milne1980} for the classical
statements. The statement on filtered colimits of rings is the case of \cref{lem:pf:steps2:limits}
for the limit of the affine schemes $\Spec A_j$.

(k) This is the conjunction of (h), (c), (d), (e) and the projective bundle formula.
\end{proof}

\begin{proposition}[Basic properties of \'etale cohomology]\label{prop:pf:steps2:etbasic}
Let $R$ be a commutative ring and $F$ a sheaf of $R$-modules on the big \'etale site.
\begin{enumerate}[label=\textup{(\alph*)},leftmargin=2.2em]
\item A quasi-isomorphism of complexes induces bijections on the cohomology with coefficients in
  $F$ computed from them.
\item (Smooth descent.) For a smooth surjective morphism $f\colon Y\to X$, the augmented \v{C}ech nerve
  of $f$ induces bijections on cohomology with coefficients in $F$.
\item A morphism of simplicial presheaves inducing bijections in every simplicial degree induces
  bijections.
\item (Mayer--Vietoris.) For compatible covers by two pieces, a morphism inducing bijections on the
  pieces and on their intersection induces bijections on the total spaces.
\item Every cohomology class of positive degree on a scheme vanishes on some \'etale surjective
  cover.
\item A short exact sequence of sheaves gives the long exact cohomology sequence.
\item Restriction of scalars along a ring map $R\to S$ does not change cohomology.
\item A morphism of sheaves which is bijective on sections over every scheme $U$ with
  $E_0(U)\ne\emptyset$ induces bijections on the cohomology of the simplicial presheaf $E$.
\item (Continuity.) For a filtered colimit of rings $A=\varinjlim A_j$ and $q$ invertible, every class
  in $H^n(\Spec A,\Z/q(i))$ comes from some $\Spec A_j$, and a class on $\Spec A_j$ vanishing on
  $\Spec A$ vanishes on some $\Spec A_k$, $k\ge j$.
\item If $A$ is a strictly henselian local ring, then $H^n(\Spec A,F)=0$ for $n>0$.
\item (Stalkwise criterion.) Let $f\colon Y\to X$ and $p\colon X\to S$ be quasi-compact and
  quasi-separated, with $q$ invertible on $S$. If for every strictly henselian local ring $A$ and
  every $\Spec A\to S$ the base change of $f$ to $\Spec A$ induces bijections, then $f$ induces
  bijections.
\item (Integral descent.) For an integral surjective morphism $f\colon Y\to X$ with $q$ invertible on
  $X$, the augmented \v{C}ech nerve of $f$ induces bijections on $H^\ast(-,\Z/q(i))$.
\item If $A$ is an absolutely integrally closed domain with $q\in A^\times$, then
  $H^n(\Spec A,\Z/q(i))=0$ for $n>0$.
\item $H^1(\Spec A,\mathbb G_m)\cong\Pic(A)$ for every commutative ring $A$.
\item If $q\in A^\times$ and $A$ contains a root of the $q$-th cyclotomic polynomial, then
  $H^n(\Spec A,\Z/q(i))\cong H^n(\Spec A,\Z/q(j))$ for all $n,i,j$.
\end{enumerate}
\formallabel{GlobalTheory.T.quasiIso_hypercohomology_bijective}\formallabel{GlobalTheory.T.etC_a2}\formallabel{GlobalTheory.T.levelwiseEquivalence_isEquivalence}\formallabel{GlobalTheory.T.mayerVietoris_equivalence_of_cover}\formallabel{GlobalTheory.T.positiveDegreeClasses_die_on_etaleCover}\formallabel{GlobalTheory.T.etaleCohomology_coefficient_longExactSequence}\formallabel{GlobalTheory.T.etaleCohomology_coefficientRing_invariance}\formallabel{GlobalTheory.T.etaleCohomology_bijective_of_support_bijective}\formallabel{GlobalTheory.T.etaleCohomology_twisted_filtered_colimit_continuity}\formallabel{GlobalTheory.T.strictlyHenselianLocalRing_acyclic_every_sheaf}\formallabel{GlobalTheory.T.etC_b3}\formallabel{GlobalTheory.T.integralSurjection_cechDescent_tateTwist}\formallabel{GlobalTheory.T.absolutelyIntegrallyClosed_domain_tateTwistAcyclic}\formallabel{GlobalTheory.T.etale_descent_picard_group}\formallabel{GlobalTheory.T.tateTwistCohomology_spec_twist_independent}
\end{proposition}

\begin{proof}
We use \cref{lem:pf:steps2:fam}(a): $H^n(E,F)$ is the cohomology of the double complex
$\Hom(\Z[E_\bullet],I_F^\bullet)$ built from the free presheaves on the terms of $E$, and for a scheme it
is $H^n(I_F(X))$.

(a) Maps into a bounded below complex of injective sheaves carry quasi-isomorphisms of bounded below
complexes to quasi-isomorphisms (standard homological algebra).
(b) Smooth surjective morphisms are morphisms of cohomological descent
\cite[Exp.~V\textsuperscript{bis}, 4.3.3(c)]{ArtinGrothendieckVerdier1972}; for the big \'etale site and
arbitrary sheaves one uses that a smooth surjective morphism has sections \'etale locally
\cite[17.16.3]{Grothendieck1967}.
(c) Filter the double complex by the simplicial degree: the map induces isomorphisms on the $E_1$ terms
of the first quadrant spectral sequences, hence on cohomology.
(d) The two covers give short exact sequences of the free complexes, hence Mayer--Vietoris sequences,
and the five lemma applies.
(e) A class of degree $n>0$ on a scheme $X$ is represented by a cocycle $a\in I_F^n(X)$; as $I_F$ is
exact in positive degrees as a complex of sheaves, $a$ is a coboundary on the members of an \'etale
covering of $X$.
(f) A short exact sequence of sheaves has injective resolutions forming a short exact sequence
(horseshoe lemma); apply $\Hom(\Z[E_\bullet],-)$, which is exact on injectives.
(g) This was proved in (j) of \cref{prop:pf:steps2:etfound}.
(h) Let $\mathcal C$ be the full subcategory of schemes $U$ with $E_0(U)\ne\emptyset$; it is a sieve
(closed under schemes mapping to it), and all terms $E_k$ are supported on it (via the face maps). The
restriction of sheaves to $\mathcal C$ is exact and has an exact left adjoint (extension by zero), so it preserves injectives, and the cohomology of $E$ only depends on the restriction of the
coefficients to $\mathcal C$, on which $\varphi$ is an isomorphism.
(i) This is \cref{lem:pf:steps2:limits} for the limit of the affine schemes $\Spec A_j$ (the values of
$\Z/q(i)$ on affine schemes commute with filtered colimits of rings).
(j) By (e) a class of positive degree vanishes on some \'etale surjective $W\to\Spec A$, which has a
section since $A$ is strictly henselian \cite[Ch.~I, \S4]{Milne1980}; so the class vanishes. (This
is \cite[Exp.~VIII, 4.7]{ArtinGrothendieckVerdier1972} for the small site.)
(k) By \cite[Exp.~VIII, 5.2]{ArtinGrothendieckVerdier1972} the hypothesis says that the map
$Rp_\ast\Lambda\to R(p\circ f)_\ast\Lambda$ is an isomorphism on all stalks; hence it is an isomorphism,
and the Leray spectral sequences give the claim.
(l) Integral surjective morphisms are morphisms of cohomological descent
\cite[Exp.~V\textsuperscript{bis}, 4.3.3(a)]{ArtinGrothendieckVerdier1972}.
(m) Every localization of $A$ at a prime is an absolutely integrally closed local domain, hence
strictly henselian, so by \cite[Exp.~VIII, 5.2, 4.7]{ArtinGrothendieckVerdier1972} the higher direct
images of $\Z/q(i)$ under the morphism from the \'etale to the Zariski site of $\Spec A$ vanish, and
$H^n(\Spec A,\Z/q(i))$ is the Zariski cohomology of the constant sheaf $\Z/q$ (the twist is trivial,
$A$ containing the $q$-th roots of unity). Constant sheaves on the irreducible space $\Spec A$ are
flasque, so this vanishes for $n>0$.
(n) This is Hilbert's Theorem 90 in the form $H^1(X,\mathbb G_m)=\Pic(X)$ \cite[Ch.~III, 4.9]{Milne1980}.
(o) A root $\zeta$ of the $q$-th cyclotomic polynomial in $A$ defines an isomorphism
$\Z/q(i)\cong\Z/q(j)$ of sheaves over $\Spec A$ (multiplication by the section defined by $\zeta$ in
$\Z/q(j-i)$).
\end{proof}
\begin{remark}\label{rem:pf:steps2:e3name}
The formal theorem proving (b) carries the word ``duality'' in its name, but its statement is the
base change statement (b). Statement (b) is the special case, for absolutely integrally closed local
bases, of the local acyclicity of smooth morphisms; it is exactly what the specialization argument
of \cref{prop:pf:steps2:specgood} needs. Similarly, one of the formal theorems proving
\cref{prop:pf:steps2:etfound}(e) carries the words ``locally acyclic'' in its name, but its statement is
(e).
\end{remark}

\begin{proposition}[Base change of towers]\label{prop:pf:steps2:towerfound}
\begin{enumerate}[label=\textup{(\alph*)},leftmargin=2.2em]
\item If $k\to\Omega$ is an embedding of algebraically closed fields, $\ell\in k^\times$ and $X$ is
  proper and smooth over $k$, then pull-back $H^n(X,\Ql(i))\to H^n(X_\Omega,\Ql(i))$ is bijective
  \cite[Exp.~XVI]{ArtinGrothendieckVerdier1972}.
\item If $\tau\colon k\to K$ is an embedding of fields of characteristic $0$, $\Omega$ an algebraic
  closure of $k$, $\Omega'$ an algebraically closed extension of $K$ and $j\colon\Omega\to\Omega'$
  extends $\tau$, then for every nice tower $H$ over $k$ and every matrix $\rho_0$ of a $\theta$-piece
  of $H$ there is a matrix $\rho_1$ of the $\theta$-piece of $\mathrm{bc}_\tau H$ with
  $\rho_1(\sigma')=\rho_0(\sigma)$ whenever $\sigma'\circ j=j\circ\sigma$.
\item For a place $A$ over $w$ and a bridge $B_w$ inducing $A$, there is a ring map
  $\stepsiiOw\to\OO_{K_w}$ compatible with $E\to K_w$ which maps uniformizers to uniformizers.
\item Covers base-change: for $\tau\colon k\to K$ and $\varphi\colon O\to O'$ compatible with $\tau$,
  every cover $C$ of $(H,\theta)$ over $O$ yields a cover $C'$ of $(\mathrm{bc}_\tau H,\theta)$ over
  $O'$ whose model is the base change of that of $C$, which is onto (good) if $C$ is, and for which
  the base changes of correspondences lying under $g$ with multiplicity $m$ lie under $g$ with
  multiplicity $m$.
\item For a finite extension $K/\Ql$ in $\Qlbar$, a nice tower $H$ over $K$ and an onto cover $C$ over
  $\OO_K$, pull-back along $Z_{\Qlbar}\to(Y_K)_{\Qlbar}$ is injective on $L\otimes\stepsiiH$, its image
  contains $\stepsiiPiece_\theta(H,\Qlbar,L)$, and it commutes with the actions of $\Gal(\Qlbar/K)$.
\item Base change along a field embedding preserves niceness of Hecke towers.
\item For a nice tower $H$ over $k$ and an algebraically closed extension $\Omega$ of characteristic
  $0$, every Hecke operator $T_g$ is an operator of the correspondence $(s_{g,\Omega},t_{g,\Omega})$
  (this follows from \cref{lem:pf:steps2:hopop}, as $\ell$ is invertible in $\Omega$).
\item $H^n(X,\Ql)=0$ for every scheme $X$ and $n<0$.
\end{enumerate}
\formallabel{UpperTheory.Gal.G3.vBcS_holds}\formallabel{UpperTheory.Gal.G3.locTower_of}\formallabel{UpperTheory.Gal.Seam.locTower_holds}\formallabel{UpperTheory.Gal.Seam.intMap}\formallabel{UpperTheory.Gal.Seam.bcCover}\formallabel{UpperTheory.Gal.Seam.pieceIn_holds}\formallabel{UpperTheory.Gal.Seam.niceBc}\formallabel{UpperTheory.Gal.Seam.hopEx_holds}\formallabel{UpperTheory.Gal.Seam.negVan}
\end{proposition}

\begin{proof}
(a) For the coefficients $\Z/\ell^k(i)$ this is the invariance of \'etale cohomology under extension of
separably closed base fields \cite[Exp.~XVI, 1.6]{ArtinGrothendieckVerdier1972} (the twist is trivial
over $k$); pass to the inverse limit and tensor with $\Ql$.
(f) Smoothness, projectivity and the property of being finite \'etale of a given degree are stable
under base change.
(g) By \cref{lem:pf:steps2:hopop}, $T_g=\bigoplus_n\stepsiitr_{s_{g,\Omega},1}\circ t_{g,\Omega}^\ast$,
which is an operator of $(s_{g,\Omega},t_{g,\Omega})$.
(h) Each $H^n(X,\Z/\ell^k)$ vanishes for $n<0$ (\cref{prop:pf:steps2:etfound}(j)), hence so do the
inverse limit and its tensor product with $\Ql$.
(b)--(e) are specific to the development; their formal proofs are not converted here
(\cref{rem:pf:steps2:deferred}).
\end{proof}

\begin{proposition}[$p$-adic Hodge theory]\label{prop:pf:steps2:padicfound}
Let $K/\Ql$ be a finite extension inside $\Qlbar$, with maximal unramified subextension $K_0$.
\begin{enumerate}[label=\textup{(\alph*)},leftmargin=2.2em]
\item $B_{\mathrm{dR}}^{\Gal(\Qlbar/K)}=K$, and every $\Gal(\Qlbar/K)$-invariant element of
  $B_{\mathrm{cris}}$ lies in $K_0$ \cite{Fontaine1994}.
\item Every tied period package admits a refined identification with the concrete period rings:
  isomorphisms of its rings with the development's $B_{\mathrm{cris}}$ and $B_{\mathrm{dR}}$ which are
  compatible with $\Zl$, with the map $\theta$ on $B^+_{\mathrm{dR}}$ and with the Galois actions.
\item The crystalline comparison theorem for smooth projective models over $\OO_K$ (crystallinity of
  $H^n(M_{\Qlbar},\Ql)$, de Rham comparison and Frobenius compatibility;
  \cite{FontaineMessing1987,Faltings1989,Tsuji1999,BhattMorrowScholze2018}) follows from its
  \emph{\'etale side}: finiteness, cup products and Poincar\'e duality for the models.
\item The following internal steps of the formal proof of the crystalline comparison theorem, whose
  role is described in \cref{sec:pf:steps1}, are proved here; we paraphrase their formal statements:
  the comparison maps are injective and strict; the classes of models used have decompositions into
  components; for all $d,j$ there is $N$ such that for every $d$-dimensional atlas the comparison maps
  $\beta$ from the Zariski--\'etale complexes are injective up to $\mu^N$ and every class is hit up to
  $\mu^N$ by a class of the filtered complex, where $\mu=\mu_\varepsilon$ is the element of the period ring attached to a compatible system $\varepsilon$ of $\ell$-power roots of unity; these
  maps $\beta$ are compatible with reduction, Galois action, Frobenius, pull-back and cup products;
  descent properties hold for the unramified, $A_{\mathrm{cris}}$ and filtered Witt-vector pieces; a
  good \'etale comparison family and a good level system of Hodge--Tate periods exist; the basic facts
  on the period rings hold; a Frobenius factorisation exists; and graded cohomology of affine pieces
  vanishes in the relevant range.
\item For a tied package with a refined identification as in (b), matrix representations of
  $\Gal(\Qlbar/K)$ that are admissible for the development's $B_{\mathrm{cris}}$ are admissible for
  the package; and, for a tied package, a representation $r\colon\Gal(\Qlbar/K)\to\GL_d(\Ql)$ admissible
  for the package becomes crystalline, in the sense of the package, after extension of scalars to any
  finite extension $L/\Ql$.
\item (The \'etale side for smooth models.) For smooth projective models $M$ over $\OO_K$: cup
  products exist; each $H^j(M_{\Qlbar},\Ql)$ is finite-dimensional; Poincar\'e duality holds for
  connected models of pure relative dimension; and every model decomposes into finitely many
  connected ones on which cohomology is detected.
\item A smooth projective $\OO_K$-scheme is a good model (projective and strictly semistable) in the
  sense of the comparison theorem.
\item Every projective morphism $X\to\Spec R$ carries a tautological line bundle, given by a cocycle of
  units on an \'etale surjective cover of $X$ obtained from an embedding in $\mathbb P^N_R$.
\end{enumerate}
\formallabel{GroundTheory.PAdicHodge.bDeRham_fixed}\formallabel{GroundTheory.PAdicHodge.ACris.bCris_fixed_mem_maxUnramified}\formallabel{UpperTheory.Lit.Cmp.locTied6}\formallabel{CrystallineSemistableComparison.cstSmooth_of_side}\formallabel{CrystallineSemistableComparison.inj}\formallabel{CrystallineSemistableComparison.comp}\formallabel{CrystallineSemistableComparison.B.main}\formallabel{CrystallineSemistableComparison.B.compat}\formallabel{CrystallineSemistableComparison.D.out}\formallabel{CrystallineSemistableComparison.E.etMap}\formallabel{CrystallineSemistableComparison.HT.levSys_ex}\formallabel{period_ring_base.base}\formallabel{CrystallineSemistableComparison.Pb.fac_ex}\formallabel{CrystallineSemistableComparison.Pb.gr_aff_pf}\formallabel{UpperTheory.Gal.Seam.isAdmP_of_isAdm}\formallabel{UpperTheory.Gal.Seam.isCrys_of_isAdmP}\formallabel{UpperTheory.Gal.KunnethLefschetz.Top.etSide_sm_holds}\formallabel{UpperTheory.Gal.Seam.HT.etSideS}\formallabel{UpperTheory.Gal.Seam.goodOf}\formallabel{CrystallineSemistableComparison.Pb.isO1_exists}\formallabel{CrystallineSemistableComparison.Pb.isO1_etale}
\end{proposition}

\begin{proof}
(a) These are Fontaine's computations of the Galois invariants of the period rings,
$(B_{\mathrm{dR}})^{\Gal(\Qlbar/K)}=K$ and $(B_{\mathrm{cris}})^{\Gal(\Qlbar/K)}=K_0$ \cite{Fontaine1994}; the
formal statements are the first equality and the inclusion of the invariants of $B_{\mathrm{cris}}$ in
$K_0$.
(e), first statement: the refined identification gives a ring isomorphism between the development's
$B_{\mathrm{cris}}$ and the corresponding ring of the package, compatible with the Galois actions; an
admissibility datum (matrices $A$, $B$ with entries in $B_{\mathrm{cris}}$, $AB=1$, and the
transformation rule under $\sigma$) is transported entrywise along this isomorphism.
(g) Locally on $M$, a smooth $\OO_K$-scheme is \'etale over an affine space over $\OO_K$
\cite[17.11.4]{Grothendieck1967}; such a chart is a strictly semistable chart with no branches, for any
uniformizer of $\OO_K$, and projectivity is part of the hypothesis.
(h), second statement: the \'etale cover is part of the definition of a tautological line datum.
(b), (c), (d), the second statement of (e), (f) and the first statement of (h) are specific to the
development; their formal proofs are not converted here (\cref{rem:pf:steps2:deferred}). For (c), the
formal proof combines \lref{CrystallineSemistableComparison.Pb.pSeam} (the comparison maps
constructed in \cref{sec:pf:steps1}) with residue data (\cref{lem:pf:steps2:resdata}) and the
statements of (d).
\end{proof}

\begin{remark}[Statements whose proofs are not converted]\label{rem:pf:steps2:deferred}
The following tagged statements of this subsection are specific to the development and their formal
proofs are not converted in this section (sizes are numbers of Lean lines of the dominated formal
proofs): \cref{prop:pf:steps2:towerfound}(b) (306), (c) (208), (d) (208), (e) (1,328);
\cref{prop:pf:steps2:padicfound}(b) (1,348), (c) (725), (d) (from 286 to 9,765 for the individual
steps; about 30,000 in total), (e) second statement (115), (f) (derived from the package
\cref{prop:pf:steps2:etfound}(k)) and (h) first statement (374). For the reduction compatibilities and
the cup product identities in \cref{prop:pf:steps2:etfound}(j) (1,338, 508 and 249 lines) only the
mechanism of the proof is indicated above.
\end{remark}
